\documentclass[11pt]{ctexart}
\usepackage[a4paper,margin=1in]{geometry}
\usepackage{graphicx}
\usepackage{xcolor}
\usepackage{hyperref}
\definecolor{refgray}{gray}{0.45}
\title{\textbf{市场最新观点汇总}\\[0.4em]{\large 中国经济正从“消费刺激依赖”转向“供给端投资驱动”，全球市场在“和平交易”与结构性供给约束之间重新定价。}}
\author{KC桌面 自动生成}
\date{260618}
\begin{document}
\maketitle
\tableofcontents
\newpage
\section{一页摘要}
\begin{itemize}
  \item 中国经济正经历从消费刺激到供给端投资的增长逻辑切换，5月零售销售转负、投资降幅扩大，但AI超周期拉动高科技生产创五年新高。
  \item 全球能源市场在美伊和平协议预期下大幅回调，但多家投行认为供给恢复存在物理约束，三季度供需缺口仍将支撑油价。
  \item 中国家庭资产正从房产向金融资产结构性迁移，房地产占比从67\%降至52\%，金融资产占比升至20\%。
  \item 欧洲数据中心并网申请飙升至480吉瓦，相当于1.5倍当前电力需求，正在重塑电力公用事业的估值逻辑。
  \item 亚洲外汇市场出现根本性分化：印尼和印度央行储备消耗接近红线，中国央行反向买入美元以放缓人民币升值节奏。
  \item 中国车企在欧洲的份额侵蚀从个别品牌扩散至全行业，欧洲车企FY27盈利仍有超过10\%的下行风险。
\end{itemize}
\section{宏观与利率：中国增长逻辑切换与全球通胀滞后效应}
\textbf{中国经济正从消费刺激转向供给端投资驱动，而全球能源冲击对消费的滞后传导将在下半年集中释放。}

\begin{itemize}
  \item MS估算二季度GDP跟踪值降至4.4\%，零售销售自2022年以来首次转负，固定资产投资单月同比跌幅扩大至10.7\%，政策重心正加速转向AI算力网络、数据中心、智能电网等“六网”基础设施投资。
  \item GS指出市场低估了能源冲击向消费者现金流的滞后传导：批发价格向零售价格的传导延迟、冬季取暖需求的季节性放大、以及美国退税红利在下半年消退，将使消费实质性放缓在下半年才显现。
  \item Citi报告显示中国5月零售销售同比-0.6\%，固定资产投资累计同比-4.1\%，CPI温和+1.2\%而PPI持续回升，国内“滞胀”风险上升，AI出口与内需疲软形成K型分化。
  \item GS对澳大利亚5月CPI的预测显示，汽油价格单月暴跌12.5\%是临时政策效果，剔除一次性因素后核心截尾均值CPI同比从3.4\%升至3.5\%，内生价格压力在加速。
  \item MS认为中国材料工业的分化在于供给侧刚性约束：粗钢、水泥、玻璃产量普遍下滑，但电解铝产量同比+1.7\%且环比上升，运行产能已接近天花板。
  \item GS在陆家嘴论坛后指出，信贷增速放缓是融资结构不可逆的转型，贷款在社融增量中的占比已从80\%降至45\%，投向“五篇大文章”的新增贷款占比跃升至70\%以上。
\end{itemize}
\begin{figure}[htbp]
\centering
\includegraphics[width=0.88\linewidth]{figures/fig_001.jpg}
\caption{Exhibit 1 - \#\# Zhipeng Cai Economist Asia Summer School 2026}
\end{figure}
\begin{figure}[htbp]
\centering
\includegraphics[width=0.88\linewidth]{figures/fig_002.jpg}
\caption{Exhibit 1 - Exhibit 1: Consumer Spending Has Held Up So Far Despite the Surge in Energy Prices}
\end{figure}
\begin{figure}[htbp]
\centering
\includegraphics[width=0.88\linewidth]{figures/fig_036.jpg}
\caption{Figure 1 - ■ Stagflation risks on the domestic side rise further. Domestic demand indicators missed already subdued expectations. Retail sales contracted for the first time since Covid at -0.6\%YoY (Citi/Mkt: 0.0/-0.2\%YoY). Investment decline deepened to -4.1\%YoY Ytd (Cit}
\end{figure}
\begin{figure}[htbp]
\centering
\includegraphics[width=0.88\linewidth]{figures/fig_018.jpg}
\caption{Exhibit 1 - Exhibit 1: We expect headline CPI and trimmed mean to increase \$4.2\textbackslash{}\%\$ yoy and \$3.5\textbackslash{}\%\$ yoy respectively in May}
\end{figure}
\begin{figure}[htbp]
\centering
\includegraphics[width=0.88\linewidth]{figures/fig_284.jpg}
\caption{Exhibit 1 - Exhibit 1: Indirect financing, predominantly loans, has seen its share in Total Social Financing (TSF) increments decline from over 80\% to 45\% by 2025}
\end{figure}
{\small\color{refgray}\textbf{References:} [R001] 摩根斯坦利：市场真正低估的不是需求疲软，而是政策转向“供给端投资”的结构性含义; [R002] GS：消费的真正压力尚未到来，市场对能源冲击的滞后效应严重低估; [R004] GS：油价传导的真正压力不在五月，而在接下来的几个月; [R005] Citi：市场真正低估的不是AI需求，而是国内“滞胀”的定价权转移; [R023] GS：银行股真正的防御性，来自资产负债表结构而非规模; [R027] 摩根斯坦利：中国材料工业的真正分化不在需求，而在供给侧的刚性约束}

\section{AI与算力：基础设施定价权重估与硬件ROE结构性上升}
\textbf{AI基础设施的供给侧物理瓶颈正在将定价权推向新高，而亚洲硬件公司正经历由AI驱动的ROE结构性上升周期。}

\begin{itemize}
  \item GS调研显示数据中心每千瓦租赁价格已接近200美元，预计年内向250美元迈进，而2021年仅约70美元，供给约束来自电力、设备和选址三个层面且短期无解。
  \item 欧洲数据中心并网申请总量飙升至约480吉瓦，相当于1.5倍当前电力需求，GS筛选的三大受益集群（变革性电气化、可再生能源、能源安全）2030年平均市盈率仅12倍。
  \item Citi引入2028年WFE牛市预测2500亿美元，核心驱动从训练侧转向推理侧，DRAM供应瓶颈正迫使产业链重新设计内存架构，为NAND设备打开新需求空间。
  \item Bernstein分析显示亚洲硬件公司ROE上升由AI需求驱动而非周期反弹，Chroma ATE的ROE从2024年25\%跃升至2026年预计45\%以上，AI芯片测试和光收发器测试业务爆发。
  \item Bernstein判断NVIDIA的SOCAMM2新型内存封装大概率局限于自身生态，不会对更广泛的服务器内存接口芯片市场构成实质性威胁，市场高估了替代风险。
  \item GS实地调研显示企业级AI基础设施规模化采用仍处早期，但Lumen测算“东-西向流量”市场有580亿美元且年增速13\%，物理光纤网络是核心护城河。
\end{itemize}
\begin{figure}[htbp]
\centering
\includegraphics[width=0.88\linewidth]{figures/fig_216.jpg}
\caption{Exhibit 1 - Exhibit 1: The EU DC pipeline has more than doubled in one year EU 28 GSe datacenter pipeline evolution, 2025-26 (GW)}
\end{figure}
\begin{figure}[htbp]
\centering
\includegraphics[width=0.88\linewidth]{figures/fig_220.jpg}
\caption{Exhibit 4 - Exhibit 4: The three clusters trade on an average P/E of c.12x by 2030E Key valuation metrics, 2026-30E (x)}
\end{figure}
\begin{figure}[htbp]
\centering
\includegraphics[width=0.88\linewidth]{figures/fig_273.jpg}
\caption{Figure 2 - KV Cache offloading - KV cache offloading refers to the architectural technique of dynamically transferring intermediate attention states (key-value tensors) from HBM into lower-cost, higher-capacity tiers such as DRAM and, increasingly, NAND flash, effectivel}
\end{figure}
\begin{figure}[htbp]
\centering
\includegraphics[width=0.88\linewidth]{figures/fig_230.jpg}
\caption{EXHIBIT 1 - EXHIBIT 1: The trend of ROE among companies reflects broader sector trends and cycles, and individual company competitiveness in each market}
\end{figure}
\begin{figure}[htbp]
\centering
\includegraphics[width=0.88\linewidth]{figures/fig_169.jpg}
\caption{EXHIBIT 9 - EXHIBIT 9: We project that NVIDIA CPU shipment volume will account for only MSD to low teens in total server CPU shipment \#\# SOCAMM2'S IMPACTS ON MEMORY INTERFACE CHIP SUPPLIERS \#\# A) SOCAMM2 EXPANDS TAM VS. SOLDERED LPDDR ON GRA}
\end{figure}
{\small\color{refgray}\textbf{References:} [R016] GS：欧洲数据中心正在重塑电力行业的估值逻辑，而市场仍按旧地图定价; [R020] Citi：半导体设备市场真正被低估的，是2028年2500亿美元WFE背后的结构性需求切换; [R018] Bernstein：亚洲硬件正在经历一轮由AI驱动的ROE结构性重定价; [R009] Bernstein：SOCAMM2不会侵蚀内存接口芯片市场，市场高估了替代风险; [R029] GS：AI基础设施定价权正在经历结构性重估，但企业级需求仍处于“等待爆发”的临界点}

\section{能源与大宗：和平交易过度定价，供给约束支撑三季度油价}
\textbf{市场对美伊和平协议的定价可能过头，供给恢复的物理约束和库存重建需求将为油价提供支撑。}

\begin{itemize}
  \item MS指出当前股价隐含的WTI油价预期约66美元/桶，显著低于12个月远期曲线约75美元，三季度全球原油仍将维持约340万桶/日的平均缺口。
  \item GS测算霍尔木兹流量恢复至战前70\%即可实现波斯湾出口正常化，通过现有绕行路线已形成每日910万桶替代流量，但船东仍在等待明确过境指引。
  \item GS发现中国库存数据存在“暗物质”缺口：5月隐含库存减少1.6mb/d，远超可见库存的0.3mb/d，暗示存在隐形去库存。
  \item JPM报告显示澳大利亚到中国的铁矿石散货运费骤降22\%至10.9美元/吨，澳大利亚矿商获得显著物流成本优势，而巴西和南非运费仍高于冲突前水平50\%以上。
  \item UBS指出中东局势缓和下化工板块的真正机会在成本端重新定价，涤纶长丝和磷化工受益于原材料降价，但需求端疲软意味着成本下降不等于利润提升。
  \item MS认为即便和平协议签署，从协议到产量回流存在数月真空期，OECD原油库存已跌至5年平均水平以下且仍在加速去库。
\end{itemize}
\begin{figure}[htbp]
\centering
\includegraphics[width=0.88\linewidth]{figures/fig_301.jpg}
\caption{Exhibit 1 - Exhibit 1: We estimate that oil E\&Ps currently discount an average WTI price of \textbackslash{}\textasciitilde{}\textbackslash{}\$65/bbl, well below the current 12-month strip of \textbackslash{}\textasciitilde{}\textbackslash{}\$75.}
\end{figure}
\begin{figure}[htbp]
\centering
\includegraphics[width=0.88\linewidth]{figures/fig_199.jpg}
\caption{Exhibit 3 - Exhibit 3: Normalization in Oil Exports From Gulf Producers to Their Pre-War Level May Be Achieved With a 12.7mb/d Increase in Hormuz Flows From Current Levels Estimating Mid June Hit to Oil Flows from Persian Gulf Countries}
\end{figure}
\begin{figure}[htbp]
\centering
\includegraphics[width=0.88\linewidth]{figures/fig_204.jpg}
\caption{Exhibit 8 - Exhibit 8: China Landed Visible Inventories Have Drawn by 0.5mb/d Over the Last 14 Days Global Visible Total Oil Inventories, Change Since Feb 27}
\end{figure}
\begin{figure}[htbp]
\centering
\includegraphics[width=0.88\linewidth]{figures/fig_156.jpg}
\caption{Figure 3 - Figure 3: Iron ore bulk shipping costs from Australia, Brazil and South Africa to China all recorded a decline in the past week with Australian bulk shipping rates -22\% WoW to \textbackslash{}\$10.9/t LHS: Iron ore price \textbackslash{}\$/mt; RHS: Shipping cost}
\end{figure}
\begin{figure}[htbp]
\centering
\includegraphics[width=0.88\linewidth]{figures/fig_179.jpg}
\caption{Figure 1 - Figure 1: China crude oil inventory}
\end{figure}
\begin{figure}[htbp]
\centering
\includegraphics[width=0.88\linewidth]{figures/fig_307.jpg}
\caption{Exhibit 8 - Exhibit 8: OECD oil stocks have fallen further below 5-year averages, driven by a large draw in the US... Total Crude Oil Inventories (mln bbls)}
\end{figure}
{\small\color{refgray}\textbf{References:} [R025] 摩根斯坦利：市场对“和平”的定价已经过头，但油价真正的支撑尚未被计入; [R014] GS：市场真正低估的不是需求崩塌，而是供给恢复的路径与节奏; [R008] JPM：铁矿石定价权正在向澳大利亚矿商倾斜，而非中国需求主导; [R013] UBS：中东局势缓和下，化工板块的真正机会不在油价，而在成本端的重新定价; [R006] UBS：航运市场的真正拐点不在运价，而在供给侧的不可逆重置}

\section{地产与消费：K型分化加剧，一线回暖难掩全国压力}
\textbf{中国楼市和消费市场均呈现K型分化，一线城市局部回暖并非全国复苏的领先指标，三季度可能迎来二次探底。}

\begin{itemize}
  \item MS数据显示5月全国销售金额同比降幅从-7.6\%扩大到-9.3\%，销售面积降幅从-9.5\%扩大到-13.2\%，三季度二手房销售预计转为同比负增长。
  \item JPM住房活动指数5月小幅回升，一线城市二手房价格环比+0.4\%，但全国新房销售面积同比跌幅从-6.5\%扩大至-8.6\%，低线城市量价齐跌。
  \item MS认为5月零售总额同比-0.6\%背后是三类分化：政策驱动型需求退潮、房地产关联需求低迷、日常消费韧性测试，补贴退坡后的真实需求压力正在显现。
  \item GS估算中国家庭资产中房地产占比从2021年高峰期的67\%降至52\%，现金及存款从16\%升至25\%，其他金融资产从15\%升至20\%，储蓄正从不动产向金融资产迁移。
  \item JPM分析一线城市回暖支撑因素并非居民购房意愿全面回升，而是二手房市场库存较低、成交量结构切换和政策资源集中。
  \item MS判断三季度销售二次探底的冲击力被市场低估，二手房从放缓转向负增长的传导效应和开发商土地市场近乎停摆的谨慎姿态将重塑下半年定价逻辑。
\end{itemize}
\begin{figure}[htbp]
\centering
\includegraphics[width=0.88\linewidth]{figures/fig_006.jpg}
\caption{Exhibit 2 - Exhibit 1: From 1997 to 2022, the household sector's net worth as a percentage of GDP increased twofold}
\end{figure}
\begin{figure}[htbp]
\centering
\includegraphics[width=0.88\linewidth]{figures/fig_010.jpg}
\caption{Exhibit 5 - Exhibit 5: Since the September 2024 policy pivot, equities have contributed more to financial asset growth Chinese households financial assets (GS estimate)}
\end{figure}
\begin{figure}[htbp]
\centering
\includegraphics[width=0.88\linewidth]{figures/fig_290.jpg}
\caption{Exhibit 2 - Exhibit 3: MSCI China vs. Retail Sales}
\end{figure}
\begin{figure}[htbp]
\centering
\includegraphics[width=0.88\linewidth]{figures/fig_291.jpg}
\caption{Exhibit 4 - Exhibit 4: China Retail Sales by Category}
\end{figure}
{\small\color{refgray}\textbf{References:} [R017] 摩根斯坦利：市场低估了三季度销售二次探底的冲击力; [R015] JPM：中国楼市真正的分水岭不在价格，而在供给侧的“窄复苏”能否撑到政策显效; [R024] 摩根斯坦利：消费复苏的真正考验不在总量，而在结构分化的再定价; [R003] GS：中国家庭资产正在经历从房产到金融资产的结构性迁移; [R001] 摩根斯坦利：市场真正低估的不是需求疲软，而是政策转向“供给端投资”的结构性含义}

\section{区域市场：亚洲外汇储备分化与东南亚互联网韧性}
\textbf{亚洲新兴经济体外汇储备消耗出现根本性分化，印尼和印度接近红线，而中国央行反向操作；东南亚互联网平台增长放缓但韧性犹存。}

\begin{itemize}
  \item NOM数据显示印尼央行5月卖出约48亿美元（占储备3.9\%），印度央行卖出115亿美元（占储备2.1\%），印尼储备充足率扣除短期负债后仅46\%。
  \item 中国央行5月净买入约287亿美元外汇储备，加上国有银行存入109亿美元，合计增持约396亿美元，主动放缓人民币升值节奏。
  \item DB报告显示出口商结汇比率连续三个月下降至62.6\%，但资本账户连续第二个月净流入，5月增量180亿美元为2021年以来最大，外资重新买入中国债券。
  \item NOM的人民币中间价定价模型显示预测误差正在系统性收窄，定价机制从“被动对冲”向“主动校准”转变，下半年关键政治议程可能影响汇率政策节奏。
  \item UBS数据显示Grab外卖订单量同比仍增长21\%，Shopee GMV同比+20.6\%，虽从3-4月高点回落，但日均订单环比仍增长，增长节奏从“爆发”切换为“韧性测试”。
  \item UBS指出东南亚外卖和电商竞争烈度环比回升，折扣力度加大，但Grab在菲律宾和马来西亚市场份额继续提升，印尼是唯一订单同比下滑的市场。
\end{itemize}
\begin{figure}[htbp]
\centering
\includegraphics[width=0.88\linewidth]{figures/fig_171.jpg}
\caption{Figure 1 - Figure 1: Exporters' FX conversion slowing for the second consecutive month}
\end{figure}
\begin{figure}[htbp]
\centering
\includegraphics[width=0.88\linewidth]{figures/fig_172.jpg}
\caption{Figure 2 - Figure 2: First net purchase of Chinese bonds by foreign investors since April 2025}
\end{figure}
\begin{figure}[htbp]
\centering
\includegraphics[width=0.88\linewidth]{figures/fig_340.jpg}
\caption{Figure 1 - Figure 1: Grab ASEAN avg daily orders MoM \%}
\end{figure}
\begin{figure}[htbp]
\centering
\includegraphics[width=0.88\linewidth]{figures/fig_347.jpg}
\caption{Figure 8 - Figure 8: Shopee's ASEAN GMV YoY \%}
\end{figure}
{\small\color{refgray}\textbf{References:} [R011] NOM：亚洲央行正在用外汇储备为货币信心买单，但印尼和印度的弹药已接近红线; [R010] DB：市场低估了资本账户对人民币的支撑力; [R022] NOM：人民币中间价定价模型释放的信号比汇率本身更重要; [R030] UBS：东南亚外卖与电商的增长故事并未结束，但竞争正在重新定价龙头门槛}

\section{汽车行业：补贴退潮暴露结构性压力，中国车企加速全球渗透}
\textbf{中国车市补贴退潮后暴露的是被透支的购买力而非需求基本面恶化，中国车企在欧洲的份额侵蚀正在从点状突破转向面状渗透。}

\begin{itemize}
  \item Bernstein估算2024-2025年补贴政策提前拉动了超过500万辆需求，5月零售同比-20\%本质上是“借来的时间”到期后的正常化过程。
  \item Citi周度订单追踪显示6月第二周新能源车订单环比下降13\%，半月累计环比仅增长4\%，远低于历史同期10\%的月环比增速，但比亚迪和吉利银河订单稳定性突出。
  \item MS报告显示欧洲OEM在EU5的市场份额将从2025年的67\%进一步下滑至2027年的62\%，此前份额流失集中在Stellantis一家，接下来将更具“普惠性”波及所有欧洲车企。
  \item Citi专家电话会披露BBA上周在华订单已降至不到7000台，而去年同期为11000-12000台，蔚来ES9稳态周订单维持在2000台，豪华车市场正在发生结构性权力交接。
  \item Bernstein指出中国车企进攻路径清晰：主攻SUV、电动车和插混车型，瞄准入门级市场，雷诺首当其冲，奔驰和宝马虽暂时受益于高端定位但面临向上突破压力。
  \item Citi判断如果中国政府满足于出口增长对冲内需疲弱，车企将被迫放弃激进策略，行业整合更温和、降价节奏放缓、成本管控更严格。
\end{itemize}
\begin{figure}[htbp]
\centering
\includegraphics[width=0.88\linewidth]{figures/fig_423.jpg}
\caption{Exhibit 2 - Exhibit 3 - Exhibit 4) EXHIBIT 2: Chinese autos retail sales came in at 1.51mn units in May 2026, -20.0\% yoy}
\end{figure}
\begin{figure}[htbp]
\centering
\includegraphics[width=0.88\linewidth]{figures/fig_064.jpg}
\caption{Exhibit 1 - Exhibit 1: European OEMs have continued to see their market share in Europe decline – consensus, optimistically, assumes this ceases}
\end{figure}
\begin{figure}[htbp]
\centering
\includegraphics[width=0.88\linewidth]{figures/fig_123.jpg}
\caption{Exhibit 84 - Exhibit 84: Chinese OEMs are significantly exposed to the C-SUV, D-SUV and B-SUV sub-segments Exhibit 85: Chinese OEM sub-segment exposure in Europe – sustained focus on D-SUV and C-SUV}
\end{figure}
\begin{figure}[htbp]
\centering
\includegraphics[width=0.88\linewidth]{figures/fig_128.jpg}
\caption{Exhibit 102 - Exhibit 102: Chinese OEMs keep gaining share in China (12M rolling monthly market share)}
\end{figure}
{\small\color{refgray}\textbf{References:} [R031] Bernstein：中国车市真正的考验不是需求萎缩，而是补贴退潮后供给侧的再定价; [R019] Citi：市场低估的不是需求疲软，而是供给端策略的结构性拐点; [R007] 摩根斯坦利：欧洲车企最坏的时刻尚未到来，中国OEM的份额侵蚀才刚刚加速; [R028] Citi：市场低估的不是需求疲软，而是供给端正在发生的结构性出清}

\section{生物医药与半导体：全球价值转化面临重估}
\textbf{中国生物医药的全球价值转化路径正面临制度性重估，而半导体设备需求的结构性切换为2028年WFE增长提供新动力。}

\begin{itemize}
  \item Bernstein数据显示中国占全球早期药物开发30\%-50\%，但美国三期临床试验占比低于10\%，FDA批准不到5\%，美国商业销售额仅1\%，“前重后轻”的结构性矛盾是真正瓶颈。
  \item Bernstein情景分析显示，若间接进入（通过欧盟、日本合作），BD收入跌50\%、全球峰值销售跌20\%；若完全封锁，BD收入跌80\%、全球销售跌60\%。
  \item Citi指出2028年WFE牛市预测2500亿美元的核心驱动正从训练侧转向推理侧，DRAM供应瓶颈迫使产业链重新设计内存架构，KV缓存卸载到NAND闪存成为新趋势。
  \item Citi对三家核心设备商给出目标价：AMAT 710美元（31x PE）、LRCX 450美元（40x PE）、KLAC 290美元（40x PE），当前估值低于历史峰值15-20\%但高于均值55-67\%。
  \item Bernstein分析显示日本龙头虽掌握PCB原材料优势但ROIC并不突出，而服务器相关PCB公司ROIC更高，说明AI才是真引擎。
  \item Citi报告强调Agentic AI推动内存需求结构性增长，多步推理工作流大幅扩大KV缓存，苹果AFM 3把大模型存在NAND而非DRAM，为NAND设备打开新需求空间。
\end{itemize}
\begin{figure}[htbp]
\centering
\includegraphics[width=0.88\linewidth]{figures/fig_173.jpg}
\caption{EXHIBIT 3 - EXHIBIT 3: Framework of the global drug discovery and development funnel EXHIBIT 4: As shown by first clinical trial registrations, China-origin assets now account for 55\% of global total, up from 38\% in 2021 Global first clinica}
\end{figure}
\begin{figure}[htbp]
\centering
\includegraphics[width=0.88\linewidth]{figures/fig_175.jpg}
\caption{EXHIBIT 6 - EXHIBIT 6: Out-licensing momentum strong - total deal value over \textbackslash{}\$90 Bn in 1H26, 70\% of 2025 total; Upfront payment \textbackslash{}\$4.8 Bn, 71\% of 2025 total China out-licensing activity: deal count, total deal value, and total upfront paymen}
\end{figure}
\begin{figure}[htbp]
\centering
\includegraphics[width=0.88\linewidth]{figures/fig_273.jpg}
\caption{Figure 2 - KV Cache offloading - KV cache offloading refers to the architectural technique of dynamically transferring intermediate attention states (key-value tensors) from HBM into lower-cost, higher-capacity tiers such as DRAM and, increasingly, NAND flash, effectivel}
\end{figure}
\begin{figure}[htbp]
\centering
\includegraphics[width=0.88\linewidth]{figures/fig_274.jpg}
\caption{Figure 3 - AMAT – We model total revenue growth of 30\%/22\% Y/Y in CY27/28, including 35\%/25\% from Silicon, and 14\%/13\% from AGS. We lift our TP to \textbackslash{}\$710 from \textbackslash{}\$550 prior, based on 31x P/E applied to CY28 EPS. The 31x P/E is consistent with our prior valuation and 55\% abo}
\end{figure}
\begin{figure}[htbp]
\centering
\includegraphics[width=0.88\linewidth]{figures/fig_238.jpg}
\caption{EXHIBIT 10 - EXHIBIT 10: PCB supply chain companies with higher exposure to server market would present better ROIC. ABF companies had relatively low ROIC last year as low-end product was still oversupplied 2025 ROIC (Y-axis) and Invested capita}
\end{figure}
{\small\color{refgray}\textbf{References:} [R012] Bernstein：中国生物医药被低估的不是创新能力，而是“全球价值转化”面临的真实阻力; [R020] Citi：半导体设备市场真正被低估的，是2028年2500亿美元WFE背后的结构性需求切换; [R018] Bernstein：亚洲硬件正在经历一轮由AI驱动的ROE结构性重定价}

\section{航运与供应链：供给约束驱动定价逻辑重置}
\textbf{全球航运市场定价逻辑正从需求驱动转向供给约束驱动，霍尔木兹海峡的船队效率损失可能已不可逆。}

\begin{itemize}
  \item UBS报告显示霍尔木兹海峡原油油轮通行量仍比冲突前水平低95\%，即便海峡完全重启，船队效率的损失可能已经不可逆，船东已调整航线规划和保险安排。
  \item UBS数据显示SCFI综合指数上周环比+9\%、同比+43\%，上海-欧洲航线环比+18\%，跨太平洋航线环比+10-12\%，亚洲到美国集装箱运量同比增速强劲。
  \item JPM报告显示澳大利亚到中国铁矿石散货运费骤降22\%至10.9美元/吨，巴西和南非运费仍高于冲突前50\%以上，澳大利亚矿商获得显著物流成本优势。
  \item UBS指出干散货BDI指数环比-10\%但年初至今仍涨61\%，好望角型散货船平均收益环比-18\%，造船新船价格指数环比持平保持高位。
  \item GS测算显示全球柴油和汽油高频可见库存已降至季节性区间底部，中国和欧洲道路燃料零售销售额同比分别下降23.3\%和3.5\%。
  \item UBS报告强调集装箱运量稳增、运费继续上行，中国主要港口集装箱吞吐量同比+5\%，但霍尔木兹恢复节奏仍是油轮运价的关键变量。
\end{itemize}
\begin{figure}[htbp]
\centering
\includegraphics[width=0.88\linewidth]{figures/fig_056.jpg}
\caption{Figure 4 - Figure 4: Crude tanker passage at Hormuz remained 95\% below pre-conflict level}
\end{figure}
\begin{figure}[htbp]
\centering
\includegraphics[width=0.88\linewidth]{figures/fig_057.jpg}
\caption{Figure 5 - Figure 5: Container throughput at China's key ports -4\% WoW, 5\% YoY last week}
\end{figure}
\begin{figure}[htbp]
\centering
\includegraphics[width=0.88\linewidth]{figures/fig_059.jpg}
\caption{Figure 7 - Figure 7: SCFI +9\% WoW and +43\% YoY Figure 9: Container ships re-routing away from Red Sea still at high levels (+4-5\% YoY)}
\end{figure}
\begin{figure}[htbp]
\centering
\includegraphics[width=0.88\linewidth]{figures/fig_156.jpg}
\caption{Figure 3 - Figure 3: Iron ore bulk shipping costs from Australia, Brazil and South Africa to China all recorded a decline in the past week with Australian bulk shipping rates -22\% WoW to \textbackslash{}\$10.9/t LHS: Iron ore price \textbackslash{}\$/mt; RHS: Shipping cost}
\end{figure}
\begin{figure}[htbp]
\centering
\includegraphics[width=0.88\linewidth]{figures/fig_205.jpg}
\caption{Exhibit 10 - Exhibit 10: Global Diesel and Gasoline High-Frequency Visible Stocks Decreased to the Bottom of Their Seasonal Range Global Visible High-Frequency Gasoline and Diesel Stocks}
\end{figure}
{\small\color{refgray}\textbf{References:} [R006] UBS：航运市场的真正拐点不在运价，而在供给侧的不可逆重置; [R008] JPM：铁矿石定价权正在向澳大利亚矿商倾斜，而非中国需求主导; [R014] GS：市场真正低估的不是需求崩塌，而是供给恢复的路径与节奏}

\section{风险偏好与政策：银行股结构溢价与香港楼市政策误读}
\textbf{银行股定价正从规模溢价转向结构溢价，而市场对内地资金买港楼的恐慌可能高估了政策冲击。}

\begin{itemize}
  \item GS在陆家嘴论坛后指出，银行资产负债表的“质量”和“结构”正在取代“增速”和“规模”成为决定估值分化的核心变量，利率走廊收窄利好中小银行。
  \item GS数据显示贷款在社融增量中的占比从超过80\%降至45\%，投向“五篇大文章”的新增贷款占比跃升至70\%以上，“更慢、更高质量的贷款增长”成为新常态。
  \item JPM专家电话会明确837号文并非针对内地居民购买香港物业，现行监管框架没有发生根本性变化，内地居民用境内资金购买境外房产“从未被允许过”。
  \item JPM分析显示内地买家占香港住宅交易约5-10\%（按金额计15-20\%），违规用境内资金买房的处罚通常为违规金额的5-10\%，实践中很少要求强制处置房产。
  \item GS指出非银流动性工具的研究旨在避免2022年底理财赎回潮重演，对银行债券投资收益是利好，尤其交易账户占比高的银行。
  \item JPM报告强调真正值得关注的是执行层面可能带来的“寒蝉效应”，尤其是对“内地税务居民”的全球收入征税风险以及CRS信息交换的潜在影响。
\end{itemize}
\begin{figure}[htbp]
\centering
\includegraphics[width=0.88\linewidth]{figures/fig_284.jpg}
\caption{Exhibit 1 - Exhibit 1: Indirect financing, predominantly loans, has seen its share in Total Social Financing (TSF) increments decline from over 80\% to 45\% by 2025}
\end{figure}
\begin{figure}[htbp]
\centering
\includegraphics[width=0.88\linewidth]{figures/fig_286.jpg}
\caption{Exhibit 3 - Exhibit 3: Among the banks we cover, Industrial Bank, Huaxia Bank, and BONB exhibit the highest interbank liability proportion, recorded at 29\%, 26\%, and 22\% respectively in 2025, in contrast to 15\% for the large four SOE banks As}
\end{figure}
\begin{figure}[htbp]
\centering
\includegraphics[width=0.88\linewidth]{figures/fig_337.jpg}
\caption{Figure 1 - Figure 1: HK private residential market – \% of “Mainland Chinese” buyers with a last name in Mandarin pinyin (by volume)}
\end{figure}
\begin{figure}[htbp]
\centering
\includegraphics[width=0.88\linewidth]{figures/fig_338.jpg}
\caption{Figure 2 - Figure 2: HK private residential market – \% of “Mainland Chinese” buyers with a last name in Mandarin pinyin (by value)}
\end{figure}
{\small\color{refgray}\textbf{References:} [R023] GS：银行股真正的防御性，来自资产负债表结构而非规模; [R026] JPM：市场对内地资金买港楼的恐慌，可能高估了政策冲击}

\section{消费行业：香水市场进入爆款竞赛模式}
\textbf{香水行业增速正常化后，增长动力从品类红利转向打造爆款的能力，品牌价值分化取决于新品转化率。}

\begin{itemize}
  \item GS数据显示香水行业增速从2023年的8\%回落至2025年的3.5\%，但新品发布数量从2019年的2500个翻倍至6000个，竞争白热化。
  \item 欧莱雅旗下Libre上市6年成为全球女性香水第一品牌超越香奈儿，普伊格Good Girl用6年拿下3\%市场份额，爆款能力成为区分赢家与输家的关键。
  \item GS分析显示Interparfums通过持续推出新支柱将Jimmy Choo做到超2亿欧元营收，普伊格计划为Rabanne和Jean Paul Gaultier推新爆款。
  \item GS认为普伊格和Interparfums目前估值未完全反映爆款潜力，若执行到位存在重估空间，欧莱雅因多元化优势享有溢价但历史估值差距仍明显。
  \item 品类增长正常化后，增长动力从“让更多人买”转向“让更多人买我的产品”，供给质量成为核心变量。
  \item GS报告指出欧莱雅在2010年代中期推出La Vie Est Belle带动女性香水两位数增长，历史证明爆款是长期增长引擎。
\end{itemize}
\begin{figure}[htbp]
\centering
\includegraphics[width=0.88\linewidth]{figures/fig_275.jpg}
\caption{Exhibit 1 - Exhibit 1: Beauty market growth rate has accelerated post Covid L'Oréal estimate of global beauty industry sales growth, 2010-25}
\end{figure}
\begin{figure}[htbp]
\centering
\includegraphics[width=0.88\linewidth]{figures/fig_278.jpg}
\caption{Exhibit 4 - Exhibit 4: Largest six players account for two-thirds of premium fragrance sales Global premium fragrances breakdown by company, 2025}
\end{figure}
\begin{figure}[htbp]
\centering
\includegraphics[width=0.88\linewidth]{figures/fig_280.jpg}
\caption{Exhibit 6 - Exhibit 6: Newness continues to resonate, with interest in fresh blockbusters tracking ahead of past launches Indexed searches for Good Girl vs La Bomba in first 12 months}
\end{figure}
\begin{figure}[htbp]
\centering
\includegraphics[width=0.88\linewidth]{figures/fig_282.jpg}
\caption{Exhibit 8 - Exhibit 8: Good Girl surpassed \$3\textbackslash{}\%\$ share in six years Good Girl: women's fragrance market share, 2022 vs 2016}
\end{figure}
{\small\color{refgray}\textbf{References:} [R021] GS：香水市场的下一个增长波，不是品类红利，而是打造爆款的能力}

\section{结语}
综合来看，市场当前的核心矛盾在于：一方面，中国经济正经历从消费刺激到供给端投资的增长逻辑切换，内需疲软与AI出口繁荣形成K型分化；另一方面，全球市场在美伊和平预期下过度定价，但能源、航运和AI基础设施的供给侧约束尚未被充分计入。资产定价逻辑正在从需求驱动转向供给约束驱动，理解这一结构性变化比预测短期数据更为重要。
\newpage
\section{报告覆盖清单}
\begin{itemize}
  \item [R001] 宏观与利率：中国增长逻辑切换与全球通胀滞后效应 / 地产与消费：K型分化加剧，一线回暖难掩全国压力 - 摩根斯坦利：市场真正低估的不是需求疲软，而是政策转向“供给端投资”的结构性含义
  \item [R002] 宏观与利率：中国增长逻辑切换与全球通胀滞后效应 - GS：消费的真正压力尚未到来，市场对能源冲击的滞后效应严重低估
  \item [R003] 地产与消费：K型分化加剧，一线回暖难掩全国压力 - GS：中国家庭资产正在经历从房产到金融资产的结构性迁移
  \item [R004] 宏观与利率：中国增长逻辑切换与全球通胀滞后效应 - GS：油价传导的真正压力不在五月，而在接下来的几个月
  \item [R005] 宏观与利率：中国增长逻辑切换与全球通胀滞后效应 - Citi：市场真正低估的不是AI需求，而是国内“滞胀”的定价权转移
  \item [R006] 能源与大宗：和平交易过度定价，供给约束支撑三季度油价 / 航运与供应链：供给约束驱动定价逻辑重置 - UBS：航运市场的真正拐点不在运价，而在供给侧的不可逆重置
  \item [R007] 汽车行业：补贴退潮暴露结构性压力，中国车企加速全球渗透 - 摩根斯坦利：欧洲车企最坏的时刻尚未到来，中国OEM的份额侵蚀才刚刚加速
  \item [R008] 能源与大宗：和平交易过度定价，供给约束支撑三季度油价 / 航运与供应链：供给约束驱动定价逻辑重置 - JPM：铁矿石定价权正在向澳大利亚矿商倾斜，而非中国需求主导
  \item [R009] AI与算力：基础设施定价权重估与硬件ROE结构性上升 - Bernstein：SOCAMM2不会侵蚀内存接口芯片市场，市场高估了替代风险
  \item [R010] 区域市场：亚洲外汇储备分化与东南亚互联网韧性 - DB：市场低估了资本账户对人民币的支撑力
  \item [R011] 区域市场：亚洲外汇储备分化与东南亚互联网韧性 - NOM：亚洲央行正在用外汇储备为货币信心买单，但印尼和印度的弹药已接近红线
  \item [R012] 生物医药与半导体：全球价值转化面临重估 - Bernstein：中国生物医药被低估的不是创新能力，而是“全球价值转化”面临的真实阻力
  \item [R013] 能源与大宗：和平交易过度定价，供给约束支撑三季度油价 - UBS：中东局势缓和下，化工板块的真正机会不在油价，而在成本端的重新定价
  \item [R014] 能源与大宗：和平交易过度定价，供给约束支撑三季度油价 / 航运与供应链：供给约束驱动定价逻辑重置 - GS：市场真正低估的不是需求崩塌，而是供给恢复的路径与节奏
  \item [R015] 地产与消费：K型分化加剧，一线回暖难掩全国压力 - JPM：中国楼市真正的分水岭不在价格，而在供给侧的“窄复苏”能否撑到政策显效
  \item [R016] AI与算力：基础设施定价权重估与硬件ROE结构性上升 - GS：欧洲数据中心正在重塑电力行业的估值逻辑，而市场仍按旧地图定价
  \item [R017] 地产与消费：K型分化加剧，一线回暖难掩全国压力 - 摩根斯坦利：市场低估了三季度销售二次探底的冲击力
  \item [R018] AI与算力：基础设施定价权重估与硬件ROE结构性上升 / 生物医药与半导体：全球价值转化面临重估 - Bernstein：亚洲硬件正在经历一轮由AI驱动的ROE结构性重定价
  \item [R019] 汽车行业：补贴退潮暴露结构性压力，中国车企加速全球渗透 - Citi：市场低估的不是需求疲软，而是供给端策略的结构性拐点
  \item [R020] AI与算力：基础设施定价权重估与硬件ROE结构性上升 / 生物医药与半导体：全球价值转化面临重估 - Citi：半导体设备市场真正被低估的，是2028年2500亿美元WFE背后的结构性需求切换
  \item [R021] 消费行业：香水市场进入爆款竞赛模式 - GS：香水市场的下一个增长波，不是品类红利，而是打造爆款的能力
  \item [R022] 区域市场：亚洲外汇储备分化与东南亚互联网韧性 - NOM：人民币中间价定价模型释放的信号比汇率本身更重要
  \item [R023] 宏观与利率：中国增长逻辑切换与全球通胀滞后效应 / 风险偏好与政策：银行股结构溢价与香港楼市政策误读 - GS：银行股真正的防御性，来自资产负债表结构而非规模
  \item [R024] 地产与消费：K型分化加剧，一线回暖难掩全国压力 - 摩根斯坦利：消费复苏的真正考验不在总量，而在结构分化的再定价
  \item [R025] 能源与大宗：和平交易过度定价，供给约束支撑三季度油价 - 摩根斯坦利：市场对“和平”的定价已经过头，但油价真正的支撑尚未被计入
  \item [R026] 风险偏好与政策：银行股结构溢价与香港楼市政策误读 - JPM：市场对内地资金买港楼的恐慌，可能高估了政策冲击
  \item [R027] 宏观与利率：中国增长逻辑切换与全球通胀滞后效应 - 摩根斯坦利：中国材料工业的真正分化不在需求，而在供给侧的刚性约束
  \item [R028] 汽车行业：补贴退潮暴露结构性压力，中国车企加速全球渗透 - Citi：市场低估的不是需求疲软，而是供给端正在发生的结构性出清
  \item [R029] AI与算力：基础设施定价权重估与硬件ROE结构性上升 - GS：AI基础设施定价权正在经历结构性重估，但企业级需求仍处于“等待爆发”的临界点
  \item [R030] 区域市场：亚洲外汇储备分化与东南亚互联网韧性 - UBS：东南亚外卖与电商的增长故事并未结束，但竞争正在重新定价龙头门槛
  \item [R031] 汽车行业：补贴退潮暴露结构性压力，中国车企加速全球渗透 - Bernstein：中国车市真正的考验不是需求萎缩，而是补贴退潮后供给侧的再定价
\end{itemize}
\section{逐篇报告摘录}
\subsection{[R001] 摩根斯坦利：市场真正低估的不是需求疲软，而是政策转向“供给端投资”的结构性含义}
{\small\color{refgray}归类：宏观与利率：中国增长逻辑切换与全球通胀滞后效应 / 地产与消费：K型分化加剧，一线回暖难掩全国压力}

摩根斯坦利：市场真正低估的不是需求疲软，而是政策转向“供给端投资”的结构性含义

5月的经济数据，表面上看是“出口强、内需弱”的老故事，但摩根斯坦利这份最新研报揭示了一个更深层的信号：中国经济正在从“消费刺激依赖”转向“供给端投资驱动”，而这一转向的资产定价含义，远未被市场充分消化。

二季度GDP跟踪值降至4.4\%，零售销售自2022年以来首次转负，固定资产投资单月同比跌幅扩大至10.7\%——这些数字本身已经足够令人警惕。但报告真正值得关注的判断，并非这些疲软数据本身，而是它们共同指向的一个政策拐点：北京正在加速放弃“以消费拉动经济”的短期路径，转而将筹码押注在AI算力网络、数据中心、智能电网等“六网”基础设施投资上。

这不是一次简单的财政加码，而是一次增长逻辑的切换。理解这一切换，比预测下个月的经济数据重要得多。

1. 消费刺激的“天花板”已经提前到来，内需疲软正在从短期波动变成结构性特征

5月零售销售同比下滑0.6\%，不仅低于市场预期的-0.5\%，更是2022年以来首次跌入负值区间。摩根斯坦利的分析指出，这并非单一因素所致，而是“以旧换新”政策效果递减与就业市场持续疲弱叠加的结果。

值得注意的一个细节是，5月的“618”电商促销提前了一天启动，但零售数据依然转负。这意味着消费的疲软已经超出了短期促销可以拉动的范围。政策工具对消费的边际效用正在快速递减。

报告进一步拆解了零售数据的内部结构：汽车销售同比下滑16.1\%，住房相关消费下滑14.1\%，黄金珠宝下滑8.9\%。剔除以旧换新相关商品和黄金后的“核心商品零售”同比增长3.5\%，虽然相对稳定，但这一增速也低于4月的4.2\%。

这些数字合在一起意味着什么？消费疲软正在从“周期性波动”演变为“结构性特征”。就业市场的持续疲弱、居民收入预期的下降、以及房地产财富效应的消退，正在形成一种负向循环，而这种循环很难通过短期的补贴政策来打破。

对于投资者而言，这意味着任何基于“消费复苏”逻辑的资产定价，都需要重新审视其假设的可靠性。消费板块的估值修复，可能比市场预期的更加漫长和不确定。

2. 出口与投资的“传导断裂”正在扩大，制造业产能过剩成为压制资本开支的核心阻力

5月工业增加值同比增长4.5\%，略高于市场预期的4.3\%，其中中下游制造业表现相对强劲，增速环比上升0.7个百分点至6.4\%，主要由高技术产业带动。这看起来是一个积极信号。

但摩根斯坦利敏锐地指出了其中的关键断裂：出口带动的生产增长，并没有有效传导至国内资本开支。制造业固定资产投资单月同比依然为负，为-4.2\%，与4月的-4.3\%基本持平。

为什么出口强劲却没有带来投资扩张？报告给出的答案是“普遍性的产能过剩”。在多数制造业领域，产能利用率已经处于低位，企业没有动力进一步扩大产能，即使出口订单表现不错。投资的传导被限制在了计算机和设备等少数领域。

这一发现具有重要的宏观含义：过去市场习惯的“出口好-投资好-就业好-消费好”的传导链条，现在出现了明显的断裂。出口的繁荣变成了“孤岛式”的增长，无法有效带动国内经济的内生循环。

这意味着，对出口相关资产的定价，需要更加谨慎地区分“出口收入增长”与“国内利润改善”之间的差异。那些依赖出口拉动国内投资和就业的行业，其景气度的可持续性值得怀疑。

3. 房地产的“局部回暖”正在被证伪，二手房亮眼数据掩盖了整体市场的脆弱性

房地产市场的表现是这份报告中最具欺骗性的部分。一些头部城市的二手房成交数据近期有所回升，这被部分市场参与者解读为“企稳信号”。

但摩根斯坦利的分析揭示了另一面：二手房销售的局部改善，与疲弱的按揭贷款增长、持续下滑的新房销售、以及仍在加速下跌的房地产投资形成了鲜明对比。5月房地产开发投资单月同比下滑24.3\%，较4月的-20

[... middle omitted ...]

要看到新房销售、土地购置、开发投资三个指标的同步企稳，而目前这三个指标均未出现明确的改善迹象。

4. 财政政策正在经历一次“沉默的转向”：从消费刺激到供给端投资的筹码转移

这是整份报告中最具前瞻性的判断，也是市场尚未充分定价的部分。

5月基础设施固定资产投资单月同比下滑10.8\%，较4月的-3.7\%大幅恶化。政府债券发行节奏明显放缓，反映出1季度前置发债后的项目储备出现空窗期。这看起来是财政扩张力度减弱的信号。

但摩根斯坦利认为，这恰恰是政策重心切换的过渡期。报告指出，目前仍有约60\%的年度政府债券额度尚未使用。预计7月的政治局会议将要求加快财政投放节奏，但重点不是新的刺激方案，而是现有额度的加速执行。

更关键的是政策投向的变化：报告明确指出，政策焦点将集中在资本开支而非消费上——特别是支持AI算力网络、数据中心和智能电网等“六网”倡议下的项目。这一选择的背后逻辑是：在地缘政治紧张加剧的背景下，能源安全和技术自主的战略必要性正在超越短期消费拉动的政策优先级。

\subsection{[R002] GS：消费的真正压力尚未到来，市场对能源冲击的滞后效应严重低估}
{\small\color{refgray}归类：宏观与利率：中国增长逻辑切换与全球通胀滞后效应}

GS：消费的真正压力尚未到来，市场对能源冲击的滞后效应严重低估

过去几周，能源市场的消息面似乎正在转好。美伊临时和平协议降低了油价的上行尾部风险，GS的大宗商品策略师已将2026年第四季度布伦特原油预期从此前的每桶90美元下调至80美元。与此同时，大西洋两岸的消费者支出数据至今表现出了令人意外的韧性。这些信号叠加在一起，很容易让市场参与者产生一种判断——能源冲击对增长的影响已经见顶，风险正在消退。

但这份来自GS全球经济学团队的最新报告给出了一个截然不同的结论。报告的署名作者包括首席经济学家Jan Hatzius，其核心判断非常清晰：现在放松警惕为时过早。真正的消费压力尚未到来，甚至可以说，主要冲击波还没有真正进入经济数据。

市场低估的不是能源价格本身，而是能源冲击向消费者现金流传导的机制、时间差和季节性放大效应。这不是一个线性过程，而是一个在第二、三季度才会集中释放的滞后冲击。对于关注宏观资产定价、行业盈利预期和消费类资产基本面的投资者而言，理解这个传导链条的节奏，比盯着油价短期走势更为重要。

GS的分析框架建立在三个相互独立的机制之上，它们共同指向同一个结论：消费的实质性放缓还没有开始。

第一条路径是批发价格向零售价格的传导延迟。汽油价格在冲突爆发后迅速跳升，但天然气和电力价格向消费者的传导更为缓慢。GS欧洲经济学家指出，固定合同和政府监管机制在历史上一直导致消费端价格调整滞后于批发价格。这意味着，相当一部分相关的消费品价格涨幅尚未实现。

第二条路径是季节性需求对现金流的放大效应。能源价格上涨对家庭现金流的冲击，在冬季取暖需求高峰时会显著放大。GS的报告用了一个非常精巧的计算来展示这一点：在固定消费结构假设下，欧洲的冲击谷底出现在2026年中；但一旦引入冬季需求的季节性波动，谷底将后移至2027年第一季度。季节性因素本身在欧洲就贡献了超过20\%的额外拖累。

第三条路径是美国特有的财政脉冲反转。今年上半年，超大规模的退税暂时将美国实际现金流推高了0.8\%（第一季度）和1.2\%（第二季度）。但这一临时性提振已经结束，第三季度同比增速将暂时转负。GS的计算显示，这将导致美国实际现金流从第二季度的正增长0.7\%急剧转为第三季度的负增长0.6\%。

这三条路径的共振窗口不在过去，不在当下，而在下半年。这是报告最核心的时间判断。

2. 历史经验给出了一个可量化的传导公式：每1\%的现金流损失，对应0.6\%的消费下降

GS基于历史统计关系建立了一个经验模型。该模型的参数简洁而有力：能源冲击导致的每1\%实际现金流下降，会在两个季度后转化为约0.6\%的消费支出下降。

这个数字本身并不惊人，但关键在于它解释了为什么到目前为止消费数据看起来还不错。GS的模型能够很好地拟合现有的有限消费影响，但更重要的是，它在新的能源价格假设下预测出下半年将出现更明显的疲软。具体而言，GS估计实际支出的峰值拖累在加拿大为0.5\%，欧元区为0.4\%，美国和英国为0.3\%。

这些数字单独看或许不算灾难性，但放在当前市场对“软着陆”的普遍预期背景下，它们意味着一个重要的边际变化。市场已经定价了经济韧性，但并未完全定价一个由能源滞后传导驱动的、结构性的消费放缓。

更值得关注的是，GS估计美国尚未出现任何实质性的支出拖累，而欧元区、英国和加拿大到第二季度末也仅实现了峰值冲击的一半左右。换句话说，消费数据中最坏的部分还在后面。

3. 消费者调查已经给出了预警，但市场可能忽视了其结构性含义

报告引用了一项前瞻性调查证据，这往往是市场最容易忽视但最具信号意义的数据。大多数发达市场的消费者预计，未来6到12个月的大额采购将少于冲突爆发前。美国的消费意愿在春季因财政支持而一度回升，但5月份已经回落。

GS的美国经济学家认为，这与他们的估计

[... middle omitted ...]

放缓叙事”的切换过程。

4. 报告尚未完全回答的关键问题：中国库存重建与欧洲冬季的叠加效应

报告中有一个看似不起眼的细节值得单独讨论。GS的大宗商品策略师虽然下调了油价预期，但并未改变天然气年末价格预测，原因在于“库存重建”。报告提到，中国在春季去库存后，将在冬季前开始重建库存。

这是一个尚未被充分讨论的变量。中国的天然气库存行为对全球价格有显著影响。如果中国在冬季前大规模补库，叠加欧洲的取暖需求，天然气的价格压力可能超出GS当前基准假设。报告没有对此进行情景分析，但读者需要意识到，当前的能源价格预测存在一个重要的上行风险。

这意味着，GS报告中对消费拖累的量化估计，可能是一个偏保守的基准情景。如果天然气价格在冬季前出现二次上涨，那么报告所描述的“滞后冲击”将不仅是幅度上的增加，还可能是时间上的延长。

对于产业决策者而言，这个未解问题意味着不应简单依赖基准预测来制定采购和库存策略。对于投资者而言，这意味着当前能源股和防御性消费股的相对吸引力可能需要重新评估。

\subsection{[R003] GS：中国家庭资产正在经历从房产到金融资产的结构性迁移}
{\small\color{refgray}归类：地产与消费：K型分化加剧，一线回暖难掩全国压力}

中国家庭资产负债表正在经历一个过去二十年从未出现过的变化：房地产从财富增长的引擎变成了拖累，而金融资产正在接替这一角色。这并非一个短期波动，而是一个结构性转折的早期阶段。

GS在最新发布的研报中构建了一个季度频率的家庭资产负债表追踪框架，填补了官方数据发布滞后的空白。根据这一框架，中国家庭总资产在2023年初见顶，经历了约六个季度的收缩后，目前在730万亿元人民币左右企稳。但比总量企稳更值得关注的，是资产构成正在发生的根本性变化。

截至2026年第一季度，GS估算房地产在中国家庭总资产中的占比已从2021年中房地产市场高峰期的67\%降至52\%，而现金及存款从16\%升至25\%，其他金融资产（含股票、债券、基金、保险等）从15\%升至20\%。直接持有的股票在总资产中的占比也从5\%微升至6\%。

这些数字本身已经说明问题，但真正重要的是它们背后的含义：中国家庭的储蓄和财富配置正在经历一个从“不动产信仰”到“金融资产再平衡”的早期阶段。这一过程才刚刚开始，其对消费、资本市场乃至宏观经济的影响，可能比多数投资者预期的更为深远。

1. 房地产从财富创造者变为财富拖累，这一转变尚未被市场充分定价

GS的数据显示，在2021年房地产见顶之前，中国家庭资产负债表持续扩张的核心驱动力是房价上涨和房产购买积累。当房价从高点下跌约30\%后，这一逻辑彻底逆转。过去几年，房价下跌造成的资产减值超过了金融资产的增值，导致家庭总资产在2023年至2024年持续收缩。

这一转变的核心含义在于：房地产不再是中国家庭财富增长的稳定器。过去二十年，买房几乎等同于“储蓄+投资”的合一体，而现在，它变成了一个需要持续消耗现金流来维护的负债项。即使房价不再继续大幅下跌，只要不涨，其对家庭财富增长的贡献就是零甚至负值。

GS报告指出，目前房地产仍占家庭总资产的约60\%，而股票占比不到10\%。更重要的是，超过90\%的家庭拥有房产，而只有约25\%的成年人参与股市。这意味着房价疲软对大众消费的影响远大于股市上涨。金融资产的增值主要惠及高收入群体，对整体消费的拉动作用相对有限。

这里有一个尚未被市场充分讨论的推论：如果房地产的财富创造功能永久性减弱，中国家庭的边际消费倾向可能会系统性下降。过去人们敢于消费，部分原因是房产升值带来的“财富幻觉”支撑了信心。当这一支柱消失后，家庭可能需要更长时间来重建资产负债表的安全感。

GS的数据显示，中国家庭债务自2024年中以来开始趋势性下行，主要驱动力是存量房贷余额的下降。提前还贷和购房需求低迷共同导致了住房相关借贷的收缩。短期消费贷款也在走弱，反映家庭部门在住房之外的消费需求同样谨慎。

从总量看，中国家庭债务占GDP比率在2024年初见顶于61\%，到2025年第三季度已回落至约59\%。这一水平低于美国（约70\%）和日本（约62\%），但高于新兴市场平均水平（约45\%）。表面上看并不算极端。

但GS报告特别指出，如果以债务收入比来衡量，中国家庭的偿债压力仍然较高。这才是问题的核心：中国家庭债务的“分母”是收入，而收入增长在过去几年明显放缓。即使债务总量在减少，如果收入增长更慢，偿债负担反而会加重。

这一判断对投资者的启示在于：政策刺激的效果可能被高估。降低房贷利率和推出消费贷利息补贴，都是在降低债务成本，但如果家庭对未来收入预期没有改善，他们更可能选择提前还贷而非增加消费。去杠杆是一种行为模式转变，不是利率工具可以轻易逆转的。

GS报告中最具前瞻性的判断是：中国家庭资产配置正处于结构性迁移的早期阶段——随着房地产在财富积累中的角色减弱，存款利率持续走低，储蓄将逐步向更广泛的金融资产迁移，股票和保险是中期最有可能受益的领域。

从数据看，这一迁移已经在发生。自2024年9月政策转向以来，股票对家庭金融资产增长的

[... middle omitted ...]

续下行的背景下，这些产品的相对吸引力在上升，但投资者对风险的接受度仍有限。

4. 报告没有完全回答的一个关键问题：财富效应的分布不均如何影响宏观

GS报告明确指出了财富效应分布不均的问题：金融资产的上涨主要惠及高收入群体，而房价下跌对大众消费的冲击更大。但报告没有深入展开的是，这种分布不均对宏观经济的结构性影响。

一个合理的推论是：如果金融资产的上涨集中在少数人手中，而房地产的下跌影响多数人，那么即使家庭总资产企稳，消费倾向也可能继续分化。高收入群体可能增加储蓄或投资，而非消费；中低收入群体则因资产负债表受损而持续压缩支出。这种“K型修复”可能导致整体消费复苏慢于资产价格修复的速度。

另一个未完全回答的问题是：家庭资产从房产向金融资产的迁移，是否会改变中国的储蓄率结构？过去二十年，高储蓄率部分与买房需求有关——人们储蓄是为了凑首付。如果购房需求系统性下降，这部分预防性储蓄是否会释放为消费，还是会以其他形式继续沉淀？这取决于社会保障体系的完善程度和收入增长预期。

\subsection{[R004] GS：油价传导的真正压力不在五月，而在接下来的几个月}
{\small\color{refgray}归类：宏观与利率：中国增长逻辑切换与全球通胀滞后效应}

通胀数据正在讲述一个比表面更复杂的故事。澳大利亚四月CPI同比降至4.2\%，环比仅微增0.4\%，看起来通胀压力在缓解。但这份GS研报的核心判断是：五月的数据很可能掩盖了一个正在积聚的结构性风险——油价下跌带来的短期通缩效应，恰好与成本端持续上升的传导压力形成错位，真正的通胀读数将在未来几个月重新显现。

市场容易犯的错误是线性外推。看到汽油价格单月下跌12.5\%，就认为通胀压力消退。但研报揭示的逻辑链条恰恰相反：油价下跌是一次性的政策与地缘因素叠加，而材料成本、化工品价格、燃料附加费的传导才刚刚进入企业定价决策的窗口。这才是真正需要关注的信号。

GS对五月澳大利亚CPI的预测是环比下跌0.4\%，同比维持4.2\%不变。这个看似温和的数字背后，是极为特殊的结构性因素在起作用。

汽油价格预计单月暴跌12.5\%，这并非需求疲软所致，而是三个月临时燃油消费税减免政策的效果集中释放。与此同时，国内机票价格在复活节后季节性回落7.2\%。这两项合计足以把整体CPI拉低约0.5个百分点。

但剔除这些一次性因素后，核心指标反而在走强。GS预计五月截尾均值CPI环比上涨0.3\%，同比从3.4\%升至3.5\%。更重要的是，其偏好的三个月对三个月截尾均值指标将升至0.88\%季度环比——这个数字已经接近RBA关注的门槛。

这意味着，如果只看五月整体CPI，你会得出通胀继续放缓的结论。但拆解后会发现，内生的价格压力不仅没有减轻，反而在加速。这是典型的“数据假摔”，也是这份报告最值得警惕的判断。

研报中一个容易被忽视但极为关键的证据，来自建筑材料供应链的数据。主要管道和建材供应商宣布的价格上涨次数和平均涨幅，已经超过了疫情后的峰值。这不是短期波动，而是成本端的结构性抬升。

GS对此做了量化测算：如果材料与化工品价格上涨10\%，对食品CPI的贡献约为2个基点，对住房CPI的贡献约为1.7个基点。这看起来不大，但关键在于传导的滞后性和累积性。

化肥价格目前仍处于相对高位，尿素和氨的价格在2026年的预测值分别为850美元/吨和750美元/吨，显著高于2024年的水平。这些成本最终会渗透到农产品价格和食品加工环节。报告明确指出，住房和食品价格将受到输入性成本上升的推动——这恰恰是CPI篮子中权重最高、消费者感受最直接的两个分项。

这里尚未完全回答的关键问题是：企业是否有足够的定价权将成本完全转嫁给消费者。如果需求端走弱，利润率将承压；如果需求韧性足够，通胀就会比预期更持久。研报暗示了前者，但没有给出明确判断。

五月数据中有一个值得追踪的新现象：餐饮行业开始出现燃料附加费。GS预计餐厅用餐和外卖价格在五月均环比上涨0.7\%，部分原因正是燃料附加费的引入。

这标志着油价传导进入第二阶段。第一阶段是汽油和柴油价格直接上涨，消费者在加油站就能感知；第二阶段是燃料成本通过运输、冷链、包装等环节，渗透到更广泛的消费品和服务价格中。新西兰的数据已经印证了这一趋势——当地餐厅用餐价格环比上涨0.7\%，创下2023年中以来的最快增速。

更值得关注的是，新西兰五月月度CPI数据显示，剔除燃料后的同比通胀率反而从2.0\%跳升至3.0\%。这意味着燃料价格下跌的同时，其他商品和服务的价格压力反而在扩大。这与“油价下跌导致通胀全面回落”的简单叙事完全相反。

报告没有充分展开的是：燃料附加费是否会从餐饮行业扩散到零售、物流、制造业等其他领域。如果这种传导模式被证实具有普遍性，那么未来几个月的核心通胀读数将面临持续上行风险。

澳大利亚住房通胀是另一个容易被低估的变量。四月数据显示，新房购买价格通胀进一步加速，这反映的是燃料附加费和材料成本上涨的滞后传导。GS预计五月新房购买价格环比将加速至0.9\%。

与此同时，广告租金增长在五月进一步回升。虽然CPI中的

[... middle omitted ...]

。

研报提供了一组值得深思的数据对比：短期通胀预期指标虽然从峰值回落，但仍远高于RBA的目标区间。消费者预期的通胀率仍在6\%-7\%的区间，市场经济学家的预期也在5\%-6\%之间。

但真正让央行安心的信号来自长期通胀预期。无论是3年期还是10年期通胀互换，还是市场经济学家的长期预测，都稳定在RBA目标区间的中点附近，大约2.5\%。这意味着市场仍然相信央行有能力在长期内控制通胀。

然而，这里存在一个微妙的博弈：如果短期预期持续高企的时间过长，它会不会最终侵蚀长期预期的锚定？研报没有给出答案，但这个问题对资产定价至关重要。一旦长期通胀预期开始松动，债券收益率将面临重新定价，而股票市场的估值逻辑也将随之改变。

6. 报告尚未完全回答的关键问题：油价下跌的“假摔”何时会被市场纠正

这份研报最值得反复阅读的部分，其实是GS自己标注的两个不确定性来源：一是燃油消费税减免和全球油价波动叠加带来的高波动性；二是国际机票价格因采集周期滞后两个月，五月数据才首次反映地缘冲突的影响。

\subsection{[R005] Citi：市场真正低估的不是AI需求，而是国内“滞胀”的定价权转移}
{\small\color{refgray}归类：宏观与利率：中国增长逻辑切换与全球通胀滞后效应}

Citi：市场真正低估的不是AI需求，而是国内“滞胀”的定价权转移

理解当前中国经济的核心框架，已经从“复苏还是放缓”切换到了“K型分化如何定价”。Citi这份5月经济指标解读报告给出了一个锋利但容易被忽视的判断：AI超周期与国内滞胀正在成为两条平行叙事，而市场对前者的关注度远高于后者，这恰恰可能隐藏着更大的预期差。

这份报告的核心洞察在于，它没有停留在“外需强、内需弱”的简单二分法，而是明确指出，国内需求侧正在经历一个自疫情以来首次出现的名义收缩——零售销售转负、固定资产投资降幅扩大，而CPI与PPI的组合却指向了滞胀风险。这不是一个短期波动，而是一个结构性信号：中国经济的“K”字，正在被AI出口和国内滞胀这两股力量拉得更开。

对于产业决策者和投资者而言，理解这个框架比预测下个月的GDP数字更重要。因为这意味着，资产的定价逻辑正在发生根本性迁移——过去几年“内需复苏”的交易逻辑正在失效，而“AI供给”和“滞胀防御”将成为新的主线。

1. AI超周期对工业生产的拉动已经超过50\%，但这是“有产量、无需求”的繁荣

Citi报告中最值得关注的数据点之一，是高新技术产业增加值同比增速达到15.1\%，创下五年新高。更关键的是，按照该机构的估算，高技术产业在整体工业增加值中的占比约为17\%，却贡献了超过50\%的工业增长。集成电路产量同比增长22.9\%，工业机器人增长27.9\%，新能源汽车增长17.8\%。

这些数字本身并不令人意外，市场对此已有充分预期。真正值得推敲的是，这种增长的质量和可持续性。Citi报告提供了一个被很多人忽略的视角：销售-生产比率。5月份这一比率为96.1\%，低于2023年和2024年同期的96.6\%和96.5\%。换句话说，生产端在加速，但销售端的消化能力并未同步跟上。

这意味着，AI驱动的出口和生产繁荣，本质上是一种“供给侧的自我强化”，而非需求拉动。企业为了抢占全球AI基建市场份额而扩大产能，但如果海外需求增速边际放缓，或者地缘政治因素导致贸易壁垒升级，这部分产能的消化压力将迅速传导回国内。这不是危言耸听，而是当前K型经济中“出口强、内需弱”结构的天然脆弱性。

2. 零售销售转负不是消费降级，而是“补贴退坡”与“收入预期”的双重挤压

5月社会消费品零售总额同比下降0.6\%，这是自疫情以来首次录得负增长。但Citi报告揭示，这背后并非单一的消费信心崩塌，而是一个结构性问题：以旧换新补贴的退坡效应正在加速显现。

报告估算，以旧换新相关品类对零售的拖累从4月的-2.0个百分点扩大至5月的-2.2个百分点。汽车销售同比下降16.1\%，家电下降15.6\%，通讯器材从6.2\%骤降至0.7\%。这些品类恰好是过去一年政策刺激的主要受益者。当补贴力度减弱、高基数效应叠加，名义增速迅速转负。

但更值得关注的是，Citi报告指出，即使剔除补贴相关品类，其他消费领域也并未显示出韧性。餐饮收入仅增长0.6\%，体育娱乐用品下降8.0\%，家居装饰下降13.6\%。这指向了一个更深层的问题：居民收入预期和企业盈利状况并未实质性改善，补贴只是暂时“借”来了未来的消费，而非创造了新的购买力。

对于投资者而言，这意味着消费板块的估值修复逻辑需要重新审视。如果补贴退坡后的真实需求比市场预期的更弱，那么当前股价所隐含的“均值回归”假设可能过于乐观。真正的拐点，或许要等到就业市场和工资增长出现实质性改善。

3. 固定资产投资月度降幅接近两位数，但“六网”投资尚未形成实物工作量

固定资产投资累计同比降幅从4月的-1.6\%扩大至5月的-4.1\%，而Citi估算的月度同比增速更是恶化至-10.7\%，为2025年四季度以来的最低水平。基础设施投资累计增速从4.3\%骤降至0.6\%，制造业投资也仅勉强维持。

这组数据中有一个

[... middle omitted ...]

房价格环比上涨的城市数量从4月的12个进一步减少至5月的10个，三线城市的价格下行压力仍在加深。

对于资产定价而言，这个格局意味着：短期内，固定资产投资相关的大宗商品、工程机械等周期板块，缺乏需求端的支撑。而房地产链条上的资产，即便在一线城市出现企稳信号，也不应被解读为全行业的拐点。真正的右侧信号，需要看到土地市场和新开工面积的连续改善。

4. 报告没有完全回答的关键问题：PPI回升是供给驱动还是需求驱动？

Citi报告在讨论滞胀风险时，有一个非常敏锐的观察：5月PPI的回升并非完全由中东冲突导致的能源价格推动。乙烯产量已经从-4.1\%回升至2.1\%，硫酸产量降幅也在收窄，这意味着上游供给约束正在缓解。但PPI仍然在上升，这指向了一个需要进一步拆解的问题。

报告没有完全展开讨论的是，当前PPI回升的驱动因素中，有多少来自AI超周期带来的特定行业资本开支需求（比如通信设备、半导体材料），有多少来自国内“反内卷”行动对部分行业供给的控制，又有多少是真正的需求拉动。

\subsection{[R006] UBS：航运市场的真正拐点不在运价，而在供给侧的不可逆重置}
{\small\color{refgray}归类：能源与大宗：和平交易过度定价，供给约束支撑三季度油价 / 航运与供应链：供给约束驱动定价逻辑重置}

这份最新发布的UBS航运供应链周报，覆盖了截至2026年6月12日的高频数据。表面上看，它是一份常规的周度数据更新，涵盖了VLCC油轮、集装箱、干散货和造船四大板块。但如果你只是把它当作运价追踪来看，就会错过一个更重要的信号：全球航运市场正在经历一轮结构性重置，而这个重置的驱动力，不是短期需求波动，而是供给端的多项不可逆变化正在同时发生。

霍尔木兹海峡的重新开放预期、红海绕航的持续、造船订单的集中爆发、以及集装箱运价的连续跳升——这些事件在市场上被分别讨论，但UBS这份报告将它们放在一起，揭示了一个更完整的图景：航运市场的定价逻辑，正在从“需求驱动”转向“供给约束驱动”。对于产业决策者和投资者来说，这意味着过去几年建立的分析框架需要重新校准。

1. 霍尔木兹海峡的“重启”已被定价，但真正影响运价的不是通航，而是船队效率的永久性损失

报告中最引人注目的数据点之一，是美伊和平协议签署后，油轮和航运股的反应积极。市场显然在押注霍尔木兹海峡重新开放后，VLCC运价将回归正常。但UBS的高频数据揭示了一个不同的现实：截至6月初，霍尔木兹海峡的原油油轮通行量仍比冲突前水平低95\%，几乎处于冻结状态。

更值得关注的是，即便海峡完全重启，船队效率的损失可能已经不可逆。过去一年，大量油轮被迫绕行好望角，这不仅拉长了航程，还改变了全球油轮的调度节奏。船东为了应对不确定性，已经调整了航线规划、船员配置和保险安排。这些调整一旦固化，即便地缘政治风险消退，船队也不会立即回到原来的效率水平。

报告显示，VLCC平均收益目前仍低于10万美元/天，中东/西非至中国的VLCC收益保持平稳。这意味着市场已经将“重启”预期部分计入价格，但尚未充分反映船队效率损失带来的供给硬约束。这是一个典型的“预期已经price in，但基本面尚未兑现”的窗口。

2. 集装箱运价的连续跳升不是季节性波动，而是“绕航常态化”正在重塑运力供需平衡

集装箱板块的数据同样值得深挖。报告显示，上海出口集装箱运价指数（SCFI）上周环比上涨9\%，同比上涨43\%。其中，上海至欧洲航线环比大涨18\%，跨太平洋航线环比上涨10-12\%。这些数字远超正常的季节性波动。

UBS的数据清晰表明，红海绕航仍在持续，且绕航量仍比去年同期高出4-5\%。这个数字看似不大，但关键在于持续时间。绕航已经持续超过一年，船公司已经将其纳入常态运营。这意味着，即便红海局势未来缓和，船公司也不会立即恢复原有航线，因为客户供应链已经适应了新的运输时间和成本结构。

报告还揭示了一个重要细节：亚洲至美国航线的集装箱运量自5月以来同比强劲增长。这背后既有传统旺季的拉动，也有“前置装运”效应——进口商为了规避关税风险，提前发货。这种前置装运会短期推高运价，但也会透支下半年的需求。UBS的“Maritime Destination Monitor”数据显示，这种前置模式已经非常明显。

对于投资者而言，关键问题不是运价能涨多高，而是这种运价水平能否持续。如果绕航常态化叠加前置装运，那么运价的支撑力可能比市场预期的更持久。

干散货板块上周出现回调，波罗的海干散货运价指数（BDI）环比下跌10\%，好望角型船平均收益环比下跌18\%。但报告同时指出，BDI年初至今仍累计上涨61\%。这个数字提醒我们，回调只是高位震荡，而非趋势逆转。

真正需要关注的，是造船市场的信号。报告中的新造船价格指数环比持平，但全球造船需求在VLCC订单的驱动下表现强劲。从历史数据看，造船订单的集中爆发往往意味着未来2-3年的运力供给将大幅增加。但这一次，情况可能不同。

一方面，环保法规（如EEXI和CII）正在加速老旧船型的淘汰。另一方面，船厂产能已经连续多年收缩，即便订单增加，交付周期也在拉长。这意味着，新订单转化为实际运力的速

[... middle omitted ...]

球经济放缓也可能压制货运需求。这些因素都有可能暂时压低运价，但长期来看，供给侧的约束才是决定运价中枢的核心变量。

UBS的报告没有给出明确的判断，但它的数据为这个判断提供了基础。对于读者来说，接下来的几周数据至关重要：如果运价在高位持续震荡而非快速回落，那么“运价中枢上移”的假设将得到验证。

5. 投资者的观察框架应该从“运价涨跌”转向“供给约束的边际变化”

基于以上分析，我们建议读者调整对航运板块的观察框架。过去几年，市场习惯用“需求是否强劲”来预测运价走势。但在当前阶段，需求端的变化已经不足以解释运价的波动。供给侧的边际变化——船队效率、绕航比例、造船订单交付周期、环保法规执行力度——才是更重要的变量。

第一，红海绕航量的变化趋势。如果绕航量持续高于去年同期水平，说明供给约束仍在强化。第二，新造船订单的交付周期是否在拉长。这将直接影响未来2-3年的运力供给。第三，VLCC和集装箱运价的周度波动是否在收窄。如果波动率下降而中枢上移，说明市场正在形成新的均衡。

\subsection{[R007] 摩根斯坦利：欧洲车企最坏的时刻尚未到来，中国OEM的份额侵蚀才刚刚加速}
{\small\color{refgray}归类：汽车行业：补贴退潮暴露结构性压力，中国车企加速全球渗透}

摩根斯坦利：欧洲车企最坏的时刻尚未到来，中国OEM的份额侵蚀才刚刚加速

这份来自摩根斯坦利欧洲与亚洲汽车团队的联合报告，试图回答一个正在被市场低估的问题：中国车企进军欧洲，究竟会给欧洲传统车企的盈利带来多大冲击？结论是明确的——共识预期仍然过于乐观，欧洲车企FY27的盈利还有超过10\%的下行风险。但更值得关注的，不是这个数字本身，而是报告揭示的结构性变化：市场份额的流失正在从个别车企扩散到全行业，而欧洲车企的防御工事，可能比它们自己想象的要脆弱。

市场当前对欧洲汽车股的定价，隐含了一个关键假设：过去几年中国车企在欧洲的份额增长已经趋于平稳，最激烈的竞争阶段已经过去。摩根斯坦利认为，这个假设是错的。报告通过拆解中国车企的产品布局、渠道扩张、本地化策略，以及欧洲各国市场的差异化渗透节奏，得出了一个更令人不安的判断：最坏的时刻尚未到来，中国车企的攻势才刚刚进入加速期。

报告最核心的量化结论是：欧洲OEM在EU5（德国、英国、法国、意大利、西班牙）的市场份额，将从2025年的67\%进一步下滑至2027年的62\%。这个数字本身或许看起来不算剧烈——五年间仅下降5个百分点。但关键在于，此前的份额流失（ 年下降5个百分点）主要集中在Stellantis一家身上，其他欧洲车企的份额基本保持稳定甚至有所增长。而摩根斯坦利认为，接下来的5个百分点流失将更具“普惠性”，波及几乎所有欧洲车企。

这意味着什么？此前市场可以将Stellantis的困境视为“个案”——一家整合不利、产品线老化、在意大利和法国本土市场受到冲击的特定公司。但中国车企正在实现从“点状突破”到“面状渗透”的转变。它们在英国、意大利、西班牙的份额已经上升了7-10个百分点，而德国和法国这两个欧洲最大的汽车市场，虽然目前看起来相对坚挺，但报告暗示，压力正在积聚。

对投资者的含义是：不能再用“Stellantis的问题”来框架化中国车企的威胁。这是一场结构性、全行业的份额再分配。任何基于“欧洲车企份额将稳定在当前水平”的盈利预测，都可能面临系统性下调。

中国车企并非盲目进入欧洲市场。它们的产品策略高度聚焦于三个维度：车型（SUV）、动力系统（纯电和混动）、价格带（入门级）。这三个维度的交叉点，恰好是欧洲车企利润最丰厚的部分。

SUV是欧洲最畅销、利润率最高的车型类别之一。大众、奔驰、宝马的核心利润都高度依赖SUV产品线。中国车企正在以极具竞争力的价格，向C级和D级SUV市场发起冲击。报告特别指出，大众在“大众市场SUV”这一细分领域面临直接竞争——这恰恰是大众集团全球利润的重要支柱。

入门级市场则是另一个被精准锁定的目标。雷诺是唯一一个在入门级市场有显著敞口的欧洲主流车企。中国车企在这个价格带的攻势最为猛烈，这意味着雷诺面临的不是边缘竞争，而是对其核心客户群的正面争夺。

动力系统的选择同样关键。中国车企初期以纯电动车（BEV）切入欧洲，但报告观察到，它们正在快速转向混合动力（HEV/PHEV），以适应当前欧洲市场对混合动力车型需求快速增长的趋势。这种灵活性——能够根据市场需求快速调整产品组合——是许多欧洲传统车企所不具备的。

对产业决策者的启示是：中国车企的竞争优势不仅仅在于成本或电动化，更在于其“精准定位+快速迭代”的市场策略。它们不是在和欧洲车企进行全面的、同质化的竞争，而是在每一个关键利润节点上集中火力。

3. 防御工事的有效性需要重新评估：品牌忠诚度、残值、政策保护能撑多久？

欧洲车企并非没有防御手段。报告列举了多个结构性保护因素：欧洲消费者对本土品牌的忠诚度、传统车企在金融和租赁服务上的渠道优势、较高的二手车残值、以及欧盟可能采取的贸易保护措施。这些因素在过去几年里确实减缓了中国车企的渗透速度。

但报告的核心洞察在于：这些防御工事的有

[... middle omitted ...]

对电动化接受度高时，中国车企的竞争力能够充分释放。

报告尚未完全回答的一个关键问题是：欧洲车企正在推进的成本削减和平台整合（如LFP电池的本地化、软件定义汽车的开发），能否在时间窗口关闭之前，有效缩小与中国车企的性价比差距？这是一个关于“追赶速度”的假设检验。

4. 盈利预测的下调才刚刚开始：为什么摩根斯坦利只调了宝马和奔驰？

值得注意的一个细节是，摩根斯坦利在此次报告中仅下调了宝马和奔驰的盈利预测（宝马2026年EPS下调4\%，奔驰2027年EPS下调6\%），而对大众、Stellantis、雷诺的预测保持不变。这并非意味着后者的风险更低。

相反，报告的逻辑是：大众和雷诺等大众市场车企面临的竞争压力，已经在此前的低于共识的预测中得到了较为充分的反映。而宝马和奔驰此前被认为受到高端品牌的保护，风险相对较低。现在，随着中国车企将产品线向上延伸至D-SUV和C-SUV领域，高端品牌也开始感受到定价压力——虽然尚未出现明显的市场份额流失，但定价能力可能先于份额被侵蚀。

\subsection{[R008] JPM：铁矿石定价权正在向澳大利亚矿商倾斜，而非中国需求主导}
{\small\color{refgray}归类：能源与大宗：和平交易过度定价，供给约束支撑三季度油价 / 航运与供应链：供给约束驱动定价逻辑重置}

JPM：铁矿石定价权正在向澳大利亚矿商倾斜，而非中国需求主导

这份来自JPMEMEA金属与采矿团队的周度渠道检查报告，表面上是在追踪中国钢铁产量的短期高频数据，但真正值得关注的信号隐藏在运费结构的变化之中。报告覆盖了截至6月10日的中国钢铁产出、铁矿石海运成本、港存变化以及钢厂利润动态，但其核心洞察并非关于中国需求是否疲软，而是关于全球铁矿石供给侧的竞争格局正在发生一次微妙但重要的再定价。

过去一周，从澳大利亚到中国的铁矿石散货运费骤降22\%，至每吨10.9美元。与此同时，巴西和南非至中国的运费虽然也录得周度环比下降，但依然高于冲突前水平超过50\%。这两个事实放在一起，意味着澳大利亚矿商正在获得显著的物流成本优势，而巴西和南非的生产商则面临持续的利润率压力。JPM据此对英美资源集团和库博铁矿给出了明确的Underweight评级，对必和必拓的澳大利亚上市主体和力拓有限公司则维持Overweight。

这不是一份关于中国钢铁需求触底的报告。这是一份关于全球铁矿石供应链如何在运费波动中重新分配利润的报告。市场如果仅仅关注中国产出数据，很可能低估了供给侧结构性差异对资产定价的影响。

1. 中国钢铁产出并未超预期，但市场对“需求见底”的叙事需要更谨慎

报告最直观的数据是：截至6月10日的10天中国日均粗钢产量年化约10.18亿吨，较前一个10天周期环比增长4\%，同比增长2\%。这一水平处于过去五年区间的下沿。乍看之下，这似乎支持了“中国钢铁需求正在企稳”的判断，但JPM并没有给出乐观结论。

关键问题在于，产出增长更多是季节性因素驱动，而非需求实质性复苏。报告提到，过去30天的滚动产量较前30天仍为环比下降1\%，同比仅增长1\%。更重要的是，中国钢铁出口在5月的年化运行率已高达1.22亿吨，处于历史均值上沿。这意味着国内产量的一部分正在被出口消化，而非国内终端需求拉动。

对于钢铁行业决策者而言，这一数据的含义是：不能简单地将产出企稳解读为需求回暖。出口的高位运行本身也面临贸易摩擦风险，一旦主要市场加征关税，国内供给压力将迅速传导至价格和利润。

这是整份报告最值得反复推敲的部分。澳大利亚至中国的散货运费周度暴跌22\%，而巴西和南非至中国的运费仅下降3\%。结果是，澳大利亚FOB铁矿石价格约为每吨90美元，较2月底（美伊冲突开始时）反而上涨了4\%。

运费分化的原因报告没有详细展开，但可以从逻辑上推断：澳大利亚的出口航线更短、船型更标准化、港口基础设施更成熟，因此在全球航运市场波动时，其运费调整的弹性和速度远优于巴西和南非。后者不仅航线更长，还面临港口拥堵、天气干扰以及地缘政治风险溢价。

这一结构变化的直接后果是：澳大利亚矿商在CFR定价体系下的边际利润正在扩大，而巴西和南非矿商则被夹在更高的运费成本和相对刚性的CFR价格之间。JPM对英美资源和库博铁矿的Underweight评级，正是建立在这一运费劣势将长期压制利润率的前提之上。

对于资产配置者而言，这意味着铁矿石板块的内部轮动逻辑已经改变。过去市场习惯于将铁矿石股票视为同质化的商品敞口，但运费差异正在制造显著的alpha机会。澳大利亚矿商的风险溢价正在下降，而巴西和南非矿商的风险溢价则被低估。

报告显示，中国港口铁矿石库存约为1.6亿吨，处于历史高位，但自3月峰值以来已下降约700万吨。这一数据容易被市场误读为“去库存开始”，从而支撑铁矿石价格。

但JPM的数据暗示，库存下降的速度和幅度并不足以改变供需格局。当前库存水平仍远高于五年区间中值，且钢厂利润在焦煤价格上涨的挤压下仍在恶化。报告中的中国热轧卷板利润已重新转负，螺纹钢利润也徘徊在盈亏线附近。

这意味着，如果钢厂利润继续承压，减产将是大概率事件。而减产一旦发生，铁矿石需求将随之下降，港口去库存的

[... middle omitted ...]

化是结构性的——例如，巴西和南非的港口基础设施瓶颈长期无法缓解，或者地缘政治风险溢价永久性上升——那么全球铁矿石的边际成本曲线将被重新定价。

这个问题之所以关键，是因为它直接决定了当前对英美资源和库博铁矿的Underweight判断是否具有中长期的合理性。如果运费分化只是暂时的，那么这些被低估的资产可能反而提供了买入机会。反之，如果分化是结构性的，那么澳大利亚矿商的估值溢价应该进一步扩大。

报告没有给出答案。这需要读者结合对航运市场、港口投资和地缘政治的独立判断来补充。这恰恰是深入阅读原始报告、追踪后续数据更新的价值所在。

5. 对读者的观察框架：如何从运费数据反推铁矿石定价权的转移

基于这份报告的逻辑，投资者和产业决策者可以建立一套更精细的跟踪框架，而非仅仅盯着中国粗钢产量或铁矿石CFR价格。

第一层：关注运费价差而非绝对价格。澳大利亚与巴西/南非至中国的运费价差，比铁矿石CFR价格更能反映利润分配的边际变化。价差扩大，意味着澳大利亚矿商的相对竞争力增强。

\subsection{[R009] Bernstein：SOCAMM2不会侵蚀内存接口芯片市场，市场高估了替代风险}
{\small\color{refgray}归类：AI与算力：基础设施定价权重估与硬件ROE结构性上升}

Bernstein：SOCAMM2不会侵蚀内存接口芯片市场，市场高估了替代风险

NVIDIA最新发布的ARM架构服务器CPU Vera采用了名为SOCAMM2的新型内存封装格式。这一变化引发了一个关键问题：如果SOCAMM2成为主流，内存接口芯片的总可寻址市场将显著缩小，因为其每模块的接口芯片价值仅为数美元，远低于DDR5 MRDIMM的50至70美元。投资者对澜起科技和瑞萨电子的担忧并非没有道理。

但Bernstein这份最新研报给出了一个清晰且反直觉的判断：SOCAMM2大概率将局限于NVIDIA自身生态，不会对更广泛的服务器内存接口芯片市场构成实质性威胁。市场真正需要关注的，不是SOCAMM2本身，而是它揭示的一个更深层结构性趋势——AI服务器架构正在从“通用标准”走向“垂直整合”，而这一变化对不同参与者的含义截然不同。

这份报告的核心价值不在于技术细节的罗列，而在于它提供了一个区分“短期噪音”与“长期结构”的分析框架。对于持有或关注澜起科技、瑞萨电子的投资者而言，真正需要判断的不是SOCAMM2的技术优劣，而是它能否跨越NVIDIA的围墙，进入x86和其他ARM生态。Bernstein的答案是否定的，且给出了令人信服的论据。

1. SOCAMM2本质上是NVIDIA的全栈优化工具，而非颠覆性的行业标准

SOCAMM2并非凭空出现。它是NVIDIA从Grace到Vera架构演进中，为解决特定问题而设计的定制化方案。Grace CPU采用焊接式LPDDR5x，虽然能效出色，但无法现场升级或更换——一旦内存故障或容量需求变化，整块主板必须更换。SOCAMM2通过可拆卸的模块化设计解决了这一痛点。

然而，这种设计的优势高度依赖于NVIDIA的机架级优化逻辑。在NVIDIA的AI服务器中，GPU计算功耗是核心约束，CPU内存的能效必须为此让步。SOCAMM2在1.05V电压下比RDIMM降低约30\%的总能耗，这恰恰是NVIDIA最看重的指标——每瓦特产生的token数。

但关键问题在于，这种优化的代价是容量和带宽的明显妥协。SOCAMM2单CPU总容量为1.5TB（8模块x192GB），而MRDIMM可轻松达到16TB以上。带宽方面，SOCAMM2约1.2TB/s，介于DDR5 RDIMM的819GB/s和MRDIMM的1.6TB/s之间。对于需要处理大规模内存数据库或HPC工作负载的x86平台而言，这种取舍不可接受。

因此，SOCAMM2不是一种“更好的内存标准”，而是NVIDIA在特定约束下做出的最优解。它的存在价值，只有在NVIDIA的全栈控制框架内才能最大化。

2. x86和其他ARM生态面临不可逾越的转换成本，SOCAMM2的扩散被天然阻断

Bernstein报告中最具洞察力的部分，是指出了SOCAMM2扩散的结构性障碍。这不仅仅是技术性能的比较，更是生态锁定成本的考量。

x86服务器CPU已深度嵌入DDR生态系统数十年。从内存控制器设计、主板布局、BIOS/固件支持，到操作系统内存管理、应用层优化，整个体系围绕DDR DIMM构建。转向SOCAMM2意味着：重新设计内存控制器、重新布线主板、重新验证整个平台的兼容性、放弃与现有DDR模块的向后兼容性。这些成本对于英特尔和AMD而言是“毁灭性”的——它们没有NVIDIA那种从零开始设计ARM CPU的“绿色field”优势。

更重要的是，x86平台的核心客户——企业级数据中心和云服务商——对容量和带宽的需求远高于NVIDIA AI服务器的典型场景。MRDIMM在容量和带宽上的绝对优势，使其成为这些客户无法替代的选择。SOCAMM2在容量上的天花板（1.5TB/CPU）对于需要8TB以上内存的数据库、虚拟化、HPC场景而言，根本不够

[... middle omitted ...]

VIDIA采用，其对内存接口芯片TAM的影响也仅限于NVIDIA CPU的出货量份额——预计在未来几年内仅为高单至低双位数百分比。这意味着，SOCAMM2带来的TAM缩减，完全可以被MRDIMM的增量采用所抵消。

更值得关注的是，SOCAMM2相对于Grace的焊接式LPDDR5x，实际上增加了内存接口芯片的使用。Grace的LPDDR5x直接焊接在主板上，几乎没有接口芯片；而SOCAMM2需要SPD和电压调节器，虽然价值量不高（数美元），但对澜起科技和瑞萨电子而言，这反而是增量而非减量。

Bernstein维持对澜起科技和瑞萨电子的“跑赢大盘”评级。对于澜起科技，市场对SOCAMM2替代风险的担忧被过度放大。真正支撑其投资逻辑的，是MRDIMM带来的结构性增长——每模块50至70美元的接口芯片价值量，以及x86平台向MRDIMM迁移的确定性趋势。对于瑞萨电子，其内存接口业务规模与澜起科技相当，但仅占公司总收入的个位数百分比，市场尚未充分定价这一高增长业务的价值。

\subsection{[R010] DB：市场低估了资本账户对人民币的支撑力}
{\small\color{refgray}归类：区域市场：亚洲外汇储备分化与东南亚互联网韧性}

人民币汇率正在经历一个罕见的“结构分化”阶段。出口商结汇意愿连续三个月下降，但人民币并未因此走弱——真正支撑汇率的，是来自债券市场的资本账户流入。DB在6月16日发布的这份报告中，用一组数据揭示了这一变化：过去两个月，证券投资净流入增加了180亿美元，是2021年以来最大规模的连续流入。其中约72亿美元已在即期市场结汇，直接构成对人民币的买盘。

这个信号之所以值得关注，是因为它指向一个正在发生的结构性切换：人民币汇率的定价逻辑，正在从“贸易顺差驱动”向“资本账户驱动”转变。而市场对人民币的定价，可能仍然停留在旧框架里。

这份报告的核心价值，不在于预测人民币会涨到多少，而在于它提供了一个理解汇率的新坐标。当出口商开始囤积美元，而外资开始买入人民币债券，这两股力量同时作用，意味着什么？以下是我们从报告中提炼出的五个层次洞察。

报告引用的国家外汇管理局数据显示，5月出口企业结汇比率降至62.6\%，连续第三个月下滑，较3月的67.5\%下降了近5个百分点。直观理解，出口商在持有更多美元而不愿结汇，通常意味着对人民币前景的谨慎。

但DB的分析没有止步于此。他们注意到，尽管结汇率下降，客户端的美元净卖出量却连续三个月增加，累计约400亿美元。这说明，推动美元卖出的力量已经从出口商切换到了其他参与者。

一个合理的推论是：出口商之所以犹豫，是因为全球不确定性上升和进口成本增加，使得他们更倾向于保留美元头寸以应对潜在风险。这不是对人民币的看空，而是一种防御性操作。如果未来不确定性消退，这部分被压抑的结汇需求可能成为人民币的额外支撑。

这里的关键判断是：结汇率的下降本身，不等于人民币贬值压力的上升。它只是反映了不同市场参与者的行为分化。真正需要追踪的，是那些正在买入人民币的新增力量。

报告最引人注目的发现，来自资本账户。SAFE的净外汇收付数据显示，证券投资净流入在5月增加了180亿美元，这是自2021年以来的最大单月流入。其中约72亿美元已在银行间即期市场结汇，直接形成了对人民币的买盘压力。

这个数据点的意义在于：它标志着人民币汇率驱动力的结构性转变。过去十年，人民币的定价权主要掌握在出口商手中——贸易顺差带来的美元流入，通过结汇转化为人民币买盘。但现在，当出口商开始减速，资本账户却接过了接力棒。

更值得关注的是，这种流入并非昙花一现。中央国债登记结算有限责任公司的数据显示，5月境外投资者净买入约130亿美元中国债券，这是自2025年4月以来的首次净买入。DB指出，考虑到全球投资者对中国债券的配置仍然偏低，这些流入有望持续。

这意味着，人民币汇率的底部正在被重新锚定——不再仅仅由贸易顺差决定，而是越来越多地受到全球资产配置决策的影响。对于投资者而言，这意味着需要调整观察汇率的方法论：只看结售汇数据可能已经不够了。

为什么外资在5月突然重返中国债券市场？报告给出了一个关键线索：中国国债年初至今的强劲表现。在全球利率环境波动加剧的背景下，中国国债的相对稳定性和收益率吸引力正在被重新定价。

但这只是表层原因。更深层的逻辑在于“低配后的回补”。全球投资者对中国债券的配置比例，长期以来显著低于中国经济在全球GDP中的权重。这种配置缺口，在特定催化剂下可能引发系统性的再平衡。

报告没有展开的是，这种再平衡的持续性取决于两个条件：一是中国债券市场的流动性和可交易性是否足够支撑大规模外资进出；二是汇率预期是否稳定。从当前数据看，这两个条件都在改善。尤其是当资本账户流入开始对汇率形成正向反馈——外资买入债券，结汇推升汇率，汇率稳定又增强外资信心——这个循环一旦启动，其影响力可能超过许多人的预期。

对于关注人民币资产定价的读者来说，这意味着一个重要的观察框架：中国债券的外资持有比例，正在从一个边缘指标，变成影响汇率定价的核

[... middle omitted ...]

这意味着，投资者在参考这份报告时，需要自己补充对地缘政治情景的概率判断。这既是报告的局限性，也是它的诚实之处——它没有试图用一个精确的预测来掩盖不确定性。

对于读者而言，这个未解问题恰恰是继续深入研究的入口：如何建立一个结合地缘政治情景的汇率分析框架？哪些先行指标可以预警资本流动的转向？这些问题，在完整报告中可能有更详细的讨论。

综合以上分析，这份报告最具价值的贡献，可能不是它的结论，而是它隐含的方法论提示：传统的汇率分析框架，以贸易收支和央行干预为核心，正在被一个更复杂的多主体、多账户模型所取代。

在新的框架下，出口商、进口商、外资机构、央行、主权基金的行为需要被同时追踪。结汇率下降和资本账户流入上升，这两个看似矛盾的信号，实际上指向同一个结论：人民币定价权的转移。

对于企业财务官和资产管理者来说，这意味着过去依赖的汇率预测模型可能需要重新校准。如果资本账户的影响力持续上升，那么影响汇率的变量清单就需要加入全球利率差异、债券市场流动性、外资配置动态等新指标。

\subsection{[R011] NOM：亚洲央行正在用外汇储备为货币信心买单，但印尼和印度的弹药已接近红线}
{\small\color{refgray}归类：区域市场：亚洲外汇储备分化与东南亚互联网韧性}

NOM：亚洲央行正在用外汇储备为货币信心买单，但印尼和印度的弹药已接近红线

2026年5月的数据，揭示了一个被市场普遍忽视的结构性变化：亚洲新兴经济体的外汇储备消耗，已经从“周期性防御”转向“结构性失血”。

NOM在6月17日发布的亚洲外汇策略报告中，给出了一个值得所有关注新兴市场资产定价的投资者认真审视的结论——印度尼西亚央行和印度储备银行在5月以极为激进的力度抛售美元以支撑本币，其外汇储备消耗速度已经触及甚至突破了IMF框架下的警戒线。与此同时，中国央行却在逆向买入美元，以阻止人民币过快升值。

这不是一个关于“谁在贬值”的简单叙事。这是亚洲经济体在美元周期尾部出现的根本性分化：那些经常账户赤字、外资依赖度高、外债负担重的国家，正在用不可再生的储备资产为汇率稳定买单；而那些拥有庞大经常账户盈余和资本管制能力的国家，则拥有在美元走弱时主动管理汇率节奏的奢侈。

这份报告最核心的洞察在于：市场对亚洲货币的定价，可能低估了“储备消耗”这一供给侧的硬约束。当印尼央行的外汇储备充足率在扣除短期负债后降至46\%，当印度央行在2026年前五个月已经消耗了近9\%的2025年末外汇储备，继续干预汇市的能力和意愿都将面临严峻考验。这不仅仅是汇率波动的问题，它可能倒逼出更激进的货币政策行动——就像印尼在4月的非周期加息，以及印度针对外币存款推出优惠掉期安排所预示的那样。

1. 印尼和印度在5月的干预力度远超市场感知，储备消耗已进入“不可持续”区间

NOM通过拆解各国央行公布的储备变动、估值效应、远期头寸调整等复杂因素，估算出5月份各央行的实际外汇干预规模。结果呈现出一个清晰的“双峰”格局：印尼央行和印度央行是绝对的干预主力，而其他亚洲经济体则相对克制。

印尼央行在5月份总计卖出约48亿美元的外汇储备（含即期和远期），相当于其4月末储备规模的3.9\%。这使得其2026年以来的累计干预规模达到212亿美元，相当于2025年末储备的15.7\%。这是一个惊人的数字。更令人担忧的是，尽管付出了如此沉重的代价，印尼盾的即期汇率在5月仍然创下17,887的历史新低，并在6月8日进一步跌至18,190。干预并未能扭转趋势，只是延缓了贬值速度。

印度央行的动作更为庞大。NOM估算其5月即期外汇干预规模达到115亿美元，占4月末储备的2.1\%。2026年至今，印度央行的累计干预已高达491亿美元，相当于2025年末储备的8.9\%。卢比即期汇率在5月20日触及96.965的历史新低，之后因央行干预而略有回弹。

NOM的分析师特别指出，IMF的四项储备充足性指标中，印尼央行截至4月的平均充足率仅为88\%。但更关键的指标是“扣除预定短期流出后的储备充足率”——这一数字骤降至46\%。这意味着，如果考虑到印尼央行已经承诺的短期外汇流出（如远期合约到期、外资汇回等），其实际可动用的“自由储备”已经不足一半。在这个水平上，继续大规模干预不仅会进一步削弱市场信心，还可能触发主权信用评级机构的关注。

印度的情况虽然总量上看似更安全（平均储备充足率234\%），但其消耗速度是惊人的。从2025年6月的储备峰值至今，印度已经消耗了1008亿美元的外汇储备，占峰值的17\%。这种消耗速度如果持续，印度储备银行将面临与印尼类似的困境——不是有没有子弹的问题，而是子弹还能打多久的问题。

这里的含义非常直接：印尼和印度央行已经陷入了“干预-贬值-再干预”的负循环。每一次干预都在消耗宝贵的储备，而市场参与者对此心知肚明，反而会利用每一次反弹来建立新的空头头寸。最终，要么央行放弃干预、允许汇率一步到位式贬值，要么被迫采取更极端的资本管制或利率手段。

与印尼和印度形成鲜明对比的是中国。NOM估算，中国人民银行在5月净买入了约287亿美元的即期外汇储备。这相当于其4月末

[... middle omitted ...]

当印尼、印度、韩国等经济体在为货币贬值而焦头烂额时，中国却在为货币升值过快而“烦恼”。这背后是经常账户地位的巨大差异：中国的月均经常账户顺差高达547亿美元，而印度则是98亿美元的逆差。这种结构性差异决定了各国央行在汇率管理上的“奢侈程度”完全不同。

对于投资者而言，这意味着在亚洲外汇市场中，人民币正在扮演一个“稳定锚”的角色，而非“风险源”。当人民币对美元保持稳定甚至升值时，亚洲其他货币的贬值压力会被部分对冲，因为人民币在亚洲贸易和金融链条中具有巨大的权重。但如果人民币因中国自身经济放缓或政策转向而出现贬值压力，那么亚洲其他货币的贬值空间将被进一步打开。

3. 韩国央行干预力度有限，外汇储备的“政治经济学”正在发挥作用

韩国是一个值得单独讨论的案例。5月，韩元对美元贬值了2.3\%，贬值压力在亚洲范围内仅次于印尼盾。然而，NOM估算韩国央行5月的即期外汇干预规模仅为14亿美元，仅占其4月末储备的0.4\%。2026年至今的累计干预也仅占2025年末储备的0.9\%。

\subsection{[R012] Bernstein：中国生物医药被低估的不是创新能力，而是“全球价值转化”面临的真实阻力}
{\small\color{refgray}归类：生物医药与半导体：全球价值转化面临重估}

Bernstein：中国生物医药被低估的不是创新能力，而是“全球价值转化”面临的真实阻力

过去一个月，港股和A股生物医药板块经历了一轮显著回调。市场的主流解读是“政策恐慌”——美国《生物安全法》推进、COINS/BINSA法案提出、以及中国国内反腐行动重启，三重压力叠加。但Bernstein在最新发布的这份研报中提出了一个更值得深思的判断：市场定价反映的并非政策恐慌本身，而是对“中国创新药全球价值转化路径”的一次根本性重估。

这轮下跌不是情绪驱动的短期波动，而是投资者在追问一个此前被忽视的问题——中国药企的早期研发优势，究竟有多少能真正转化为全球商业回报？答案取决于一整套正在收紧的制度安排，而市场目前可能定价了过于悲观的极端情景。

这份研报的价值，不在于罗列政策清单，而在于它为这个问题搭建了一个分析框架：从研发管线、交易结构到情景假设，层层拆解“价值转化”这一核心变量。以下是我们从中提炼出的五个关键洞察。

1. 中国在早期研发中的角色已不可逆，但“前重后轻”的结构性矛盾才是真正瓶颈

Bernstein的数据显示，中国在全球早期药物开发中的占比已经相当可观：约30\%-40\%的临床前项目、约50\%的早期临床管线、以及2026年至今约55\%的首次临床试验注册来自中国。这组数字直观地说明，中国生物医药的“创新供给”能力已经完成了质的跃迁。

然而，这种增长并未等比转化为后期阶段的全球存在感。中国原研资产在美国三期临床试验中的占比仍低于10\%，FDA批准占比不到5\%，美国商业销售额仅约1\%。报告将这一落差归结为两个因素：一是时间滞后——近年加速的对外授权交易才刚刚开始进入美国后期管线；二是竞争力差距——美国FDA的审评标准和全球竞争强度，意味着比国内市场环境更高的淘汰率。

这里的核心含义是：中国生物医药当前面临的瓶颈不是“造不出分子”，而是“造出来的分子能否跨过全球商业化的门槛”。这个门槛由监管、临床设计、商业化能力、以及地缘政治共同决定，而最后一项正在成为最不确定的变量。

2. 美国政策不是“一刀切”，而是针对价值转化链条的逐层收紧

许多投资者将美国政策视为一个黑箱，但Bernstein将其拆解为三个层次：战略指引（美国优先投资政策）、架构设计（NSCEB报告）、以及具体立法（BINSA和生物安全法）。这种分层视角有助于理解政策的真实意图——它不是要全面封锁中国生物医药，而是要系统性地增加中国资产进入美国市场的摩擦。

具体而言，BINSA法案要求对跨境授权交易进行强制审查，限制中美合资企业和股权投资。值得注意的是，近期Innovent/辉瑞和恒瑞/BMS的合作已被美国议员公开称为“危险”交易。这表明，监管敏感性已经从CDMO（合同定制研发生产机构）扩展到了创新药授权层面。

同时，生物安全法的覆盖面也在扩大。药明康德在2026年6月被列入美国国防部1260H清单，这意味着它自动符合生物安全法的约束范围。如果OMB在年底发布正式的“关注生物技术公司”名单，中国CDMO与西方药企之间的结构性脱钩可能加速。

这些政策叠加在一起，指向一个共同的方向：中国生物医药进入美国市场的“通道”正在变窄、变贵、变不确定。

3. 中国国内政策对创新药总体友好，但“IP外流”监管正在成为新变量

在美国政策收紧的同时，中国国内政策环境对创新药总体保持支持。2026年4月发布的药品价格形成机制意见允许高创新度药物在上市时定高价；5月的NRDL调整方案引入了商业健康保险目录，为创新药提供更多支付渠道；罕见病和儿科药物的市场独占期也首次以法规形式确定。

但报告也指出一个尚未被充分定价的风险：中国政府对“技术外流”的监管正在加强。在Meta-Manus交易被叫停后，市场开始担忧生物医药领域的对外授权是否会受到类似限制。从目前情

[... middle omitted ...]

个股的影响差异显著。在完全封锁情景下，信达生物、恒瑞医药、翰森制药等全球敞口较大的公司估值可能面临约50\%的下行空间。科伦博泰相对最具韧性，因为其旗舰产品sac-TMT是主要估值锚点，且已有17项美国临床试验在进行中，风险相对已释放。

关键判断在于：Bernstein认为完全封锁属于低概率情景，更可能的路径是“中间地带”——跨境合作持续但摩擦增加。而当前的市场定价，似乎已经反映了比这个基准情景更严重的假设。换句话说，板块可能被过度惩罚了。

这份报告在情景分析上做得相当扎实，但也留下了一些开放性问题，这些恰恰是决定投资时机的关键。

第一个问题是：如果市场定价已经反映了比基准情景更悲观的假设，那么触发估值修复的催化剂是什么？是BINSA法案在年底NDAA立法中未通过？是某笔重大授权交易顺利完成？还是中国创新药在美国三期临床取得突破性数据？报告没有给出明确的催化剂时间表，但暗示了几个关键节点：2026年8月的国家安全评估、年底的BCC名单发布、以及药明康德诉讼的结果。

\subsection{[R013] UBS：中东局势缓和下，化工板块的真正机会不在油价，而在成本端的重新定价}
{\small\color{refgray}归类：能源与大宗：和平交易过度定价，供给约束支撑三季度油价}

UBS：中东局势缓和下，化工板块的真正机会不在油价，而在成本端的重新定价

市场当前正在消化一个关键信号：美国与伊朗就霍尔木兹海峡重新开放达成协议。这一事件的影响远不止于油价本身。UBS最新发布的每周跟踪报告给出了一个值得产业决策者反复推敲的判断——如果局势快速缓和，布伦特原油价格在2026年第三季度平均可能回落至85美元/桶。但这份报告最有价值的洞察，并非油价预测，而是它揭示了化工板块内部正在发生的结构性分化：哪些公司会在成本端受益，哪些公司反而可能面临短期压力。

对于化工行业投资者而言，当前需要回答的核心问题不是“油价会跌多少”，而是“当油价回归常态，哪些企业的盈利弹性被市场低估了”。UBS的报告为此提供了清晰的观察框架。

UBS在报告中明确列出了受益于地缘政治缓和的化工企业类型。第一类是受原材料价格影响较大的公司，包括涤纶长丝（桐昆）和磷化工企业。这背后的逻辑并不复杂：这些企业的成本结构中，原油或原油衍生物占比较高，油价回落直接意味着成本下降。

但值得注意的，并非所有成本敏感型企业都会同等受益。报告给出了一组重要的运营数据作为参照：上周中国PTA开工率环比上升2个百分点至65\%，而涤纶长丝开工率环比下降4个百分点至76\%。这意味着下游需求端并未同步回暖。如果仅仅因为油价下跌就买入化工股，可能会忽略需求端的疲软。真正的受益者，必须是那些成本下降速度快于产品价格下降速度的企业。UBS的报告实际上在提醒：成本端受益，不等于利润端受益，中间还有一条“需求传导链条”需要验证。

UBS特别提到了两类具备重新评级潜力的化工龙头：万华化学和扬农化工。这一判断值得深入拆解。

万华化学的MDI开工率上周环比大幅上升9个百分点至83\%，这是报告中最引人注目的产能数据之一。在整体化工开工率普遍低迷的背景下——乙烯开工率环比下降2个百分点、PP开工率同比下降14个百分点——MDI的逆势提升显示出万华在聚氨酯领域的定价权和需求韧性。UBS将万华列为“re-rating potential”标的，其隐含逻辑是：当油价回落、成本端压力缓解时，万华这类拥有技术壁垒和规模优势的企业，其盈利改善幅度可能远超行业平均水平。

扬农化工的情况类似。作为农药龙头，其成本结构中原油衍生品占比较高，但产品定价更多取决于全球农产品周期和自身产品结构。UBS将其列为受益者，暗示的是当外部成本冲击消退后，市场可能重新聚焦其内生增长能力。

这里需要补充一点：UBS的报告没有完全展开的是，万华和扬农的re-rating逻辑是否成立，还取决于它们能否在行业景气低谷期维持产能利用率不降。MDI开工率的环比提升是一个积极信号，但需要更多季度的数据来确认趋势。

3. 替代路线生产商面临的短期压力，是报告里最容易被忽略的风险提示

UBS在报告中指出，拥有替代生产路线的烯烃生产商（宝丰能源和卫星化学）可能在短期内面临压力。这一判断初看反直觉：油价下跌，使用煤制烯烃（CTO）或乙烷裂解路线的企业，成本优势不是应该扩大吗？

但UBS的逻辑恰恰相反。当油价下跌时，石脑油裂解路线的成本随之下降，这会压缩煤制烯烃和乙烷裂解路线原本享有的成本溢价。换句话说，油价下跌反而削弱了这些替代路线的竞争优势。报告数据显示，上周石脑油基乙烯开工率环比下降2个百分点至77\%，而MTO路线开工率持平于83\%。这一数据暗示，替代路线的开工率并未因油价预期而提升，市场已经在提前定价这种竞争格局的变化。

对于持有宝丰能源和卫星化学的投资者来说，这是一个需要认真对待的短期风险。UBS的表述是“短期面临压力”，但报告没有给出压力持续时间的判断。这取决于油价回落的幅度和速度，以及这些企业能否通过产品结构差异化来对冲成本优势的收窄。

4. 库存数据揭示的真相：去库存周期尚未结束，但结构性差异正

[... middle omitted ...]

环比27\%的增幅可能意味着下游需求不足以支撑当前开工率水平。钛白粉开工率上周持平于79\%，如果库存继续累积，开工率面临下行压力。

涤纶长丝库存环比下降16\%，但同比上升66\%。这一数据组合暗示：当前库存下降更多是季节性因素或短期订单驱动，而非需求趋势性回暖。UBS将涤纶长丝企业列为成本端受益者，但如果需求端不能同步改善，成本下降带来的利润弹性可能被高估。

5. 报告尚未完全回答的关键问题：AI驱动的化工材料需求，能否成为新的增长极？

UBS在报告末尾提到了几个值得关注的选股方向：杰瑞股份作为AI发展的受益者，以及部分MLCC企业和氟化工材料企业。这一部分内容在报告中着墨不多，但可能是未来最值得跟踪的变量。

报告提到UBS本周将与可能受益于AI发展的化工材料公司举行电话会议。这说明该机构正在积极评估AI对化工材料需求的影响。从逻辑上推演，AI基础设施建设（数据中心、服务器、冷却系统）对高端化工材料的需求正在上升，包括电子特气、高性能氟材料、MLCC用陶瓷粉体等。

\subsection{[R014] GS：市场真正低估的不是需求崩塌，而是供给恢复的路径与节奏}
{\small\color{refgray}归类：能源与大宗：和平交易过度定价，供给约束支撑三季度油价 / 航运与供应链：供给约束驱动定价逻辑重置}

布伦特原油现货价格本周跌破80美元/桶。触发因素是一次传闻中的美伊临时协议——如果协议落地，霍尔木兹海峡的封锁可能迅速解除，超过5000万桶伊朗海上浮仓原油将立即进入市场。GS在最新一期《石油追踪》报告中，用一个数字重新定义了市场正在定价的叙事：70\%的战前霍尔木兹流量，可能成为新的“100\%”。

这个数字背后，不是简单的供需再平衡，而是一套关于供给恢复速度、库存真实消耗、以及需求结构性错位的全新框架。市场目前的下跌，更多是在定价“协议达成”这一事件本身，但GS的测算表明，真正决定中期价格中枢的，是协议达成后供给恢复的路径、节奏，以及隐藏在库存数据之下的结构性变化。

这份报告的价值不在于给出一个方向性判断，而在于拆解了三个市场尚未充分定价的关键变量：霍尔木兹流量的实际恢复斜率、中国库存数据的“暗物质”缺口、以及IEA与GS在供需平衡表上的系统性分歧。理解这三个变量，比猜测油价下一个整数关口重要得多。

1. 霍尔木兹流量恢复至战前70\%，可能就是一个新的供需平衡点

GS的核心测算框架建立在一个反直觉的假设上：波斯湾出口恢复至战前水平，可能不需要霍尔木兹海峡流量回到100\%。报告指出，通过现有绕行路线——包括沙特的东西管道（Yanbu）、阿联酋的阿布扎比原油管道（Fujairah）、以及土耳其的基尔库克-杰伊汉管道——已经形成了每日750万桶的替代流量。加上阿曼湾的“暗渡”流量，当前实际绕过霍尔木兹的总流量约为每日910万桶。

这意味着，要让波斯湾总出口恢复到战前水平，霍尔木兹本身只需要额外增加每日1270万桶的流量，即恢复到战前约70\%的水平。这个“70\%即100\%”的判断，直接降低了供给恢复的门槛。市场此前隐含的假设是，霍尔木兹必须完全恢复正常，供给压力才能缓解。GS的数据表明，即使海峡流量只恢复到七成，配合现有绕行能力，全球供应就能回到战前格局。

这个结论对资产定价的含义是：即便谈判出现波折，只要现有绕行路线不中断，供给冲击的烈度已经在边际上减弱。市场对“封锁解除”这一事件的定价，可能高估了其实际增量效果。

供给恢复的另一个潜在约束是运输能力。GS测算了霍尔木兹海峡两侧5天航程内的空载油轮运力，结果显示有8.61亿桶的装载能力。从物理运力角度看，不存在瓶颈。

但报告也直言不讳地指出，许多船东对过境仍持谨慎态度，缺乏明确的过境指引是主要障碍。这意味着，即便政治协议达成，从“协议签署”到“船东愿意正常航行”，中间可能还有一段风险溢价消散的时间。船东的风险偏好，而非运力本身，将成为供给恢复速度的边际约束。

此外，伊朗在未来60天核谈判中的地缘政治目标，也可能影响其配合协议执行的意愿。供给恢复的路径不是一条直线，而是一个包含政策落地、商业信心重建、地缘博弈的多阶段过程。市场对“协议即通航”的线性外推，可能低估了中间环节的不确定性。

3. IEA与GS的供需缺口分歧，折射出对需求破坏程度的不同判断

报告中最值得关注的交叉验证，来自GS与国际能源署（IEA）在供需平衡表上的分歧。IEA最新《石油市场报告》估计二季度全球石油赤字为每日310万桶，而GS的估计是每日500万桶。差距的核心，在于对需求破坏规模的判断不同。

IEA认为二季度全球需求破坏高达每日550万桶，其中5月达到峰值每日610万桶。GS的对应数字是每日490万桶。这个约每日60万桶的差异，看似不大，但在一个全球供需本就紧平衡的市场中，足以改变库存变化的斜率。

更关键的是需求破坏的结构。IEA的数据显示，需求损失高度集中在石化原料领域：液化石油气和乙烷需求较战前预期下调每日170万桶（降幅11\%），石脑油需求下调每日110万桶（降幅15\%）。这与GS的分析一致——当前需求破坏并非全面性的经济衰退，而是特定产业链环节的调整。这意味着，一

[... middle omitted ...]

据）仅显示每日30万桶的去库。

两者之间每日130万桶的缺口，报告提出了两种可能解释：一是存在大量不可见的战略储备释放（如地下储备）；二是实际炼厂开工率低于官方数据，即真实需求比统计数字更弱。GS倾向于后一种解释，因为这与5月中国汽油及其他零售油品销量同比下降23\%的数据相吻合。

这个“缺口”对全球市场的含义是：中国作为全球最大的原油进口国，其真实需求可能比官方数据暗示的更低。如果不可见去库持续，意味着中国在供给恢复后的补库需求可能弱于预期。反之，如果缺口来自统计口径问题，那么一旦需求真实回暖，补库力度可能超出市场想象。无论哪种情况，这个每日130万桶的“暗物质”都是当前全球石油平衡表上最大的不确定性来源。

5. 报告尚未完全回答：IEA对2027年巨大过剩的预测，是否过于激进？

IEA在最新报告中首次给出了2027年的供需平衡预测，预计全球将出现每日500万桶的过剩，而GS的预测仅为每日320万桶。分歧主要来自IEA假设2027年供给增长高达每日790万桶。

\subsection{[R015] JPM：中国楼市真正的分水岭不在价格，而在供给侧的“窄复苏”能否撑到政策显效}
{\small\color{refgray}归类：地产与消费：K型分化加剧，一线回暖难掩全国压力}

JPM：中国楼市真正的分水岭不在价格，而在供给侧的“窄复苏”能否撑到政策显效

市场正在被两个相互矛盾的现象拉扯。一方面，5月JPM住房活动指数小幅回升，一线城市二手房价格连续录得环比正增长，部分核心区甚至出现“抢房”传闻。另一方面，全国新房销售面积同比跌幅从-6.5\%扩大至-8.6\%，房地产开发投资跌幅进一步加深至-24.3\%，低线城市量价齐跌的格局毫无松动迹象。

这份由JPM中国房地产研究团队发布的月度跟踪报告，其核心判断并不复杂，但含义却比多数市场解读更为深远：当前中国楼市正在经历的，不是一轮周期性复苏的前夜，而是一场“窄基K型”的结构性再平衡。一线城市的回暖并非全国复苏的领先指标，而是政策资源、人口流动和存量资产重新定价的结果。真正决定行业走向的，不是房价何时见底，而是“十五五”城市更新规划能否在供给端建立起一个不同于过去二十年“卖地-开发-销售”循环的新均衡机制。

报告用一个词概括了这种状态——“量价调整的长期化”。这不是悲观，而是一种必要的清醒。

1. 一线城市二手房回暖的真实驱动不是需求爆发，而是库存出清和供给结构切换

市场最容易犯的错误，是把一线城市的局部量价改善直接等同于全国需求的底部确认。JPM的报告提供了两个关键数据，足以让这种乐观情绪降温。

第一，5月70城新房价格环比跌幅稳定在-0.20\%，但一线城市录得+0.2\%的环比正增长。更值得关注的是二手房市场：全国二手房价格环比跌幅从-0.23\%小幅扩大至-0.26\%，而一线城市却维持了+0.4\%的环比涨幅。这意味着，一线城市和低线城市之间的差距不仅没有收敛，反而在拉大。

第二，报告指出，一线城市回暖的支撑因素并非居民购房意愿的全面回升，而是三个更具体的结构性条件：二手房市场库存较低、成交量改善带动换手率上升、以及部分业主开始试探性提价。此外，报告还含蓄地提到了“IPO活动和股市财富效应”可能对高端需求产生的边际影响——这是一个容易被忽视的变量，但它的可持续性高度依赖资本市场表现。

更深层的含义在于，一线城市二手房正在加速替代新房成为交易主力。报告明确提到“二手房市场正在从新房销售中不断夺取份额”。这意味着，即便整体交易量企稳，开发商层面的新屋销售、施工和投资仍会持续承压。对于投资者而言，一线城市二手房回暖更多是存量资产流动性的修复，而非开发商基本面的拐点信号。

2. 低线城市的困境不是周期性滞后，而是库存结构和信用传导机制的断裂

与一线城市的窄幅改善形成鲜明对比，低线城市面临的问题已经超出了简单的“需求不足”。JPM的图表清晰地展示了两条令人不安的曲线：在建库存月数从2021年的35个月一路攀升至2026年的75个月，而竣工库存月数也从3个月翻倍至7个月。这是历史级别的库存压力。

更关键的是，低线城市的问题不是“卖不掉”，而是“卖不掉之后”的连锁反应。报告指出，居民信贷需求疲弱、房地产开发投资持续下降，形成了一个负向循环：库存高企导致开发商减少新开工，新开工下降拖累土地出让收入，土地收入减少又削弱地方政府的财政能力和城市更新投入，最终进一步压低居民购房意愿和资产价格预期。

这种传导机制使得低线城市的调整具有极强的惯性。即便政策层面出台再多的“稳市场”措施，只要居民对未来收入和资产价格的预期没有根本性扭转，需求就很难被政策刺激有效激活。报告将其称为“时间不一致性陷阱”——政策制定者希望降低经济对房地产的依赖、改善居民可负担性，但又不敢让价格过快出清，以免冲击居民财富和地方财政，最终只能通过量的收缩来消化压力。

这意味着，低线城市的调整可能会比多数人预期的更为漫长。对于布局在这些市场的开发商、金融机构和地方政府而言，真正的挑战不是“何时见底”，而是如何在不引发系统性风险的前提下完成这一轮出清。

3. “十五五”城市更新

[... middle omitted ...]

著扩容：财政预算投资、保障房专项债、超长期特别国债、税收优惠、REITs和社会资本都被纳入支持框架；土地政策方面，闲置土地再利用、混合用途开发、临时使用、用途转换和存量资产盘活的交易摩擦降低，都得到了明确支持。

这些变化的本质，是政策从“卖地-开发-销售”的增量循环，转向“存量盘活-资产运营-公共服务”的可持续模式。过去二十年，城市更新的核心驱动力是土地增值预期；未来，城市更新的核心驱动力将是财政可持续性和资产运营效率。

但报告也清醒地指出，这一规划的定位是“稳定而非复苏房地产市场”。换句话说，政策的目标不是让房价重新上涨，而是让市场在更低的量价水平上找到新的均衡。对于投资者而言，这意味着需要重新定义“利好”——不是看政策能否刺激需求，而是看政策能否加速供给出清、降低系统性风险溢价。

4. 报告没有完全回答的关键问题：城市更新的资金闭环能否跑通？

尽管“十五五”规划在政策工具上做出了显著升级，但报告并未详细展开一个核心问题：这些资金能否在实际操作层面形成闭环？

\subsection{[R016] GS：欧洲数据中心正在重塑电力行业的估值逻辑，而市场仍按旧地图定价}
{\small\color{refgray}归类：AI与算力：基础设施定价权重估与硬件ROE结构性上升}

GS：欧洲数据中心正在重塑电力行业的估值逻辑，而市场仍按旧地图定价

欧洲的电网运营商正在收到一个前所未有的信号。截至2026年5月，GS追踪的欧洲数据中心并网申请总量已飙升至约480吉瓦，相当于整个欧洲当前电力需求的1.5倍。六个月前这个数字是290吉瓦，十二个月前是170吉瓦。

这不是一个线性增长，而是一个指数级跳升。即便考虑到其中一部分申请可能带有投机性质，这个规模本身已经构成一个判断：欧洲正在进入一个数据中心大规模建设的物理阶段，而电力公用事业公司将成为这个阶段最核心的枢纽角色。

GS在最新发布的这份研报中，给出了一个清晰的逻辑链条：数据中心建设拉动电力需求，电力需求拉动电网和发电投资，投资拉动公用事业公司的盈利增长。但真正值得关注的，不是这个链条本身，而是这个链条如何改变市场对这些公司的定价方式。

当前GS筛选出的三大受益集群——变革性电气化故事、可再生能源、能源安全提供商——在2030年的平均市盈率仅为12倍，低于行业平均的12.9倍。这些公司拥有更高的增长前景，却以更低的估值在交易。这组数字指向一个核心问题：市场是否在用过去的框架，定价一个正在被数据中心重塑的未来？

1. 480吉瓦的并网申请意味着电力需求增长可能比市场预期的快一倍

GS的数据追踪方法比较直接：汇总欧洲主要电网运营商收到的数据中心并网请求。480吉瓦这个数字最惊人的地方不是它的大小，而是它的增速。从2025年1月的170吉瓦到2026年5月的480吉瓦，一年半时间增长了近两倍。

这个数字意味着什么？欧洲电力消费在过去十年基本停滞，年均增速在零附近。而数据中心一个单一品类，就可能在2028至2029年间推动欧洲电力需求年增长1.5个百分点。这还没有计入电气化、氢能、制造业回流等其他结构性因素。

GS还模拟了一个“超电气化”情景：如果欧洲同时实现其2030年电气化目标并加速数据中心建设，电力需求年增速可能达到5\%。5\%的年增速对于欧洲这样一个成熟电力市场而言，是一个代际级别的变化。

这个判断的关键含义是：电力公用事业公司的收入增长假设需要被重新校准。过去十年，市场对这些公司的定价基于一个“低增长、高股息”的稳定预期。如果电力需求进入一个结构性加速期，这些公司的资本开支周期、盈利增速和估值中枢都需要被重新审视。

2. 欧洲数据中心建设有较高的可见度，但480吉瓦中的大部分仍需验证

第一个层次是欧洲数据中心协会的预测：到2031年新增20吉瓦，使总容量从当前的15吉瓦提升至35吉瓦。这20吉瓦中，约60\%已经处于在建或已获得最终投资决定的状态，其余也基本完成了审批。这一部分的可预见性很高，对应的电力需求增量大约在每年1.5个百分点。

第二个层次是GS自己的预测：到2035年，欧洲数据中心IT容量将达到65至80吉瓦。这个预测的核心变量是人工智能代理的普及速度。GS引用行业文献指出，一次人工智能代理查询的能耗可能是传统人工智能聊天机器人的13到50倍。即便GS采用了相对保守的假设——将每次代理查询的能耗从15瓦时下调至2瓦时——每10个百分点的代理查询渗透率提升，仍将推动数据中心需求增加25\%至30\%。

这里有一个尚未完全回答的问题：480吉瓦的并网申请中，有多少最终会转化为实际建设？GS坦诚地指出部分申请可能具有投机性质，但报告没有给出具体的转化率假设。这个转化率将直接影响电力需求曲线的陡峭程度，也是投资者在阅读完整报告时需要重点关注的假设变量。

3. 受益公司的估值折价背后，是市场对“盈利质量”的定价分歧

变革性电气化故事（Naturgy、Enel、Engie）的2030年平均市盈率为11倍，可再生能源集群（RWE、EDPR、Solaria、Orsted、Vestas、Nordex）为12.5倍，能源安全提供

[... middle omitted ...]

的电网投资和发电投资，而这些投资往往伴随着监管审批、建设周期和融资成本等风险。市场可能认为，即便需求旺盛，公用事业公司的盈利增长也会被资本开支扩张和监管回报率限制所抵消。

GS对此的判断是：当前的估值折价提供了一个有吸引力的入场窗口。但这里仍需验证的一个关键问题是，哪些公司真正具备将规模转化为议价权的能力。

GS的研报在逻辑上完整，但有一个重要的现实约束没有被充分展开：欧洲电网的实际承载能力。

480吉瓦的并网申请是一个理论数字，但欧洲的电网基础设施能否在2035年前消化这么多新增容量，是一个需要独立评估的问题。电网扩建的审批周期通常以年为单位，变压器和开关设备等关键设备的供应链也面临瓶颈。如果电网建设跟不上数据中心建设的节奏，那么一部分并网申请可能只是“纸上容量”。

另一个未被充分讨论的问题是电力定价机制的变化。数据中心的大规模接入将改变电力负荷曲线的形状，可能导致高峰时段的电价波动加剧。这对依赖固定电价或长期PPA的可再生能源开发商而言，既是风险也是机会。

\subsection{[R017] 摩根斯坦利：市场低估了三季度销售二次探底的冲击力}
{\small\color{refgray}归类：地产与消费：K型分化加剧，一线回暖难掩全国压力}

5月的房地产销售数据，打破了市场此前两个月的微弱乐观。摩根斯坦利最新发布的这份研报，用一组清晰的数字给出了一个明确的判断：5月销售跌幅重新扩大，而更值得警惕的是，三季度可能迎来年内第二轮深度下滑。这不是一次简单的数据波动，而是行业结构性压力的再次确认。

这份报告的价值不在于告诉你“市场不好”——这已经是共识。它的核心洞察在于：市场正在低估两个关键变量——二手房销售从放缓转向负增长的传导效应，以及开发商在土地市场上近乎停摆的谨慎姿态。这两个因素叠加，正在重塑下半年的行业定价逻辑。

为什么现在必须认真对待这份判断？因为5月的数据并非孤立事件。全国销售金额同比降幅从4月的-7.6\%扩大到-9.3\%，销售面积降幅从-9.5\%扩大到-13.2\%。更关键的是，二手房市场的实时高频数据在5月中旬开始走弱，且近几周减速明显加快。摩根斯坦利据此判断，三季度二手房销售将转为同比负增长，这将直接拖累新房价格和销售。

这意味着什么？意味着此前市场期待的“政策底”到“市场底”的传导路径，可能比预期更长、更曲折。行业的出清过程，正在从“量”的收缩进入“价”的压力阶段。

4月的数据曾让部分投资者燃起希望。销售金额和面积的同比降幅双双收窄，市场开始讨论“触底”的可能性。但5月的数据无情地打破了这种预期。

摩根斯坦利在报告中明确写道，5月销售跌幅的重新扩大“符合预期”。这个措辞本身就值得深思——这意味着该机构的预测模型早已将这种反复纳入考量。全国销售金额同比降幅从-7.6\%扩大到-9.3\%，销售面积从-9.5\%扩大到-13.2\%。累计来看，前5个月销售金额同比下降13.5\%，销售面积下降10.8\%。

更值得关注的是结构性的分化。国家统计局70城房价指数显示，5月新房价格环比下降0.2\%，二手房下降0.3\%，降幅较4月略有扩大。但一线城市却呈现出相反的走势：新房环比上涨0.2\%，二手房上涨0.4\%。这种K型分化——一线城市微弱回暖，二线及以下城市持续承压——正在成为本轮下行周期最显著的特征之一。

这里的含义很明确：市场不能再用“整体数据”来掩盖结构性风险。一线城市的韧性并不能代表全行业的健康状况。对于大多数开发商而言，核心战场在二线及以下城市，而这些区域的价格和销售压力仍在加剧。

如果说销售数据还受到政策刺激和季节性因素的干扰，那么施工活动则更真实地反映了开发商对未来的判断。因为施工是开发商用真金白银做出的决策，而非市场情绪的短期波动。

5月的数据令人担忧。新开工面积同比下降25\%，竣工面积下降20\%。前5个月累计，新开工和竣工面积均下降约23\%。房地产投资同比降幅从4月的-13.7\%进一步扩大到5月的-16.2\%。

这些数字的背后是一个清晰的逻辑链条：销售疲软导致回款困难，回款困难导致开发商无力或不愿新增投资，投资收缩又进一步恶化市场信心，形成负向循环。而打破这个循环的关键变量——土地市场——正在发出最刺耳的警报。

摩根斯坦利的数据显示，前5个月百强开发商的土地购置金额同比下降超过40\%。这不是一个边际变化，而是近乎断崖式的收缩。开发商在土地市场上的集体撤退，意味着未来一到两年的可售资源将进一步减少。这不仅影响短期的销售规模，更意味着行业供给侧的深度收缩将持续更长时间。

对于投资者而言，这意味着什么？意味着即使需求端出现边际改善，供给端的制约也会限制销售反弹的幅度和持续性。行业正在经历的不是一个短周期的调整，而是供给侧的重新定价。

摩根斯坦利这份报告中最具前瞻性的判断，是对三季度的展望。该机构预计，三季度二手房销售将转为同比负增长，新房销售降幅将在低基数基础上进一步扩大。

这个判断的底层逻辑值得仔细拆解。二手房市场在过去几个月一直扮演着“先行指标”的角色。当新房市场因限价、交付风险等因素而扭曲时，二手房价格和成交量往

[... middle omitted ...]

。而降价又会进一步侵蚀本就微薄的利润空间，并加剧购房者的观望情绪。

报告提到，三季度可能呈现“K型”房价走势：整体环比下跌，但部分一线城市因去库存进展较好，可能出现温和上涨。这种分化将更加考验投资者的选股能力——不是所有开发商都能在这一轮分化中幸存，更不用说受益。

尽管摩根斯坦利对三季度的判断已经相当清晰，但报告中有几个关键变量并未完全展开。这正是值得深入探讨的地方。

第一个问题是二手房销售转为负增长的速度和幅度。报告指出“近期减速明显加快”，但没有给出具体的量化预测。如果二手房销售在三季度初就出现两位数的同比下滑，其对新房市场的冲击将远超预期。反之，如果下滑幅度较为温和，市场可能仍有时间调整预期。

第二个问题是政策应对的空间和时点。报告没有详细讨论如果三季度数据恶化，政策层面可能采取哪些措施。从历史经验看，当二手房市场出现系统性走弱时，放松限购、降低房贷利率、加大政府收储力度都是可能的选项。但这些政策能否有效阻断负向循环，取决于执行力度和市场信心的恢复程度。

\subsection{[R018] Bernstein：亚洲硬件正在经历一轮由AI驱动的ROE结构性重定价}
{\small\color{refgray}归类：AI与算力：基础设施定价权重估与硬件ROE结构性上升 / 生物医药与半导体：全球价值转化面临重估}

Bernstein：亚洲硬件正在经历一轮由AI驱动的ROE结构性重定价

这份Bernstein研报的核心判断并不复杂，但它的支撑框架值得产业决策者和投资者仔细推敲：亚洲科技硬件公司正在经历一轮由AI需求驱动的ROE（净资产收益率）结构性上升周期，而市场对这一变化的定价，可能才刚刚开始。报告通过长达15年的资产负债表与现金流分析，揭示了不同细分赛道在盈利能力、资本效率和财务稳健性上的根本差异——这些差异在AI浪潮下正在被急剧放大。

报告覆盖的七家公司中，Chroma ATE、Delta Electronics、Unimicron等AI直接受益标的的ROE预计将在 年达到30\%甚至40\%以上，而消费电子相关公司（Luxshare、Largan、Sunny Optical）的改善则温和得多，仅能恢复到中双位数至低20\%的水平。这不仅是增速差异，更是商业模式护城河的重新定价。

以下是我们从这份研报中提炼出的五个结构性洞察，以及一个报告尚未完全回答的关键问题。

1. 这轮ROE上升不是周期性的，而是由AI需求驱动的商业模式质变

报告最值得关注的一点，是它区分了“周期性反弹”和“结构性跃迁”。过去15年，亚洲硬件公司的ROE波动大多跟随全球PCB市场周期或消费电子换机周期。但当前这一轮上升，核心驱动力来自AI数据中心对测试设备、电源管理、散热方案和高端PCB的持续需求，而非库存回补或终端需求脉冲。

以Chroma ATE为例，其ROE从2024年的25\%左右跃升至2026年预计的45\%以上，背后是AI芯片测试、光收发器测试、以及AI数据中心电源测试三大业务的同时爆发。报告预计其 年营收和营业利润的复合年增长率分别达到46\%和65\%——这不是传统设备供应商能实现的增速，而是站在AI资本开支浪潮的浪尖上。

Delta Electronics同样如此。其ROE从2025年的24\%左右预计上升至 年的40\%以上，核心驱动是AI数据中心相关收入已占其总营收的约40\%。这意味着，对于这些公司而言，AI不再是“增量业务”，而是正在成为“主干业务”。

对投资者的含义：传统上，市场对硬件公司的估值框架往往基于周期市盈率。但如果ROE的上升是结构性的，那么合理的估值中枢也应该上移。报告对Chroma给出了38倍市盈率的目标价，对Delta则没有明确给出倍数，但其隐含估值逻辑显然已经脱离了传统硬件公司的定价范式。

2. 杜邦分析揭示：高ROE的路径不止一条，但AI正在重塑每一条路径的权重

报告通过杜邦分析拆解了各公司的ROE驱动因素，结果极具启发性。

Chroma代表了“高利润率驱动”的模式：其近三年平均净利率约30\%，远超同行。这得益于其定制化产品聚焦于利基市场，技术壁垒高、竞争强度低。它的资产周转率（58\%）和杠杆率（147\%）均处于中等水平，但净利率的绝对优势使其ROE达到27\%。

Luxshare和Quanta则代表了另一条路：“高周转+高杠杆驱动”。作为ODM厂商，它们的净利率极低（Luxshare约3\%，Quanta约2\%），但资产周转率（Luxshare约140\%，Quanta约200\%）和权益乘数（Luxshare约300\%，Quanta约400\%）显著高于其他公司，从而实现了同样甚至更高的ROE。

Largan是一个有趣的对照案例。它的净利率高达约39\%，与Chroma相当，但其资产周转率极低（约27\%），且权益乘数也较低（约117\%），导致ROE仅为12\%左右。报告指出，这主要是因为Largan持有大量现金，且其高端镜头生产需要较高的资本投入（尤其是满足苹果的严格要求），导致资产效率偏低。

对投资者的含义：在AI需求爆发的背景下，不同ROE路径的可持续性正在被重新评估。高净利率模式（Chroma、

[... middle omitted ...]

例外。随着AI服务器业务的扩张，其杠杆率正在上升。报告给出的评级是“Underperform”，目标价250新台币，而当前股价为374新台币，隐含约33\%的下行空间。

这一判断值得认真对待。Quanta是全球最大的PC ODM厂商，AI服务器组装业务正在成为其第二增长引擎。但问题在于，AI服务器业务的商业模式是“买进卖出”（buy-sell），即公司需要先垫付大量资金采购关键零部件（如GPU、内存等），再组装后交付给客户。这导致其报告营收大幅膨胀，但同时也带来了营运资金压力和杠杆率上升。报告指出，Quanta的营收波动性是覆盖公司中最高的，这正是买进卖出模式的直接后果。

对投资者的含义：在AI硬件投资热潮中，营收增速和订单能见度往往是最受关注的指标。但资产负债表的变化同样重要。Quanta的案例提醒我们，当一家公司的业务模式从“轻资产代工”向“重资产采购+组装”转变时，其财务风险特征也会随之改变。投资者需要关注现金流、杠杆率和营运资本效率，而不仅仅是营收和利润表。

\subsection{[R019] Citi：市场低估的不是需求疲软，而是供给端策略的结构性拐点}
{\small\color{refgray}归类：汽车行业：补贴退潮暴露结构性压力，中国车企加速全球渗透}

Citi：市场低估的不是需求疲软，而是供给端策略的结构性拐点

这份来自Citi的周度订单追踪报告，表面上是在讨论6月第二周电动车订单环比下降13\%的短期波动。但如果只读到“需求不及预期”就停下，就错过了这份报告真正有价值的判断。

真正值得关注的，不是订单数字本身，而是Citi在报告后半段提出的一个推演：如果中国政府满足于出口增长来对冲内需疲弱，不再加码消费刺激，那么车企将被迫放弃过去几年的激进策略——这意味着市场格局、价格战烈度、成本控制逻辑，都将进入一个与过去三年完全不同的阶段。

订单波动是表象，供给端策略的系统性调整才是内核。而市场对这个内核的定价，目前仍不充分。

1. 周度订单的“不及预期”本身不是新闻，但结构信号值得拆解

Citi数据显示，6月第二周（6月8日至14日）行业整体订单环比下降13\%，半月累计环比增长仅4\%，远低于历史同期10\%的月环比增速。单独看这个数字，很容易得出“需求走弱”的简单结论。

但Citi的观察框架不止于此。报告指出，4月燃油车销量崩盘之后，新能源车需求其实并未显著受益于燃油车用户转化，而是维持了自身相对平坦的节奏。换句话说，新能源车需求并没有变差，它只是没有变好——而市场此前隐含的预期是“应该变好”。

更值得关注的是品牌间的分化。半月环比来看，小鹏（+52\%）、蔚来（+52\%）、吉利银河（+45\%）、比亚迪（+11\%）跑赢行业；但周度环比来看，只有比亚迪（-4\%）和吉利银河（+16\%）保持了相对稳定，其余品牌周度波动显著放大。

这种分化意味着什么？订单的稳定性正在成为比订单绝对值更重要的竞争维度。比亚迪和吉利银河的订单波动最小，说明它们的用户心智占据和渠道效率已经形成了某种“抗干扰”能力。而蔚来、小鹏尽管半月环比亮眼，但周度环比大幅下滑，暗示订单更多来自促销或新车型脉冲，而非稳定的自然流入。

2. Citi真正的核心判断：政府刺激缺位将迫使车企放弃“烧钱换份额”模式

这份报告最有价值的洞察，出现在Citi对下半年政策情景的推演中。

Citi提出一个关键假设：如果中国政府认为强劲的整车出口增长足以弥补国内需求疲弱，从而不再推出大规模消费刺激，那么车企将不得不调整过去几年“不惜代价抢份额”的激进策略。具体而言，Citi预期将出现三个变化：

第一，市场整合将变得温和而渐进，而非剧烈出清。这意味着过去那种“今年不赢就出局”的叙事需要修正——不是没有整合，而是整合的速度和烈度都会下降。

第二，价格战策略将趋于软化。不是不打价格战，而是降价幅度和频率都会收敛。车企会意识到，在缺乏政策催化的环境下，降价对需求的拉动边际递减，反而加速利润损耗。

第三，成本控制将变得前所未有的严格。2026下半年到2027年，车企的核心KPI将从“份额增长”转向“单位经济性改善”。

这三个变化合在一起，指向一个判断：过去三年“烧钱换份额”的竞争模式，正在接近尾声。而市场对“竞争格局改善”的定价，目前仍停留在预期阶段，尚未反映到估值中。

Citi提供的订单数据表格中，有一个容易被忽视的信号：品牌间的订单波动率正在系统性分化。

以比亚迪为例，5月至6月七周订单量分别为51.2k、43.7k、40.7k、42.3k、48.7k、47.7k、45.7k，周度波动幅度在15\%以内。吉利银河的周度波动也相对可控。相比之下，蔚来、小鹏、华为鸿蒙等品牌的周度波动幅度可达30\%-50\%。

订单稳定性背后反映的是三重能力：品牌认知的牢固程度、渠道网络的覆盖深度、以及产品矩阵的宽度。当一个品牌能够在不依赖单款爆品或短期促销的情况下，持续吸引稳定的周度订单流入，说明它已经建立起了某种“需求基础设施”。

这个指标的产业含义是：在行业从“增量竞赛”转向“存量博弈”的过程中，订单稳定性高的企业将拥有更

[... middle omitted ...]

是否可持续？如果主要出口市场（如欧洲、东南亚）出现贸易壁垒或需求放缓，出口的支撑力就会减弱。

第二，出口的利润结构如何？如果出口以低价车型为主，即使量增长，对车企利润表的贡献也有限。

第三，政府是否会因为就业压力或消费数据进一步走弱而改变态度？Citi的“满足于出口”假设是一种情景推演，而非确定性判断。

这几个问题决定了Citi推演的三个变化（温和整合、价格战软化、成本控制强化）是否真的会发生，以及发生的节奏和幅度。如果出口不及预期或国内消费进一步恶化，政府刺激重回视野，那么竞争格局的改善逻辑可能被推迟甚至逆转。

基于Citi这份报告的信号，一个更务实的分析框架应该是：放弃对“谁增长最快”的追逐，转向对“谁最稳定”的评估。

第一，订单稳定性指标。关注周度订单的波动率，而非月度总量的高低。波动率越低，说明品牌对促销和脉冲的依赖越小。

第二，成本控制能力。在价格战软化的假设下，成本控制能力将直接转化为利润弹性。关注单车折旧摊销、研发费用率、以及供应链管理效率。

\subsection{[R020] Citi：半导体设备市场真正被低估的，是2028年2500亿美元WFE背后的结构性需求切换}
{\small\color{refgray}归类：AI与算力：基础设施定价权重估与硬件ROE结构性上升 / 生物医药与半导体：全球价值转化面临重估}

Citi：半导体设备市场真正被低估的，是2028年2500亿美元WFE背后的结构性需求切换

半导体设备投资者的注意力，长期集中在2026年和2027年。这可以理解：AI基建正在加速，台积电、三星和英特尔都在扩产，WFE（晶圆厂设备支出）的增长路径看起来足够清晰。但Citi这份6月17日发布的最新研报，提出了一个值得认真对待的判断：市场对2028年的WFE预期可能显著偏低，而且，真正驱动增长的逻辑正在发生变化——不是简单的AI需求延续，而是一种由DRAM瓶颈引发的、对NAND设备需求的系统性重估。

Citi在报告中引入了2028年WFE的牛市预测——2500亿美元。这个数字本身并非关键，关键的是支撑它的两个结构变化：第一，AI基建的资本开支增速虽然会从2026年的84\%放缓至2028年的38\%，但绝对规模仍然巨大，且正在从训练侧转向推理侧；第二，也是更值得关注的，是DRAM供应瓶颈正在迫使产业链重新设计内存架构，从而为NAND设备打开了一个此前未被充分定价的需求空间。

这不是一个“AI继续增长”的故事，而是一个“AI的增长形态正在改变，设备需求结构必须随之调整”的故事。

这份报告的核心贡献，不是给出了三个目标价——AMAT 710美元、LRCX 450美元、KLAC 290美元——而是提供了一个新的分析框架：当市场开始以2028年的视角审视这些公司时，哪些假设需要被重新审视，哪些风险尚未被定价。

1. 2028年WFE达到2500亿美元，这个预测的底气来自哪里

Citi基于其超大规模云服务商资本开支模型，给出了 年WFE的牛市路径：1450亿美元、2000亿美元、2500亿美元。2028年相对于2027年仍有25\%的增长。在设备投资周期通常被认为会在 年见顶的共识下，这是一个相当大胆的预测。

支撑这一预测的，是三个相互独立的驱动力。第一，AI相关的WFE仍然是最大的增量来源，Citi预计其将从2025年的350亿美元增长至2028年的1670亿美元（牛市情景），占2028年总WFE的66\%。第二，中国市场的WFE被假设为基本持平，在360-420亿美元的区间内波动，既不构成拖累，也不提供额外弹性。第三，非AI领域的WFE被假设为每年增长5\%，从390亿美元增至450亿美元。

真正有意思的部分，是Citi对2028年WFE为何“更乐观”的解释：TSMC和存储制造商的持续产能限制与扩张，以及Intel和三星代工厂的最新进展。报告提到了Intel代工在Google和Apple获得不同规模的订单，三星加速泰勒工厂建设并与Google合作TPU v10的2nm内存I/O芯片，以及与比亚迪在2nm和4nm自动驾驶芯片上的讨论。

这些信息意味着什么？意味着AI芯片的制造，正在从台积电一家独大，向一个多供应商的格局演变。每一家新进入的代工厂都需要大量设备来建立产能，而这个过程不会在2027年结束。2028年的WFE增长，本质上是在为2027年之后仍然存在的产能缺口做对冲。

2. DRAM瓶颈不是需求疲软的信号，而是NAND设备需求爆发的起点

这份报告最值得深读的部分，是Citi对DRAM瓶颈与NAND需求之间关系的分析。这不是一个常规的供需平衡分析，而是一个关于架构演变的判断。

Citi的核心逻辑是：Agentic AI的兴起正在结构性推高NAND需求，因为多步推理工作流大幅扩大了KV缓存（key-value cache）的占用空间，使得总内存需求远远超出了高成本的HBM和DRAM能够高效支持的范围。在DRAM供应受限、价格高企的环境下，这种压力已经在推动架构层面的取舍。

报告引用了一个关键证据：Nvidia在其Vera Rubin NVL72系统中，将SoCAMM2 DRAM容量削减了约50\%，原因是

[... middle omitted ...]

存储”向“结构性AI基础设施组件”的角色转变。如果这个判断成立，那么Lam Research（LRCX）作为NAND设备领域的主要玩家，其估值逻辑需要被重新审视——这也是Citi将其目标PE从37倍提升至40倍、目标价从315美元提升至450美元的重要原因。

3. KV缓存卸载正在重塑内存层次结构，NAND从存储变成推理内存

为了理解DRAM瓶颈如何转化为NAND需求，需要理解一个正在快速发展的技术趋势：KV缓存卸载（KV cache offloading）。

简单来说，KV缓存卸载是指将LLM推理过程中产生的中间注意力状态（key-value张量）从HBM动态迁移到更低成本、更高容量的存储层级——先是DRAM，然后是NAND闪存。随着Agentic AI工作负载的增长，KV缓存可以迅速占满GPU内存，而卸载通过创建一个分层系统来缓解这个“内存墙”问题：热数据留在GPU的HBM中，较冷、较少访问的上下文被迁移到DRAM，再进一步迁移到SSD，并在需要时重新加载。

\subsection{[R021] GS：香水市场的下一个增长波，不是品类红利，而是打造爆款的能力}
{\small\color{refgray}归类：消费行业：香水市场进入爆款竞赛模式}

GS：香水市场的下一个增长波，不是品类红利，而是打造爆款的能力

香水行业的增长正在回归常态。疫情后那波由消费者基数扩张和使用频率提升驱动的增长，已经让位于一个更考验内功的阶段。GS欧洲消费品团队在最新发布的研报中给出了一个清晰判断：当品类增速从2023年的8\%回落至2025年的3.5\%，竞争强度却从2019年的2500个新品翻倍至2025年的6000个时，香水品牌的价值分化将不再取决于“是否站在风口上”，而取决于能否持续创造出新的爆款支柱。

这份报告的核心贡献，不是预测行业增速，而是提供了一个分析框架：在供给驱动的品类里，增长本质上是一场“新品转化率”的竞赛。谁能在未来三年内推出并规模化新的爆款，谁就能在估值中获得重新定价。

1. 品类增长正常化后，真正的变量从“需求扩张”转向“供给质量”

疫情后的香水市场经历了一轮罕见的量价齐升。消费者首次购买香水的年龄提前了，使用频率也显著增加。GS引用的波士顿咨询数据显示，从2000年代到2020年代，美国消费者首次购买香水的平均年龄从13岁下降到了11.5岁。这个“年轻化+扩圈”的组合，是过去几年行业高增长的底层逻辑。

但这一结构性红利正在边际递减。当渗透率已经提升、消费者基数不再快速扩大时，增长动力就必须从“让更多人买”转向“让更多人买我的产品”。这正是GS报告的核心逻辑切换点：当行业进入存量博弈阶段，供给质量——即新品能否成为爆款——将成为区分赢家与输家的关键。

数据显示，2019年欧莱雅推出Libre等三款新品，当年香水销售增长8\%。六年后，Libre成为全球女性香水第一品牌，超越香奈儿。这不是偶然。在2010年代中期，欧莱雅推出La Vie Est Belle，带动女性香水在2014年实现两位数增长，2015年增长超过9\%，至今仍是兰蔻的支柱产品。GS将这类产品定义为“爆款支柱”——不是一次性促销，而是能够持续贡献销售、提升品牌势能的长期资产。

所以呢？投资者需要关注的核心指标，不再是行业增速或市场份额的静态数据，而是一家公司新品转化为爆款的“命中率”和“规模化能力”。

GS对欧莱雅、Puig和Interparfums未来三年的新品节奏给出了具体预判，这是报告中最具操作价值的部分。

欧莱雅方面，阿玛尼自2020年My Way之后没有推出过新的爆款支柱，未来三年内推出新品的概率很高。华伦天奴自2019年Born in Roma之后也需要拓宽产品矩阵。兰蔻自2019年Idole之后同样没有重大新品。此外，欧莱雅还肩负着培育Kering Beauty香水组合（包括葆蝶家、巴黎世家）的任务，最晚到FY28需要推出古驰香水。GS特别提到，欧莱雅美国奢侈品部门CEO已暗示2026年下半年将有新品发布。

Puig的看点更为集中。Jean Paul Gaultier将在2026年下半年推出女性爆款，以缩小其在女香领域的弱势。Rabanne预计在2027年推出新的男性支柱，因为自2021年Phantom之后该品牌在男香领域没有重大动作。Carolina Herrera的La Bomba刚刚上市，GS认为其搜索热度在上市前12个月的表现已超过Good Girl同期水平，这是一个积极信号。

Interparfums的战略则更加稳健——在其四大品牌Jimmy Choo、Coach、Montblanc和Lacoste之外，即将推出Longchamp香水。Jimmy Choo的案例尤其值得关注：通过持续推出新支柱（如I Want Choo），该品牌销售额已从2011年的2900万欧元增长到2025年的2.28亿欧元，即使在上市四年后，I Want Choo在2025年仍实现了27\%的增长。

但这里有一个关键问题报告没有完全回答：这些新品究竟是在吸引新消费者，还是在蚕

[... middle omitted ...]

分由其无与伦比的执行力和更广泛的品类布局支撑，但与其历史水平相比，这一溢价空间仍然可观。

为什么市场可能低估了这些公司？原因在于，香水行业的新品效应并非线性。一个成功的爆款可以在上市后持续贡献增长数年，甚至十年以上。Good Girl用了六年时间才达到3\%的市场份额，但此后一直保持在高位。1 Million在第一年就达到1\%的份额，随后十年持续攀升。这种“长尾效应”意味着，当前投入的新品，其价值可能要到三到五年后才在财务数据中充分体现。而市场往往只关注短期业绩。

GS给出的催化剂明确：Puig将于10月28日举行资本市场日，预计将阐述其在女性香水领域的战略以及重振Rabanne的计划。欧莱雅的资本市场日则定于12月1日，预计将聚焦AI议程。这两个事件可能成为市场重新定价的触发点。

4. 报告尚未完全回答的关键问题：爆款的可复制性是否存在边界？

尽管GS对几大头部公司的新品管线给出了乐观预判，但报告本身也隐含了一个尚未完全解答的问题：爆款的可复制性是否存在边界？

\subsection{[R022] NOM：人民币中间价定价模型释放的信号比汇率本身更重要}
{\small\color{refgray}归类：区域市场：亚洲外汇储备分化与东南亚互联网韧性}

对于关注人民币汇率走势的投资者而言，盯住每日中间价的变化已经成为一种条件反射。但NOM最新发布的这份USD/CNY定价模型报告揭示了一个更深层的判断：真正重要的不是中间价在某个交易日调升或调降了多少个基点，而是定价模型本身的误差正在系统性收窄，以及逆周期因子正在扮演一个更微妙的角色。

这份报告的核心主张是，人民币中间价的定价机制正在经历一个从“被动对冲”向“主动校准”的转变。市场习惯于将中间价视为政策意图的风向标，但NOM的模型拆解显示，当前定价的驱动力已经从单一的政策维稳转向了更为复杂的多因子均衡。这意味着，投资者需要更新自己的分析框架，不能再用过去的逻辑来解读中间价的每一个波动。

为什么现在需要关注这个问题？因为2026年下半年中国将迎来一系列关键的政治和经济议程：7月底的政治局经济工作会议、11月在深圳举办的APEC领导人会议、以及年底的中央经济工作会议。在这些关键时点前后，中间价的定价逻辑可能会发生微妙但重要的切换。NOM的模型恰恰为提前捕捉这种切换提供了一个量化工具。

以下是我们从这份报告中提炼出的四个关键洞察，以及一个报告尚未完全回答的关键问题。

1. 定价模型的预测误差正在系统性收窄，这暗示市场定价与政策意图正在趋于一致

NOM模型显示，在不考虑逆周期因子调整的情况下，其投影值为6.7761，较前一日的6.8096大幅调低了335个基点。但更值得关注的不是单日变化，而是图2中展示的模型误差走势。

从2025年初到2026年初，模型的日度误差从约-1800个基点的高位持续收窄，到2026年一季度已经趋近于零，甚至在4月份出现了约+600个基点的正向误差。这种误差的系统性收敛，意味着NOM的定价模型——一个基于隔夜主要货币对变动、韩元、澳元等一篮子货币加权贡献的量化框架——正在越来越准确地拟合实际的中间价设定。

这对投资者的含义是双重的。一方面，它表明市场预期与央行实际操作的偏离度正在缩小，人民币汇率的定价正在变得更加可预测。另一方面，误差从负值转向正值，暗示模型在某个阶段可能低估了中间价的调升幅度，这反过来又说明央行在特定时点可能比模型预期的更加积极。

这里需要强调的是，误差收窄本身并不代表央行放松了对汇率的管理，而是表明管理的逻辑正在变得更加透明和可量化。对于交易员而言，这意味着过去那种“中间价超预期”带来的交易机会正在减少，但同时也意味着模型本身可以成为更可靠的参考基准。

2. 韩元成为仅次于欧元的关键驱动因子，这揭示了人民币定价的“区域锚”效应正在增强

在图1中，NOM列出了对投影变化贡献最大的四个货币：欧元贡献60个基点，韩元贡献46个基点，澳元13个基点，英镑9个基点。韩元的权重仅次于欧元，且远超其他非美主要货币。

这一发现值得深入解读。传统上，人民币中间价的定价更多参考美元指数和欧元/美元走势。但韩元权重的显著上升，暗示人民币定价正在越来越多地吸收区域性的竞争性贬值压力。韩国作为中国在出口市场上的直接竞争对手，其汇率变动对中国制造业的竞争力有着直接的影响。

当韩元大幅贬值时，如果人民币不做出相应调整，中国出口商将在价格竞争中处于劣势。因此，模型赋予韩元高权重，可以被视为央行在定价时主动纳入了“出口竞争力平价”的考量。

这一逻辑的延伸含义是，未来人民币中间价的波动将不仅受制于中美利差和美元走势，还会越来越多地受到整个亚洲出口板块汇率变动的影响。对于关注中国出口企业盈利能力的投资者而言，盯住韩元、新台币等亚洲货币的走势，可能比单纯盯住美元指数更有预测价值。

3. 逆周期因子的调整幅度正在变小，但它的存在本身就是一个关键变量

NOM模型在加入逆周期因子后的投影值为6.7992，比前一日实际定盘价低了104个基点。这意味着，模型估算的逆周期因子调整幅度为104个基点

[... middle omitted ...]

。它既不是完全放任市场定价，也不是强力干预。这种状态可能反映了央行在当前阶段的政策权衡：既要避免人民币过快升值损害出口，又要防止贬值预期失控引发资本外流。

4. 2026年下半年的政治日程，将是检验中间价定价逻辑稳定性的关键窗口

NOM的报告特别列出了一个日历事件表，其中包括7月底的政治局会议、10月国庆黄金周、11月APEC深圳峰会、以及年底的中央经济工作会议和可能的习近平访美行程。

这些事件为中间价的定价逻辑设置了明确的测试节点。在APEC这样的多边外交场合前后，汇率稳定往往是政策优先目标。而在中央经济工作会议期间，汇率政策基调的表述可能会直接影响中间价的定价中枢。

报告没有直接给出这些事件对中间价的具体影响预测，但提供了一个分析框架：投资者需要关注的是，在关键事件前后，模型误差是否会突然扩大，以及逆周期因子的调整幅度是否会显著变化。如果误差在事件前迅速收窄，说明市场预期与政策意图提前达成了共识；如果误差在事件后突然扩大，则可能意味着定价逻辑正在发生切换。

\subsection{[R023] GS：银行股真正的防御性，来自资产负债表结构而非规模}
{\small\color{refgray}归类：宏观与利率：中国增长逻辑切换与全球通胀滞后效应 / 风险偏好与政策：银行股结构溢价与香港楼市政策误读}

6月17日，国家金融监督管理总局、人民银行、证监会、外汇管理局等主要金融监管机构负责人在陆家嘴论坛上发表了系列讲话。GS在随后的研报中提炼了一个核心判断：市场对银行股的定价，可能正在经历从“规模溢价”向“结构溢价”的切换。

这不是一次普通的政策解读。这份报告的价值在于，它把监管层释放的分散信号，拼成了一幅关于未来三到五年中国金融体系演变的结构性地图。对投资者而言，真正重要的不是某条政策是松是紧，而是这些政策共同指向了一个方向——银行资产负债表的“质量”和“结构”，正在取代“增速”和“规模”，成为决定估值分化的核心变量。

GS报告中最具穿透力的观察，是人民银行对融资结构变化的定量描述。

数据显示，以贷款为主的间接融资在社融增量中的占比，已从超过80\%降至2025年的45\%；在社融存量中的占比，也从接近100\%降至三分之二。这不是一个短期现象。新发放贷款的结构变化更为剧烈：投向房地产和基础设施领域的新增贷款占比，从超过60\%骤降至当前的10\%；而投向科技金融、绿色金融、普惠金融、养老金融、数字金融这“五篇大文章”的新增贷款占比，已跃升至70\%以上。

这意味着什么？GS援引人行的判断称，随着经济结构从依赖中长期贷款的重资产传统行业，转向依赖多元化融资方式的轻资产新兴行业，维持过去的信贷增速既困难也无必要。“更慢、更高质量的贷款增长”正在成为新常态。

对投资者而言，这个判断的含金量在于：它把信用增速放缓从“周期性利空”重新定义为“结构性常态”。那些仍在押注信贷脉冲反弹来驱动银行股估值的逻辑，可能需要根本性的修正。在资产负债表扩张放缓的环境中，银行自身的资产质量韧性和负债结构稳定性，将变得前所未有的重要。

GS报告指出，人民银行宣布将常备借贷便利利率调整为7天逆回购操作利率加减25个基点，相比此前的加50个基点/减20个基点，利率走廊从70个基点收窄至50个基点。以当前1.4\%的政策利率计算，新的利率走廊区间为1.15\%-1.65\%。

这个看似技术性的调整，对不同类型的银行影响迥异。GS的测算逻辑非常清晰：短期利率上限降低和走廊收窄，有助于平抑短期利率波动，从而稳定银行的同业负债成本。尤其是在季末、税期等流动性收紧时段，这种稳定效应会更加显著。

关键的分化在于：中小银行通常存款基础较弱，负债结构中同业融资占比更高，因此受益程度更大。GS覆盖的银行中，兴业银行、华夏银行和宁波银行的同业负债占比分别高达29\%、26\%和22\%，而四大国有银行平均仅为15\%。

这个数据背后是一个重要的投资启示：在利率走廊收窄的政策环境下，同业负债占比较高的银行将获得更直接的净息差保护。这不是一个短期交易信号，而是一个结构性优势。那些能够有效利用这一政策窗口优化负债结构的银行，将在未来几个季度的盈利竞争中占据先机。

GS报告披露，人民银行已授权工商银行、农业银行、中国银行、建设银行、交通银行和中信银行等六家银行，在上海自贸区开展离岸人民币外汇交易试点，使用中国外汇交易中心平台。

从财务影响看，GS的测算非常克制：以国际业务最领先的中国银行为参照，其整个国际结算业务占手续费收入不到7\%，占总收入不到1\%。因此，这项新业务短期内对银行盈利的贡献有限。

但如果我们把视角拉长，这个试点的战略含义远超其财务比重。它标志着在岸和离岸人民币外汇市场的双向开放迈出了实质性的一步。对参与试点的银行而言，这不仅意味着新的手续费收入增长点，更重要的是能够通过满足境外投资者的在岸外汇对冲需求，增强客户粘性，拓展国际业务纵深。

这里有一个报告没有完全展开的问题：随着人民币国际化进程加速，哪些银行能够将先发优势转化为持续的护城河？试点资格本身是起点，但真正决定胜负的，是银行在跨境金融服务能力、产品设计能力和风险管理体系上的综合实力。这个问题，值得投

[... middle omitted ...]

细读。GS强调，央行明确表示这并非常态化的流动性供给，而且非银机构需要提供高质量抵押品以防止道德风险。这意味着，这是一个“保险”而非“提款机”。

它的关键意义在于：过去当债券市场出现系统性压力时，非银机构的主要流动性来源是银行。一旦银行收紧资金拆借，非银机构往往被迫抛售高流动性资产，进一步加剧市场波动和净值下跌。2022年底的理财产品赎回潮就是一个典型案例。如果这个新工具到位，央行的应急流动性注入将帮助稳定债券市场，进而稳定银行的债券投资收益，尤其是交易性金融资产和可供出售金融资产的估值。

GS的数据显示，宁波银行、南京银行和中国银行的交易账户占投资资产比例最高，分别为67\%、67\%和55\%。这意味着，这些银行对债券市场稳定性的敏感度最高，也将是这项新工具最直接的受益者。

GS报告对资本市场政策的判断同样清晰：供给端工具将加速扩容，同时监管约束将进一步收紧。对券商而言，核心含义是增量机会将集中于那些同时具备强大机构业务能力、跨境平台和稳健风险管理框架的头部机构。

\subsection{[R024] 摩根斯坦利：消费复苏的真正考验不在总量，而在结构分化的再定价}
{\small\color{refgray}归类：地产与消费：K型分化加剧，一线回暖难掩全国压力}

摩根斯坦利：消费复苏的真正考验不在总量，而在结构分化的再定价

2026年5月的中国零售数据，是疫后首个同比负增长的月份——同比下降0.6\%，低于市场预期的-0.2\%和4月的+0.2\%。这个数字本身并不令人震惊，但它揭示了一个比“消费疲软”更值得深挖的结构性事实：总量数字掩盖了供给端和渠道端正在发生的重新定价，而市场对这一重新定价的认知仍停留在需求侧叙事上。

这份摩根斯坦利研报的价值，不在于它告诉你消费不好，而在于它提供了足够精细的颗粒度，让读者看清“哪里在加速出清、哪里在积蓄动能、哪里只是暂时性扰动”。以下是我们从中提取的五个关键判断。

1. 总量负增长背后的真正信号：补贴退坡后的真实需求压力测试

5月零售总额同比下降0.6\%，但更值得关注的是扣除汽车后的同比增速仍为正（+1.1\%），以及对比2019年的CAGR从4月的2.9\%小幅回升至3.2\%。这说明总量下滑的主因并非消费意愿全面崩塌，而是特定品类的基数效应和补贴政策到期产生了集中拖累。

电子与家电品类同比-15.6\%，与4月的-15.1\%基本持平，连续多月处于深度负增长区间。这一品类的疲软可以追溯到两个叠加因素：一是 年以旧换新补贴的刺激效应正在快速衰减，二是房地产相关需求（装修、家电配套）尚未出现实质性回暖。汽车品类同比-16.1\%，同样受到去年高基数和新一轮价格战的不确定性压制。

从MECE的角度看，5月数据可以被拆解为三类：政策驱动型需求的退潮、房地产关联需求的持续低迷、以及日常消费的韧性测试。三者的分化程度远超市场此前预期。

这意味着，对于投资者而言，盯住总量数字做判断的风险正在上升。真正重要的是识别哪些品类正在经历“真实需求支撑的调整”，哪些只是“政策退出的噪音”。

2. 黄金珠宝的跌幅收窄并非消费回暖，而是价格效应掩盖了量价双杀

黄金珠宝品类从4月的-21.3\%收窄至5月的-8.9\%，看起来是一个积极的边际改善。但这份研报的数据结构提示我们：这个改善需要谨慎解读。

黄金珠宝的零售额同比增速同时受金价和销量两个变量影响。2025年二季度金价处于高位调整期，而2026年5月的金价基数有所回落，这意味着同比跌幅收窄可能更多来自价格效应，而非销量复苏。结合该品类对比2019年的CAGR从4月的3.0\%微降至5月的2.8\%，量的动能仍在走弱。

更重要的是，黄金珠宝属于典型的“可选消费+资产属性”双重定位品类。当消费者信心处于修复初期，这类品类的购买决策往往被推迟而非取消。摩根斯坦利在报告中提到的最新消费者调查显示“情绪温和改善但谨慎依旧”，这与黄金珠宝数据反映的“边际改善但动能不足”高度一致。

这一品类的真正拐点，需要等到量价同时企稳的信号出现，而不是单纯依赖同比跌幅的算术变化。

3. 线上零售的加速是消费降级还是渠道重构？答案决定投资方向

5月线上零售同比增长3.4\%，较4月的2.3\%明显加速，而同期线下零售（限额以上企业）同比-5.2\%。这一反差在过去几个月中持续存在，但5月的差距进一步拉大。

一种常见的解读是“消费者在降级，转向更便宜的线上渠道”。但这份研报的数据并不完全支持这个结论。对比2019年的CAGR，线上零售从4月的9.8\%提升至11.8\%，而线下限额以上企业从3.0\%微降至3.8\%。两者都在改善，只是线上改善的斜率更陡。

更合理的解释是：渠道结构本身正在发生不可逆的变化。疫情后消费者养成的线上购物习惯并未消退，叠加直播电商、即时零售等新业态的渗透率仍在上升，线上渠道的增速高于线下并非周期性现象，而是结构性趋势。真正值得关注的是：哪些品类正在被线上渠道侵蚀利润，哪些品类正在通过线上渠道实现增量。

报告中没有直接展开的是，不同消费品类在线上线下的利润率差异巨大。对于品牌商而言，线上占比的提升可能意味着营销费

[... middle omitted ...]

市场中创造了新的消费场景和需求。这份研报中提到的蒙牛（2319.HK）和伊利（600887.SS）被列为“供给重新校准与需求改善”的代表，正是基于这一逻辑。

软饮料品类的案例揭示了一个更通用的投资框架：在消费总量增长放缓的环境下，能够通过产品迭代、渠道变革或品牌升级来“创造需求”的企业，将获得超越行业平均水平的回报。反之，依赖行业β增长的企业将面临持续的压力测试。

5. 报告尚未完全回答的关键问题：补贴退坡后的“真空期”会持续多久？

摩根斯坦利在报告中明确指出了补贴政策对电子、家电、汽车等品类的支撑作用正在减弱，但并未给出一个清晰的“政策真空期”持续时间判断。这是一个值得深入讨论的未解问题。

从历史经验看，消费品补贴政策的退坡通常会经历三个阶段：第一个月出现明显的销量下滑，第二至第三个月出现渠道库存调整（经销商主动去库存），第四个月后真实需求开始显现。如果真实需求足够强劲，市场会在第三至第四个月自发企稳；如果真实需求只是被补贴提前透支，则下滑可能持续更长时间。

\subsection{[R025] 摩根斯坦利：市场对“和平”的定价已经过头，但油价真正的支撑尚未被计入}
{\small\color{refgray}归类：能源与大宗：和平交易过度定价，供给约束支撑三季度油价}

摩根斯坦利：市场对“和平”的定价已经过头，但油价真正的支撑尚未被计入

这份摩根斯坦利研报的标题本身就是一句判断：“Is Peace Mispriced?”——和平是否被错误定价？在美伊谅解备忘录即将签署、WTI自4月初已下跌29\%的背景下，市场似乎已经在为“和平溢价”买单。但摩根斯坦利的分析揭示了一个更微妙的图景：当前股价所隐含的油价预期（约66美元/桶WTI）已经显著低于12个月远期曲线（约75美元），而基本面——尤其是三季度仍将存在的供需缺口——并没有被充分反映。真正被市场低估的，不是冲突持续的概率，而是即便和平到来，供给恢复的速度与库存重建的规模。

这份报告的核心价值不在于判断地缘政治走向，而在于提供了一个量化框架：把“和平”拆解为可追踪的变量——产量恢复节奏、库存回补需求、以及不同油价情景下的企业现金流。对于产业决策者和长期投资者而言，当前的回调可能不是风险出清，而是重新评估资产定价的窗口。

1. 和平协议签署只是第一步，供给恢复的物理约束才是三季度油价的真正支撑

市场对美伊达成谅解备忘录的反应是直观的：油价下跌。但摩根斯坦利原油策略师Martijn Rats的假设条件值得仔细拆解。该机构预计，即便霍尔木兹海峡近期重新开放，油轮运输恢复正常需要数周时间；到9月伊朗产量恢复50\%，到12月恢复80\%。这意味着三季度全球原油市场仍将维持约340万桶/日的平均缺口（仅OECD国家看，缺口约110万桶/日）。

这个测算的隐含意义是：和平协议本身并不等于供给立即释放。从协议签署到实际产量回流，中间存在一个数月的“真空期”。而在此期间，全球库存已经因冲突期间的大规模释放而处于低位。摩根斯坦利的图表显示，OECD原油库存已经跌至5年平均水平以下，且仍在加速去库。三季度的高温需求、叠加供给恢复的时滞，意味着油价在90美元附近仍将获得支撑，即便在和平情景下。

对投资者的含义：当前市场对“和平”的线性外推可能过度。三季度供需基本面并不支持油价持续下行，而库存重建的需求将在四季度及之后形成新的价格底部。

2. 当前股价隐含的油价预期已低于远期曲线13\%，但市场并未为“库存重建”定价

摩根斯坦利通过逆向测算发现，美国油气生产商的当前估值隐含的WTI均价约为66美元/桶，比12个月远期曲线（约75美元）低约13\%。这个差距在历史上并不罕见，但当前的特殊性在于：远期曲线本身可能已经包含了过多的“和平折价”。

更关键的是，该机构指出，即便在和平情景下，全球仍需补充自冲突以来消耗的大量库存。这并非一个线性过程。库存重建意味着额外的需求增量，而这一增量在当前的远期曲线中可能并未被充分定价。摩根斯坦利的敏感性分析显示，WTI每变动10美元/桶，覆盖范围内生产商的2027年自由现金流收益率将变动约3个百分点。在72美元的远期曲线下，美国油气生产商的2027年自由现金流收益率中位数为13\%，这一水平在历史上属于有吸引力的区间。

这意味着：当前股价已经为“和平”支付了溢价，但尚未为“和平后的库存重建”定价。如果库存重建推动实际油价高于当前远期曲线，则当前估值可能低估了企业的真实盈利潜力。

3. 美国油气生产商的现金流韧性被低估，但这一判断依赖于两个尚未验证的假设

摩根斯坦利对美国油气生产商的整体判断是偏积极的：强劲的自由现金流叠加有吸引力的估值，回调创造了增加配置的机会。但这一判断依赖于两个尚未充分验证的假设。

第一，需求韧性。报告中的图表显示，美国隐含成品油需求迄今表现良好，汽油消费仍在正常范围内。但中国海运石油净进口量持续下降，6月可能进一步走低。全球需求的分化——美国强、中国弱——能否持续，是影响油价中枢的关键变量。如果中国需求疲软进一步蔓延，或美国经济出现放缓迹象，当前的需求假设可能需要下调。

第二，生产商资

[... middle omitted ...]

在于：过去几个月，SPR释放是抑制油价上涨的重要力量；而一旦释放大幅减少，市场将失去一个重要的缓冲机制。

结合前述三季度仍将存在的供需缺口，SPR释放的减少意味着市场将更加依赖商业库存和新增产量来满足需求。而商业库存目前已经处于低位，新增产量（尤其是伊朗产量）又需要时间才能完全恢复。这个组合在三季度末至四季度初可能形成一个供给紧张的窗口。

对投资者的含义：不应只关注地缘政治事件的短期影响，而应跟踪SPR释放节奏、库存水平、以及产量恢复的月度数据。这些“慢变量”的累积效应可能比和平协议本身对油价有更持久的影响。

5. 一个尚未被充分讨论的关键问题：天然气市场的结构性变化是否被低估？

摩根斯坦利的报告主要聚焦于原油，但对天然气生产商也进行了敏感性分析。在3.50美元/百万英热单位的远期曲线下，天然气生产商的2027年自由现金流收益率中位数为9\%，略低于原油生产商。但天然气市场的定价逻辑与原油存在根本差异：天然气更多受区域供需、LNG出口能力、以及季节性因素的影响。

\subsection{[R026] JPM：市场对内地资金买港楼的恐慌，可能高估了政策冲击}
{\small\color{refgray}归类：风险偏好与政策：银行股结构溢价与香港楼市政策误读}

一份来自JPM的专家电话会纪要，正在香港地产投资圈内流传。触发点是中国国务院发布的《境外直接投资条例》（即“837号文”）。这份文件让不少投资者担忧：内地居民购买香港物业的通道是否会被彻底堵死？香港楼市刚刚开始的复苏周期，是否会被政策“截胡”？

JPM在6月16日发布的这份研报，核心判断出乎许多人的意料：837号文并非针对内地居民购买香港物业，现行监管框架实质上没有发生根本性变化。 但报告同时指出，真正值得关注的，不是政策本身，而是执行层面可能带来的“寒蝉效应”——尤其是对“内地税务居民”的全球收入征税风险，以及CRS信息交换对资产透明度的潜在影响。

这篇研报的价值，不在于它给出了一个简单的“利好”或“利空”结论，而在于它通过一场与北京律师的深度对话，拆解了一个高度模糊的监管灰色地带。本文尝试提炼这份报告中最具洞察力的六个层次，并留下一个关键问题供读者自行判断。

1. 837号文并非针对香港楼市，它的真正靶点是跨境科技并购与个人境外投资监管空白

许多市场参与者将837号文直接解读为“打压内地人买港楼”，这可能是误读。JPM邀请的北京律师明确指出，这份文件的出台背景，源于一个标志性跨境案例：一家美国公司拟收购某中国AI公司，交易暴露出两个监管漏洞——个人境外持股的合规性，以及核心技术未经审批的跨境转移。

837号文本质上是一次“顶层设计”的补丁。它首次在国家层面将个人境外投资纳入监管框架，并强化了技术出口管制。对于香港物业购买，该文件并未设置新的禁令。律师的观点很清晰：内地居民用境内资金购买境外房产，在现行外汇管制框架下“从未被允许过，未来也不太可能被允许”。837号文只是重申了现状，而非收紧。

这意味着，市场此前对“政策突变”的定价可能过度。但“未收紧”不等于“无风险”——接下来的问题在于，执行层面的穿透力度是否会加强。

2. “资金源头”是唯一的合规分水岭，但灰色地带的清理正在加速

JPM研报反复强调一个核心逻辑：监管关注的是“资金是否从境内流出”，而非“资产是否在香港”。只要购买香港物业的资金来源是合法的境外收入（如港股IPO收益、境外薪酬、股票期权行权所得），内地居民在香港购房并不受限制。反之，任何涉及人民币兑换后跨境转移用于购房的行为，均属资本项目下的违规操作。

实践中，大量内地买家通过“对敲”（境内人民币换境外港币）、分拆5万美元购汇额度、或通过关联公司服务贸易转移资金等方式完成交易。律师指出，这些方式大多处于灰色地带，合规风险正在上升。银行的跨境汇款合规审查已明显趋严，“对敲”等模式的集中资金流动更容易触发监管关注。

这里的关键洞察是：监管的“靶心”不是资产本身，而是资金跨境流动的合规性。 对于已经持有香港物业的内地买家，只要资金来源问题不涉及重大违规，被强制处置资产的可能性极低。罚款通常为违规汇出金额的5\%-10\%，且不要求资金回流。

3. “内地税务居民”身份才是最大的潜在风险变量，而非837号文本身

研报中最具前瞻性的洞察，可能不是对外汇管制的分析，而是对税收风险的提示。律师指出，一个持有香港物业的内地居民，如果被界定为“内地税务居民”，其全球收入（包括香港物业的租金收入和资本利得）理论上需按20\%的税率在中国缴纳个人所得税。

问题在于，很多人并未意识到自己仍是“内地税务居民”。只要持有内地户口（户籍），即使长期在香港工作生活，税务居民身份通常不会自动改变。这意味着，未来如果内地税务机关加强对境外资产的税务稽查，持有香港物业的隐形成本可能显著上升——20\%的税率足以大幅降低香港物业的净回报率。

这一风险目前尚未被市场充分定价。它与837号文无关，但可能成为比外汇管制更实质的“灰犀牛”。

4. CRS短期内不会直接穿透房地产，但间接监控手段已经存在

另一个让市场紧

[... middle omitted ...]

渐进的。 真正需要担心的，不是某一天房地产突然被纳入CRS，而是税务机关通过资金流与公开信息的交叉验证，逐步构建起对个人境外资产的监控能力。

5. 对香港楼市的影响：短期扰动而非趋势逆转，但需警惕利率与股市的叠加效应

JPM估算，非香港本地的内地买家占香港住宅交易量的5\%-10\%（按金额计约10\%-15\%）。即使这部分需求因政策不确定性受到抑制，报告认为也不足以逆转香港楼市的复苏周期。真正需要关注的，是另外两个在研报中并列提及的逆风因素：利率上升预期的升温，以及恒生指数持续低迷对楼价的传导效应。

历史数据显示，香港房价与恒生指数高度相关。当前港股表现疲弱，叠加美联储降息预期推迟，这两个因素对楼市的压制可能远大于837号文本身。JPM预期，香港地产股短期将维持区间震荡，直到市场看到：一手房销售率、二手成交量与房价未受明显冲击；或者利率预期出现松动。

在选股层面，报告给出了相对清晰的排序：短期看好长实集团与会德丰，看淡恒基地产与新世界发展，逢低建议吸纳新鸿基地产。

\subsection{[R027] 摩根斯坦利：中国材料工业的真正分化不在需求，而在供给侧的刚性约束}
{\small\color{refgray}归类：宏观与利率：中国增长逻辑切换与全球通胀滞后效应}

摩根斯坦利：中国材料工业的真正分化不在需求，而在供给侧的刚性约束

五月工业数据出炉，市场目光大多落在房地产的持续疲软上。新开工面积同比下滑24.7\%，销售面积降幅扩大至14.1\%，固定资产投资同比下跌12.5\%——这些数字确实令人忧虑。但如果只盯着需求端的压力，就很可能错过这份摩根斯坦利研报中最具价值的信号：中国材料工业正在经历一场深刻的分化，而分水岭不在于谁的需求更好，而在于谁的供给已经触及天花板。

摩根斯坦利的分析师团队在这份五月工业产出数据报告中，给出了一个清晰的图景：在粗钢、水泥、玻璃、煤炭产量普遍同比下滑的背景下，电解铝产量却实现了同比1.7\%的增长，且环比仍在上升。这不是偶然的月度波动，而是结构性差异的体现。当多数材料行业还在为过剩产能和低迷需求挣扎时，铝行业已经走到了供给约束的临界点。

这份报告值得深读的原因，不在于它披露了五月的生产数据，而在于它揭示了不同材料子行业在未来12个月将走向截然不同的定价逻辑。对于产业决策者和投资者而言，理解这种分化的根源，比跟踪月度数据波动重要得多。

1. 房地产的拖累仍在深化，但三季度可能比数据本身更值得警惕

五月房地产数据延续了疲弱态势。新开工面积同比下滑24.7\%，虽较四月27.1\%的降幅略有收窄，但绝对水平依然处于深度收缩区间。销售面积同比下降14.1\%，较四月10.3\%的降幅进一步恶化。竣工面积同比下降19.6\%，与四月的19\%基本持平。

摩根斯坦利的中国房地产团队给出了一个值得高度重视的判断：预计三季度将更加疲弱。这个判断基于三个因素的叠加——居民信心仍然脆弱、政策效应和积压需求正在消退、可售资源减少。团队进一步预测，二手房销售将在三季度转为同比负增长，新房销售则在低基数上继续恶化。

这意味着什么？材料需求的下行压力远未结束。房地产链条上的水泥、玻璃、粗钢等品种，三季度可能面临比上半年更严峻的需求环境。五月水泥产量同比下滑8.1\%，玻璃产量同比下降6.3\%，这些数字在三季度可能继续承压。对于持有这些资产或相关头寸的读者而言，三季度不是“见底”的时间窗口，而是需要更谨慎评估风险暴露的时期。

在五月所有主要材料品种中，电解铝是唯一实现产量同比正增长的品种。产量达到390万吨，同比增长1.7\%，环比增长0.5\%。前五个月累计产量1920万吨，同比增长3.5\%。驱动因素来自辽宁产能复产和内蒙古新产能投放。

但报告传递的关键信息不在于增长本身，而在于增长已经触及上限。摩根斯坦利明确指出，随着运行产能达到天花板，预计整体铝产量将在年内剩余时间保持高位，前提是行业盈利性保持稳健。这句话翻译过来就是：产量已经无法再大幅增加了。

这是铝与其他材料品种最本质的区别。粗钢、水泥、玻璃面临的是需求不足导致的被动减产，而铝面临的是供给约束下的主动限产。当需求端出现任何边际改善时，铝的价格弹性将远高于其他品种。这种供给刚性为铝行业的盈利提供了结构性支撑，即使需求端没有明显起色，供给侧的约束本身就是一个定价支撑因素。

这里仍需验证的一个问题是：如果下游需求（尤其是建筑和交通领域）进一步走弱，铝的库存是否会开始累积？报告没有展开讨论这个情景。但即便出现需求走弱，铝的供给天花板意味着价格下跌空间也比其他品种有限得多。

五月煤炭产量3.972亿吨，同比下降1.7\%，但环比增长3\%。这个环比回升符合季节性规律——电厂为迎峰度夏开始补库。但报告同时指出，五月下旬的一起矿难以及六月加强的安全检查，将使国内煤炭产量在短期内继续受到约束。

这里有一个值得关注的结构性变化。安全监管正在从临时性措施演变为持续性的供给约束。五月煤炭产量同比仍为负增长，说明即使有季节性补库需求，供给端也无法完全释放。火电发电量同比增长2.1\%，占总发电量的60\%，这个比例较四月的62\%有所下降，但

[... middle omitted ...]

1\%，降幅较四月的10.8\%有所收窄，但这主要来自低基数效应。报告提到，华东地区水泥价格在六月初仍然疲软，原因是下游需求疲弱、潮湿天气以及高考影响施工。这三个因素叠加，说明短期内看不到价格修复的催化剂。

玻璃的情况更为严峻。产量同比下降6.3\%，环比增长3.6\%主要来自复产产线的贡献。但报告直言不讳：浮法玻璃基本面依然疲弱，库存高企，需求低迷。过剩压力将继续压制浮法玻璃价格和利润率。

水泥和玻璃的困境在本质上是相同的：需求端受到房地产和基建双重拖累，而供给端又缺乏有效的出清机制。行业的盈利修复只能依靠更大幅度的减产或需求端的实质性好转。从摩根斯坦利的行业评级来看，水泥为“中性”（In-Line），这实际上意味着缺乏向上弹性。

5. 报告尚未完全回答的关键问题：基建投资的下行是暂时的还是趋势性的？

五月固定资产投资同比下降12.5\%，较四月9.4\%的降幅进一步扩大。公路固定资产投资同比下降13.8\%。这份报告披露了这些数据，但没有深入分析基建投资持续下行的原因。

\subsection{[R028] Citi：市场低估的不是需求疲软，而是供给端正在发生的结构性出清}
{\small\color{refgray}归类：汽车行业：补贴退潮暴露结构性压力，中国车企加速全球渗透}

Citi：市场低估的不是需求疲软，而是供给端正在发生的结构性出清

5月中国乘用车市场销量不及预期，燃油车大幅下滑，新能源车环比持平。如果只看这个表面数据，很容易得出“需求不行”的结论。但Citi在6月16日组织的一场产业专家电话会中，提供了一个更具穿透力的观察框架：真正值得关注的不是短期订单波动，而是供给端正在发生的几重结构性变化——合资品牌面临生存抉择、增程式电动车的技术天花板开始显现、以及头部企业正在通过海外渠道和成本优势重构竞争边界。

这份研报的核心价值不在于给出了多少家车企的周度订单数据，而在于它揭示了这些数据背后的一个关键判断：中国汽车市场的竞争逻辑，正在从“谁能造出更好的车”转向“谁能把规模优势和成本优势转化为定价权和出清能力”。

1. 豪华车市场正在经历的不是周期波动，而是份额的永久性转移

Citi专家电话会披露的一组数据值得反复推敲：BBA（奔驰、宝马、奥迪）上周在华订单已降至不到7000台，而去年同期为 台。这意味着德系豪华品牌在一年内失去了超过三分之一的订单量。同期，蔚来ES9稳态周订单维持在2000台左右，问界系列周订单约7500台，理想L系列约2000台。

这不是简单的“消费降级”或“观望情绪”。这是豪华车市场正在发生的一次结构性权力交接。BBA的订单下滑与国内高端新能源品牌订单的韧性同时发生，说明高端消费需求并未消失，只是转移了去向。过去支撑BBA溢价的品牌认知、经销商网络和服务体系，在电动化与智能化的双重冲击下，正在失去对消费者的说服力。

Citi专家更指出，BBA的周订单已降至“低于7000台”的水平，且没有看到任何企稳迹象。这意味着德系三强在中国市场可能已经越过了某个临界点——当订单量持续低于某个阈值，经销商的现金流、库存周转和定价体系都会进入负向循环。这为后续的合资品牌战略抉择埋下了伏笔。

2. 增程式电动车的增长故事正在触及天花板，这不是一个假设，而是正在发生的事实

Citi专家在问答环节提出了一个容易被忽视却极具冲击力的判断：增程式电动车（EREV）的渗透率下降趋势将持续，因为其电池容量已经触及技术天花板——纯电续航约400公里。而入门级纯电动车（约10万元人民币）已经可以实现600-700公里的续航里程。

这个判断的逻辑链条非常清晰：增程式电动车本质上是一种“过渡方案”，其存在的前提是纯电续航不足以覆盖大部分日常场景。当10万元级别的纯电动车都能做到600公里以上续航时，增程式的“里程焦虑解决方案”价值就被大幅削弱了。更关键的是，增程式同时背负着内燃机和电池两套系统的成本与重量，在纯电续航优势不再明显的情况下，其性价比优势正在被侵蚀。

这一趋势对理想汽车、问界等依赖增程技术路线的品牌影响深远。Citi专家电话会中提及理想L系列周订单约2000台，其中L9不到1000台，并明确表示问界M9和蔚来ES9对L9的销售形成了压制。这并非孤立的竞争事件，而是技术路线演进带来的系统性压力。

3. 比亚迪正在通过海外市场建立第二增长曲线，这是竞争格局中尚未被充分定价的因素

Citi专家电话会揭示了一个有趣的细节：比亚迪上周5万台的订单中，包含了“BYD从经销商回购用于海外销售的小电池插电混动车型”，数量不到5000台。这意味着比亚迪的海外需求不仅真实存在，而且已经强到需要从国内经销渠道“回购”车辆来满足海外订单。

这个操作看似是临时性的库存调配，但背后透露的信号远比数据本身重要。当一家车企的海外需求强到需要“回购”国内经销商库存时，说明其海外渠道的销售速度已经超过了产能的全球调配节奏。更重要的是，这种回购行为实际上为国内经销商提供了一个“价格托底”——如果国内市场需求波动，海外渠道可以成为有效的缓冲池。

Citi专家还提到，比亚迪第二代刀片电池的月产能正在以

[... middle omitted ...]

可能面临停产，它们需要在“完全本地化以适应中国市场”和“彻底退出中国市场”之间做出选择。

这个判断的背景是5月燃油车销量的“崩盘式”下滑——Citi专家用“collapsed”来形容燃油车销量，而新能源车仅环比持平，说明燃油车下滑的份额并没有被新能源车完全承接，而是直接消失了。这背后的含义是：燃油车市场的需求萎缩并非简单的“消费者转向新能源”，而是整体购车意愿下降与燃油车竞争力丧失的叠加效应。

对于合资品牌而言，传统的“拿来主义”模式——将海外车型引入中国、进行小幅本地化调整——已经彻底失效。中国市场的电动化、智能化节奏远超全球其他市场，合资品牌母公司的产品规划周期（通常3-5年）与中国市场的迭代速度（通常12-18个月）之间存在巨大的时间错配。如果不能在中国建立独立的研发体系和供应链能力，合资品牌将无法跟上竞争节奏。

Citi专家的“二选一”判断虽然直接，但逻辑上成立。值得注意的是，报告本身并未给出具体的退出时间表或名单，这恰恰是需要投资者自行判断的关键变量。

\subsection{[R029] GS：AI基础设施定价权正在经历结构性重估，但企业级需求仍处于“等待爆发”的临界点}
{\small\color{refgray}归类：AI与算力：基础设施定价权重估与硬件ROE结构性上升}

GS：AI基础设施定价权正在经历结构性重估，但企业级需求仍处于“等待爆发”的临界点

市场对AI基础设施的讨论，大多聚焦于GPU算力是否过剩、资本开支何时见顶。但一份来自GS的最新实地调研报告，给出了一个更具穿透力的判断：当前真正被低估的，不是需求总量，而是供给侧约束正在重塑的定价结构。

这份报告基于GS在丹佛举办的第二届年度AI基础设施投资者之旅，核心内容包括对私有数据中心运营商Flexential的实地考察，以及对光纤网络运营商Lumen Technologies管理层的深度访谈。报告揭示了一个清晰但尚未被市场充分定价的信号：供给侧的物理瓶颈——从电力到设备到选址——正在将数据中心的定价权推向一个前所未有的水平，而这一趋势的持续性，可能远超市场预期。

与此同时，企业级AI基础设施的规模化采用，仍然处于“早期乐观但尚未落地”的阶段。这构成了当前AI基础设施投资的核心张力：一边是供给侧驱动的定价重估，一边是需求侧的结构性等待。理解这两者之间的错位，才是理解未来12-18个月相关资产表现的关键。

1. 供给约束已从短期扰动演变为结构性壁垒，数据中心定价权进入加速上行通道

GS调研中来自Flexential的信息，可能是这份报告最值得关注的信号。Flexential是一家专注于多租户和批发业务的数据中心运营商，其管理层的表述直接且具体：当前每千瓦的租赁价格已接近200美元，预计到今年晚些时候将向250美元迈进。作为对比，2021年这一价格仅为70美元左右。

这意味着，在过去四年中，数据中心单位容量的租赁价格已经上涨了近两倍，且涨价仍在加速。这不是一个周期性的波动，而是一个由供给侧结构性变化驱动的趋势。

第一，电力基础设施的交付周期已拉长至不可忽视的程度。Flexential指出，变电站的建设周期已达四年，发电机组的交付周期为两年。这意味着，即便今天做出投资决策，新增电力容量也需要数年才能就位。第二，数据中心设备本身的供应链瓶颈仍然存在，导致新建项目从启动到交付的时间已从2021年的约18个月延长至约三年。第三，选址难度在增加——社区反对、环保审批、土地竞争等因素都在收窄可开发空间。

这些约束叠加在一起，产生了一个关键的商业后果：数据中心运营商正在系统性地改变其销售策略。Flexential明确表示，为了避免在价格快速上升的环境中“贱卖”产能，其新产能的预销售窗口被限制在9个月以内。换句话说，他们宁可承担暂时的空置风险，也不愿锁定在未来的更高价格之下。

这一策略转变的行业含义是深远的。它意味着，数据中心的定价将从“成本加成”模式，逐步转向“稀缺定价”模式。对于拥有现成产能或快速交付能力的运营商而言，这直接转化为利润率的结构性提升。

2. 租赁续约中的价格爬升机制，正在成为数据中心运营商最稳定的利润引擎

如果说新产能定价体现的是市场对未来的预期，那么存量产能的续约定价，则直接反映了当前供需关系的紧张程度。Flexential的数据提供了一个清晰的量化基准：批发业务的年度续约价格爬升为3-4\%，多租户业务则达到5-6\%。

这些数字在传统商业地产租赁中并不常见。通常，商业地产的年租金涨幅在1-2\%之间，且往往受限于长期合约条款。数据中心之所以能够实现更高的续约涨幅，根本原因在于客户的粘性和转换成本。一旦企业的IT系统部署在一个数据中心内，物理迁移的成本和风险极高，这使得运营商在续约谈判中拥有显著的议价优势。

更重要的是，这一续约价格爬升机制具有自我强化的特征。随着新产能的定价不断走高，存量客户在续约时面对的“市场参考价”也在提升。Flexential管理层提到，续约谈判通常在租约到期前6-8个月启动，这意味着当前市场的高定价将很快传导到未来12-18个月的续约合同上。

对于投资者而言，这意味着数据

[... middle omitted ...]

East-West）连接市场的总可寻址市场规模为580亿美元，年复合增长率约为13\%。所谓“东西向流量”，指的是数据中心之间的数据交换，而非传统的用户到数据中心的“南北向流量”。AI训练和推理任务，特别是大规模分布式训练，会产生大量的东西向流量，这正是Lumen光纤网络的核心价值所在。

Lumen的差异化优势在于其物理网络的不可复制性。管理层强调，其最近签署的PCF（Private Connectivity Fabric）协议，使其能够在超大规模客户几乎全额预付的情况下扩建物理网络，从而大幅降低了建设成本。同时，其光纤与超大规模数据中心的地理接近性，使其能够提供多云端直连服务，绕过传统的托管中心交叉连接费用，显著降低延迟并提升安全性。

Lumen的Networking-as-a-Service平台是这一转型的载体。数据显示，2026年第一季度，超过20\%的NaaS净新增用户是新客户，超过60\%的现有客户在扩展其使用规模。同时，NaaS客户的流失率明显低于传统业务。

\subsection{[R030] UBS：东南亚外卖与电商的增长故事并未结束，但竞争正在重新定价龙头门槛}
{\small\color{refgray}归类：区域市场：亚洲外汇储备分化与东南亚互联网韧性}

UBS：东南亚外卖与电商的增长故事并未结束，但竞争正在重新定价龙头门槛

市场对东南亚互联网平台的担忧并非没有根据。宏观放缓、季节性退潮、以及折扣力度回升——这些信号叠加在一起，很容易让人产生“增长见顶”的直觉判断。但UBS最新发布的5月外卖与电商收据追踪数据，提供了一个截然不同的叙事框架：真正需要被关注的，不是增长是否在放缓，而是这种放缓的结构性含义——它正在检验哪些平台真正拥有定价权与用户粘性。

这份基于UBS Evidence Lab高频数据的报告，覆盖Grab与Shopee在东盟六国的核心运营指标。结论简洁但有力：两家公司均维持了20\%以上的同比GMV/订单增速，但竞争烈度确实在环比回升。市场低估的不是它们能否继续增长，而是它们在增长放缓时，能否将规模优势转化为不可逆的竞争壁垒。

1. 5月数据揭示了一个关键拐点：增长节奏从“爆发”切换为“韧性测试”

5月历来是东南亚消费的季节性低点，叠加今年开斋节时间调整，同比基数效应更为复杂。但UBS数据表明，Grab的外卖订单量同比仍增长21\%，虽然从3-4月的26\%-28\%有所回落。Shopee的GMV增速同样健康，同比+20.6\%，较4月的33\%和3月的23\%有所放缓。

这里需要区分两个概念：减速与恶化。减速是节奏调整，是季节性、基数效应和宏观噪音共同作用的结果；恶化则是结构性的需求崩塌或竞争失衡。UBS的数据支持前者。Grab的日均订单量环比仍增长2.8\%，Shopee日均订单环比增长6\%，说明用户活跃度并未出现断崖式下降。

关键在于，这种韧性是否能持续。印尼市场是Grab唯一录得订单同比下滑（-2.6\%）的地区，UBS将其归因于宏观疲软。如果印尼的疲软从个别市场扩散为东盟趋势，那么当前的韧性测试可能演变为真正的压力测试。这是报告没有完全展开，但值得密切跟踪的变量。

竞争烈度上升是5月数据中最容易引起误读的信号。Grab的折扣占AOV比例环比上升5-140个基点，Shopee的平台补贴同样环比走高。乍看之下，似乎回到了 年烧钱换增长的旧逻辑。

但UBS的解读提供了一个更精细的视角：Grab的AOV环比仅微降0.3\%，同比降幅也仅为1.1\%。这说明折扣的增加并没有导致客单价系统性下降。更可能的解释是，平台正在利用折扣工具调整用户结构——通过补贴吸引价格敏感型用户，扩大可触达市场总量（TAM），同时维持核心用户的客单价稳定。

Shopee的情况类似。其AOV同比下降5.6\%，但订单量同比增长28\%。量增价减，意味着平台在主动向低客单价、高频次的消费场景渗透。这种策略在电商领域并不新鲜，但关键在于执行效率：如果折扣能有效转化为复购率和用户生命周期价值，那么当前的补贴就是投资而非成本。

报告没有直接讨论的另一个问题是：折扣回升是否反映了来自TikTok Shop等新进入者的竞争压力？UBS的收据数据主要覆盖传统电商和外卖平台，对社交电商的渗透率变化缺乏直接观测。这可能是竞争格局中最大的盲区。

UBS对Grab的评级为买入，目标价5.9美元。支撑这一判断的核心论据，并非简单的增长数字，而是Grab在关键市场的份额变化。

在菲律宾和马来西亚，Grab分别从foodpanda手中夺取了69和136个基点的市场份额。这并非偶然。Grab在菲律宾的折扣环比大幅下降390个基点，说明其增长并非依赖补贴，而是产品力或运营效率的提升。在印尼，虽然订单同比下滑，但Grab的折扣环比仅微增5个基点，显示出其在弱势市场并未盲目跟投补贴。

这组数据指向一个更重要的结论：Grab正在从“依赖补贴维持份额”的阶段，过渡到“通过运营效率自然获得份额”的阶段。如果这一趋势持续，那么Grab的利润率扩张将不再是线性改善，而是可能加速。

但这里仍有一个未解问题：Grab在新加

[... middle omitted ...]

示Shopee在这两个核心市场的用户粘性已经足够强，不需要持续的补贴刺激。同时，卖家端补贴在所有东盟市场环比上升，说明Shopee正在从“补贴消费者”转向“赋能卖家”，这是一种更可持续的竞争策略。

报告没有深入讨论的是，Shopee的AOV持续下行（同比-5.6\%）是否会对单位经济模型产生压力。如果订单增长完全依赖低客单价商品，那么即使GMV增速健康，单均履约成本和收入结构可能面临挑战。这是判断Shopee估值逻辑是否成立的核心变量。

UBS在报告中多次提及“宏观疲软”和“需要密切关注”，但并未给出具体的压力测试场景。这是当前分析框架中最薄弱的环节。

具体而言，有三个问题值得进一步拆解：第一，印尼的订单下滑是周期性的还是结构性的？如果是结构性，Grab是否需要调整其资源分配策略？第二，如果东盟整体消费增速从当前水平进一步放缓，平台是否有能力在不牺牲市场份额的前提下削减补贴？第三，TikTok Shop等社交电商平台对传统电商的份额侵蚀，在收据数据中是否被低估了？

\subsection{[R031] Bernstein：中国车市真正的考验不是需求萎缩，而是补贴退潮后供给侧的再定价}
{\small\color{refgray}归类：汽车行业：补贴退潮暴露结构性压力，中国车企加速全球渗透}

Bernstein：中国车市真正的考验不是需求萎缩，而是补贴退潮后供给侧的再定价

2026年5月的中国汽车零售数据，看起来像一份“软着陆”报告。Bernstein追踪的交强险数据显示，当月国内零售销量同比下滑20.0\%，年化销量（SAAR）降至1920万辆，低于4月的1960万辆，更远低于该机构估算的2200万辆正常年需求水平。

但问题不在于数字本身。这份报告真正值得关注的判断是：市场目前的定价框架仍然在沿用“周期波动”的逻辑，而实际上中国汽车产业正在经历一次由政策退出、成本上升和出口结构变化共同驱动的供给侧再定价。Bernstein在维持行业谨慎展望的同时，对BYD和Xiaomi给出了Outperform评级，却同步下调了XPeng和Li Auto的目标价——这种分化本身就是一个信号。

1. 补贴退潮暴露的不是需求缺口，而是过去两年被透支的购买力

市场习惯用“需求疲软”来解释销量下滑，但Bernstein的数据指向了一个更结构性的问题。该机构估算， 年的补贴政策提前拉动了超过500万辆需求。这意味着5月的-20\%零售同比，本质上是“借来的时间”到期后的正常化过程。

更关键的是，在没有新政策支持的情况下，Bernstein判断国内需求将保持疲软，至少要到11-12月基数效应恢复正常。这不是消费者不愿意买车，而是他们已经在补贴窗口期把未来几个月的购买力提前释放了。

这对投资框架的含义是：如果市场把当前的销量下滑解读为“需求基本面恶化”，那么当基数效应在年底逆转时，可能出现预期差。但前提是，企业能够在不依赖补贴的情况下维持定价能力。

2. 电动车渗透率突破61\%背后的真正信号：消费者正在为更高价格买单

5月电动车渗透率升至61.0\%，其中纯电占41.9\%，插混占19.1\%。这组数字在表面上看起来是“电动化趋势不可逆”的证据，但Bernstein的解读更深一层。

从2026年1月1日起，电动车购置税率从零提高到5\%。理论上，这应该压制需求。但Bernstein观察到，消费者正在逐步适应更高的购车成本，同时新推出的电动车型在定价和产品竞争力上更强，汽油价格上涨也改善了电动车的使用经济性。

这意味着什么？电动车市场正在从“补贴驱动”转向“产品力驱动”。那些能够把成本上涨转嫁给消费者、同时通过产品差异化维持溢价的公司，将在这一轮调整中胜出。而那些依赖补贴价格优势的品牌，将面临利润率压缩的严峻考验。

3. 出口增长73\%不是锦上添花，而是部分企业生存的关键变量

5月出口同比增长73\%，其中电动车出口增长119\%，占整体出口量的约55\%。Bernstein的数据显示，出口已占乘用车批发总量的35.9\%，部分抵消了国内需求的疲软。

这个数字背后有两个值得拆解的含义。第一，出口结构正在从燃油车主导转向电动车主导，这意味着中国电动车企的全球化竞争不再是“低价倾销”，而是产品力和供应链优势的溢出。第二，对于像Chery、BYD、Geely这样出口量领先的公司，海外市场已经成为利润增长的第二曲线。

但问题在于，出口增长的可持续性取决于目标市场的贸易政策、本地化要求和竞争烈度。Bernstein没有在报告中展开的是：当更多中国车企涌入海外市场时，价格竞争是否会重演国内市场的剧本？这可能是未来12-18个月需要密切跟踪的风险。

4. Bernstein下调XPeng和Li Auto目标价的真正理由不是销量，而是盈利可见性

这是报告中最值得细读的部分。Bernstein将XPeng的XPEV目标价从22美元下调至20美元，估值倍数从1.0倍EV/Sales降至0.8倍，理由是“在AI投资高企的背景下盈利可见性减弱”。对于Li Auto，目标价从19美元降至15.5美元，基于盈利预测下调和不变的估值框架。

注意

[... middle omitted ...]

的估值框架下，市场对“故事”的耐心正在耗尽。AI投入、Robotaxi、人形机器人这些概念，如果不能转化为可量化的收入和利润，估值倍数就会被压缩。而产品线更新和品牌下沉，虽然能带来短期销量支撑，但不一定能改变盈利轨迹。

Bernstein在报告中展示了出口的高速增长，但没有深入拆解出口业务的利润率结构。这里存在一个值得追问的假设：中国车企在海外市场是否真的能维持与国内类似的毛利率？

从全球OEM的估值对比表中可以看到，中国EV车企的EV/Sales倍数（0.9倍）远高于欧洲OEM（0.1倍）和美国OEM（0.4倍），但EBIT利润率（5.4\%）却低于欧洲（5.4\% vs 13\%）和美国（5.4\% vs 10\%）。这意味着市场对中国车企的估值溢价主要来自增长预期，而不是当前盈利质量。

如果出口市场的价格竞争加剧，或者贸易壁垒提高导致本地化成本上升，那么当前的估值溢价将面临修正风险。这个问题在Bernstein的报告中只是被提及，但没有给出明确的结论，值得持续跟踪。

\section{Disclaimer}
{\scriptsize\color{refgray}
This community is a paid learning and information-sharing community created for general educational purposes only. The membership fee is charged solely for access to community discussions, educational content, organization, and operational costs. It is not an advisory fee, management fee, performance fee, brokerage fee, or compensation for personalized financial, investment, legal, tax, or professional advice.\par
I participate and share information solely as an individual and for educational discussion among community members. I am not acting as a financial adviser, investment adviser, broker, dealer, portfolio manager, legal adviser, tax adviser, accountant, fiduciary, or any other licensed professional adviser. Nothing shared in this community constitutes, and should not be interpreted as, financial advice, investment advice, legal advice, tax advice, a recommendation, solicitation, offer, endorsement, instruction, or guarantee to buy, sell, hold, short, trade, or invest in any security, cryptocurrency, fund, derivative, strategy, or other financial product.\par
All content shared in this community, including messages, discussions, links, files, charts, examples, market commentary, watchlists, AI-generated or AI-polished summaries, and third-party materials, is general, impersonal, and educational in nature. It does not take into account any individual's financial situation, investment objectives, risk tolerance, investment horizon, liquidity needs, tax status, legal restrictions, or suitability requirements. No content should be relied upon as suitable for any specific person, account, portfolio, or financial decision.\par
Information shared in this community may be incomplete, inaccurate, outdated, speculative, AI-generated, AI-edited, or based on third-party internet sources that have not been independently verified. I make no representation or warranty as to the accuracy, completeness, timeliness, reliability, or suitability of any information shared.\par
Financial markets involve substantial risk, including the possible loss of principal. Past performance, historical data, hypothetical examples, simulated results, backtests, personal opinions, or educational examples do not guarantee future results. Each member is solely responsible for conducting their own research, exercising independent judgment, and making their own decisions. Before making any financial, investment, legal, or tax decision, members should consult qualified and licensed professionals.\par
Participation in this community, payment of any membership fee, or communication with me or other members does not create any adviser-client, fiduciary, professional, contractual, agency, partnership, employment, or other advisory relationship. By joining or participating in this community, you acknowledge and agree that you are solely responsible for your own decisions, actions, transactions, profits, losses, risks, and consequences, and that I shall not be liable for any direct, indirect, incidental, consequential, financial, legal, tax, or other loss or damage arising from or related to any content, discussion, information, or material shared in this community.\par
}
\end{document}
